% AI+电力电子技术课程笔记 —— 主文件（正文排版在 0.样式.tex，章节文件只放内容）
% 编译：xelatex 1.main.tex 两遍；成品见 build.ps1
\documentclass[11pt,a4paper,oneside,openany]{book}
\usepackage[UTF8,fontset=windows,scheme=plain]{ctex}
\usepackage{fontspec}
\setmainfont{Times New Roman}
\setsansfont{Microsoft YaHei}
\newcommand{\YHBbooktrue}{}
\newcommand{\YHBcourse}{AI+电力电子技术}
\newcommand{\YHBversion}{v0.1}
\newcommand{\YHBquote}{用开关，把电从一种形态搬到另一种形态。}
\newcommand{\YHBcoverSize}{40pt}
\newcommand{\YHBcoverLead}{48pt}
\usepackage{xcolor}
\definecolor{yhbAccent}{HTML}{2A3A6B}
\definecolor{yhbAccent2}{HTML}{5169A8}
% ============================================================
% 0.样式.tex —— 课程笔记统一正文排版
% 用法：主文件先定义课程信息与主色，再 % ============================================================
% 0.样式.tex —— 课程笔记统一正文排版
% 用法：主文件先定义课程信息与主色，再 % ============================================================
% 0.样式.tex —— 课程笔记统一正文排版
% 用法：主文件先定义课程信息与主色，再 \input{0.样式.tex}
%   \newcommand{\YHBcourse}{微积分A}      % 课程名（页眉与书签用）
%   \newcommand{\YHBversion}{v1.2}        % 版本号（封面用）
%   \newcommand{\YHBquote}{……}            % 封面题句
%   \definecolor{yhbAccent}{HTML}{24465C} % 主色（深）
%   \definecolor{yhbAccent2}{HTML}{347A80}% 辅色（浅）
%   \YHBbooktrue / \YHBbookfalse          % 是否 book 文档类（有 \chapter）
% 章节文件只写内容，不要重复 \usepackage、不要改字号。
% ============================================================

% ---------- 页面与字体 ----------
\usepackage{geometry}
\geometry{a4paper,left=24mm,right=24mm,top=23mm,bottom=24mm,headheight=15pt,headsep=7mm,footskip=13mm}
\usepackage{microtype}
\usepackage{amsmath,amssymb,amsfonts,mathtools}
\usepackage{esint}
\usepackage{graphicx}
\usepackage{adjustbox}
\usepackage{array,booktabs,multirow,makecell,tabularx,longtable}
\usepackage{xcolor,colortbl}
\usepackage{enumitem}
\usepackage{needspace,etoolbox}
\usepackage{titlesec,titletoc}
\usepackage{fancyhdr}
\usepackage{caption}
\usepackage[most]{tcolorbox}

% ---------- 颜色别名（章节里用的旧名继续有效） ----------
\colorlet{mainblue}{yhbAccent}
\colorlet{secondcyan}{yhbAccent2}
\colorlet{thirdteal}{yhbAccent2}
\colorlet{highlight}{yhbAccent!6!white}
\colorlet{tablehead}{yhbAccent!12!white}
\definecolor{textgray}{HTML}{4A5057}

% ---------- 课程标签 ----------
\providecommand{\YHBdoclabel}{课程笔记}
% ---------- 正文节奏 ----------
\linespread{1.3}
\setlength{\parindent}{0pt}
\setlength{\parskip}{6pt plus 1.5pt minus 1.5pt}
\setlength{\emergencystretch}{2.5em}
\raggedbottom
\widowpenalty=10000
\clubpenalty=10000
\setlength{\tabcolsep}{7pt}
\renewcommand{\arraystretch}{1.3}
\allowdisplaybreaks[2]
\AtBeginDocument{%
  \setlength{\abovedisplayskip}{10pt plus 2pt minus 2pt}%
  \setlength{\belowdisplayskip}{10pt plus 2pt minus 2pt}%
  \setlength{\abovedisplayshortskip}{6pt}%
  \setlength{\belowdisplayshortskip}{6pt}%
}
\captionsetup{font=small,labelfont=bf,skip=6pt}

% ---------- 列表 ----------
\setlist{leftmargin=2em,labelsep=0.6em,itemsep=4pt,parsep=2pt,topsep=6pt,partopsep=0pt}
\setlist[itemize,2]{label=\textbullet,leftmargin=1.6em}

% ---------- 标题层级 ----------
\ifdefined\YHBbooktrue
  \titleformat{\chapter}[display]
    {\sffamily\bfseries\color{yhbAccent}}
    {\normalsize\sffamily 第\,\thechapter\,章}
    {6pt}
    {\LARGE}
    [{\vspace{5pt}{\color{yhbAccent!55}\titlerule[1pt]}}]
  \titlespacing*{\chapter}{0pt}{14pt}{24pt}
  \titleformat{\section}{\Large\sffamily\bfseries\color{yhbAccent}}{\thesection}{0.7em}{}
  \titlespacing*{\section}{0pt}{18pt plus 3pt minus 2pt}{8pt}
\else
  \titleformat{\section}{\Large\sffamily\bfseries\color{yhbAccent}}{\thesection}{0.7em}{}[{\vspace{3pt}{\color{yhbAccent!45}\titlerule[0.7pt]}}]
  \titlespacing*{\section}{0pt}{18pt plus 3pt minus 2pt}{9pt}
\fi
\titleformat{\subsection}{\large\sffamily\bfseries\color{yhbAccent!85!black}}{\thesubsection}{0.7em}{}
\titleformat{\subsubsection}{\normalsize\bfseries\color{black!85}}{\thesubsubsection}{0.7em}{}
\titlespacing*{\subsection}{0pt}{12pt plus 2pt minus 1pt}{5pt}
\titlespacing*{\subsubsection}{0pt}{9pt plus 2pt minus 1pt}{4pt}
\pretocmd{\section}{\Needspace{6\baselineskip}}{}{}
\pretocmd{\subsection}{\Needspace{5\baselineskip}}{}{}
\pretocmd{\subsubsection}{\Needspace{4\baselineskip}}{}{}

\ifdefined\YHBbooktrue
  \setcounter{secnumdepth}{2}
  \setcounter{tocdepth}{1}
\else
  \setcounter{secnumdepth}{3}
  \setcounter{tocdepth}{2}
\fi
\renewcommand{\contentsname}{目录}
% 目录编号颜色由 titletoc 的 titlecontents 控制，不再重定义 \numberline

% ---------- 目录样式 ----------
\ifdefined\YHBbooktrue
  \titlecontents{chapter}[0pt]
    {\addvspace{6pt}\sffamily\bfseries\color{yhbAccent}}
    {\contentslabel{2.2em}}{}{\titlerule*[0.6em]{.}\contentspage}
  \titlecontents{section}[2.2em]
    {\addvspace{2pt}}
    {\contentslabel{2.2em}}{}{\titlerule*[0.6em]{.}\contentspage}
\else
  \titlecontents{section}[0pt]
    {\addvspace{5pt}\sffamily\bfseries\color{yhbAccent}}
    {\contentslabel{2.2em}}{}{\titlerule*[0.6em]{.}\contentspage}
  \titlecontents{subsection}[2.2em]
    {\addvspace{2pt}}
    {\contentslabel{2.4em}}{}{\titlerule*[0.6em]{.}\contentspage}
\fi

% ---------- 页眉页脚 ----------
\pagestyle{fancy}
\fancyhf{}
\fancyhead[L]{\footnotesize\sffamily\color{textgray}\YHBcourse\ \YHBdoclabel}
\ifdefined\YHBbooktrue
  \fancyhead[R]{\footnotesize\sffamily\color{textgray}\nouppercase{\leftmark}}
  \renewcommand{\chaptermark}[1]{\markboth{#1}{}}
\else
  \fancyhead[R]{\footnotesize\sffamily\color{textgray}\nouppercase{\rightmark}}
  \renewcommand{\sectionmark}[1]{\markright{#1}}
\fi
\fancyfoot[C]{\footnotesize\color{textgray}\thepage}
\renewcommand{\headrulewidth}{0.35pt}
\renewcommand{\headrule}{{\color{yhbAccent!45}\hrule width\headwidth height\headrulewidth}}
\fancypagestyle{plain}{%
  \fancyhf{}%
  \fancyfoot[C]{\footnotesize\color{textgray}\thepage}%
  \renewcommand{\headrulewidth}{0pt}%
}

% ---------- 公式与考点框 ----------
\newcommand{\keyformula}[1]{%
  \begin{tcolorbox}[enhanced,breakable,colback=highlight,colframe=yhbAccent!30!white,boxrule=0.5pt,arc=2.5pt,left=11pt,right=11pt,top=9pt,bottom=9pt,boxsep=0pt,before skip=11pt,after skip=11pt,halign=center,width=\linewidth]
  \adjustbox{max width=\linewidth-2em}{$\displaystyle #1$}
  \end{tcolorbox}}
\newcommand{\corepoint}[1]{%
  \begin{tcolorbox}[enhanced,breakable,colback=highlight,colframe=yhbAccent2,boxrule=0pt,borderline west={2.4pt}{0pt}{yhbAccent2},arc=0pt,left=11pt,right=11pt,top=8pt,bottom=8pt,before skip=11pt,after skip=11pt]
  \textbf{\color{yhbAccent}核心考点}\quad #1
  \end{tcolorbox}}
\newcommand{\symb}[2]{\textbf{#1}：#2\par}

% ---------- 超链接（务必最后加载） ----------
\usepackage[unicode=true,colorlinks=true,linkcolor=yhbAccent,citecolor=yhbAccent,urlcolor=yhbAccent!85!black,pdfborder={0 0 0},bookmarksnumbered=true,pdfstartview=FitH]{hyperref}
\usepackage{hypcap}
\hypersetup{pdftitle={\YHBcourse 课程笔记},pdfauthor={浮萍},pdfsubject={课程学习与复习笔记},pdfcreator={XeLaTeX}}

%   \newcommand{\YHBcourse}{微积分A}      % 课程名（页眉与书签用）
%   \newcommand{\YHBversion}{v1.2}        % 版本号（封面用）
%   \newcommand{\YHBquote}{……}            % 封面题句
%   \definecolor{yhbAccent}{HTML}{24465C} % 主色（深）
%   \definecolor{yhbAccent2}{HTML}{347A80}% 辅色（浅）
%   \YHBbooktrue / \YHBbookfalse          % 是否 book 文档类（有 \chapter）
% 章节文件只写内容，不要重复 \usepackage、不要改字号。
% ============================================================

% ---------- 页面与字体 ----------
\usepackage{geometry}
\geometry{a4paper,left=24mm,right=24mm,top=23mm,bottom=24mm,headheight=15pt,headsep=7mm,footskip=13mm}
\usepackage{microtype}
\usepackage{amsmath,amssymb,amsfonts,mathtools}
\usepackage{esint}
\usepackage{graphicx}
\usepackage{adjustbox}
\usepackage{array,booktabs,multirow,makecell,tabularx,longtable}
\usepackage{xcolor,colortbl}
\usepackage{enumitem}
\usepackage{needspace,etoolbox}
\usepackage{titlesec,titletoc}
\usepackage{fancyhdr}
\usepackage{caption}
\usepackage[most]{tcolorbox}

% ---------- 颜色别名（章节里用的旧名继续有效） ----------
\colorlet{mainblue}{yhbAccent}
\colorlet{secondcyan}{yhbAccent2}
\colorlet{thirdteal}{yhbAccent2}
\colorlet{highlight}{yhbAccent!6!white}
\colorlet{tablehead}{yhbAccent!12!white}
\definecolor{textgray}{HTML}{4A5057}

% ---------- 课程标签 ----------
\providecommand{\YHBdoclabel}{课程笔记}
% ---------- 正文节奏 ----------
\linespread{1.3}
\setlength{\parindent}{0pt}
\setlength{\parskip}{6pt plus 1.5pt minus 1.5pt}
\setlength{\emergencystretch}{2.5em}
\raggedbottom
\widowpenalty=10000
\clubpenalty=10000
\setlength{\tabcolsep}{7pt}
\renewcommand{\arraystretch}{1.3}
\allowdisplaybreaks[2]
\AtBeginDocument{%
  \setlength{\abovedisplayskip}{10pt plus 2pt minus 2pt}%
  \setlength{\belowdisplayskip}{10pt plus 2pt minus 2pt}%
  \setlength{\abovedisplayshortskip}{6pt}%
  \setlength{\belowdisplayshortskip}{6pt}%
}
\captionsetup{font=small,labelfont=bf,skip=6pt}

% ---------- 列表 ----------
\setlist{leftmargin=2em,labelsep=0.6em,itemsep=4pt,parsep=2pt,topsep=6pt,partopsep=0pt}
\setlist[itemize,2]{label=\textbullet,leftmargin=1.6em}

% ---------- 标题层级 ----------
\ifdefined\YHBbooktrue
  \titleformat{\chapter}[display]
    {\sffamily\bfseries\color{yhbAccent}}
    {\normalsize\sffamily 第\,\thechapter\,章}
    {6pt}
    {\LARGE}
    [{\vspace{5pt}{\color{yhbAccent!55}\titlerule[1pt]}}]
  \titlespacing*{\chapter}{0pt}{14pt}{24pt}
  \titleformat{\section}{\Large\sffamily\bfseries\color{yhbAccent}}{\thesection}{0.7em}{}
  \titlespacing*{\section}{0pt}{18pt plus 3pt minus 2pt}{8pt}
\else
  \titleformat{\section}{\Large\sffamily\bfseries\color{yhbAccent}}{\thesection}{0.7em}{}[{\vspace{3pt}{\color{yhbAccent!45}\titlerule[0.7pt]}}]
  \titlespacing*{\section}{0pt}{18pt plus 3pt minus 2pt}{9pt}
\fi
\titleformat{\subsection}{\large\sffamily\bfseries\color{yhbAccent!85!black}}{\thesubsection}{0.7em}{}
\titleformat{\subsubsection}{\normalsize\bfseries\color{black!85}}{\thesubsubsection}{0.7em}{}
\titlespacing*{\subsection}{0pt}{12pt plus 2pt minus 1pt}{5pt}
\titlespacing*{\subsubsection}{0pt}{9pt plus 2pt minus 1pt}{4pt}
\pretocmd{\section}{\Needspace{6\baselineskip}}{}{}
\pretocmd{\subsection}{\Needspace{5\baselineskip}}{}{}
\pretocmd{\subsubsection}{\Needspace{4\baselineskip}}{}{}

\ifdefined\YHBbooktrue
  \setcounter{secnumdepth}{2}
  \setcounter{tocdepth}{1}
\else
  \setcounter{secnumdepth}{3}
  \setcounter{tocdepth}{2}
\fi
\renewcommand{\contentsname}{目录}
% 目录编号颜色由 titletoc 的 titlecontents 控制，不再重定义 \numberline

% ---------- 目录样式 ----------
\ifdefined\YHBbooktrue
  \titlecontents{chapter}[0pt]
    {\addvspace{6pt}\sffamily\bfseries\color{yhbAccent}}
    {\contentslabel{2.2em}}{}{\titlerule*[0.6em]{.}\contentspage}
  \titlecontents{section}[2.2em]
    {\addvspace{2pt}}
    {\contentslabel{2.2em}}{}{\titlerule*[0.6em]{.}\contentspage}
\else
  \titlecontents{section}[0pt]
    {\addvspace{5pt}\sffamily\bfseries\color{yhbAccent}}
    {\contentslabel{2.2em}}{}{\titlerule*[0.6em]{.}\contentspage}
  \titlecontents{subsection}[2.2em]
    {\addvspace{2pt}}
    {\contentslabel{2.4em}}{}{\titlerule*[0.6em]{.}\contentspage}
\fi

% ---------- 页眉页脚 ----------
\pagestyle{fancy}
\fancyhf{}
\fancyhead[L]{\footnotesize\sffamily\color{textgray}\YHBcourse\ \YHBdoclabel}
\ifdefined\YHBbooktrue
  \fancyhead[R]{\footnotesize\sffamily\color{textgray}\nouppercase{\leftmark}}
  \renewcommand{\chaptermark}[1]{\markboth{#1}{}}
\else
  \fancyhead[R]{\footnotesize\sffamily\color{textgray}\nouppercase{\rightmark}}
  \renewcommand{\sectionmark}[1]{\markright{#1}}
\fi
\fancyfoot[C]{\footnotesize\color{textgray}\thepage}
\renewcommand{\headrulewidth}{0.35pt}
\renewcommand{\headrule}{{\color{yhbAccent!45}\hrule width\headwidth height\headrulewidth}}
\fancypagestyle{plain}{%
  \fancyhf{}%
  \fancyfoot[C]{\footnotesize\color{textgray}\thepage}%
  \renewcommand{\headrulewidth}{0pt}%
}

% ---------- 公式与考点框 ----------
\newcommand{\keyformula}[1]{%
  \begin{tcolorbox}[enhanced,breakable,colback=highlight,colframe=yhbAccent!30!white,boxrule=0.5pt,arc=2.5pt,left=11pt,right=11pt,top=9pt,bottom=9pt,boxsep=0pt,before skip=11pt,after skip=11pt,halign=center,width=\linewidth]
  \adjustbox{max width=\linewidth-2em}{$\displaystyle #1$}
  \end{tcolorbox}}
\newcommand{\corepoint}[1]{%
  \begin{tcolorbox}[enhanced,breakable,colback=highlight,colframe=yhbAccent2,boxrule=0pt,borderline west={2.4pt}{0pt}{yhbAccent2},arc=0pt,left=11pt,right=11pt,top=8pt,bottom=8pt,before skip=11pt,after skip=11pt]
  \textbf{\color{yhbAccent}核心考点}\quad #1
  \end{tcolorbox}}
\newcommand{\symb}[2]{\textbf{#1}：#2\par}

% ---------- 超链接（务必最后加载） ----------
\usepackage[unicode=true,colorlinks=true,linkcolor=yhbAccent,citecolor=yhbAccent,urlcolor=yhbAccent!85!black,pdfborder={0 0 0},bookmarksnumbered=true,pdfstartview=FitH]{hyperref}
\usepackage{hypcap}
\hypersetup{pdftitle={\YHBcourse 课程笔记},pdfauthor={浮萍},pdfsubject={课程学习与复习笔记},pdfcreator={XeLaTeX}}

%   \newcommand{\YHBcourse}{微积分A}      % 课程名（页眉与书签用）
%   \newcommand{\YHBversion}{v1.2}        % 版本号（封面用）
%   \newcommand{\YHBquote}{……}            % 封面题句
%   \definecolor{yhbAccent}{HTML}{24465C} % 主色（深）
%   \definecolor{yhbAccent2}{HTML}{347A80}% 辅色（浅）
%   \YHBbooktrue / \YHBbookfalse          % 是否 book 文档类（有 \chapter）
% 章节文件只写内容，不要重复 \usepackage、不要改字号。
% ============================================================

% ---------- 页面与字体 ----------
\usepackage{geometry}
\geometry{a4paper,left=24mm,right=24mm,top=23mm,bottom=24mm,headheight=15pt,headsep=7mm,footskip=13mm}
\usepackage{microtype}
\usepackage{amsmath,amssymb,amsfonts,mathtools}
\usepackage{esint}
\usepackage{graphicx}
\usepackage{adjustbox}
\usepackage{array,booktabs,multirow,makecell,tabularx,longtable}
\usepackage{xcolor,colortbl}
\usepackage{enumitem}
\usepackage{needspace,etoolbox}
\usepackage{titlesec,titletoc}
\usepackage{fancyhdr}
\usepackage{caption}
\usepackage[most]{tcolorbox}

% ---------- 颜色别名（章节里用的旧名继续有效） ----------
\colorlet{mainblue}{yhbAccent}
\colorlet{secondcyan}{yhbAccent2}
\colorlet{thirdteal}{yhbAccent2}
\colorlet{highlight}{yhbAccent!6!white}
\colorlet{tablehead}{yhbAccent!12!white}
\definecolor{textgray}{HTML}{4A5057}

% ---------- 课程标签 ----------
\providecommand{\YHBdoclabel}{课程笔记}
% ---------- 正文节奏 ----------
\linespread{1.3}
\setlength{\parindent}{0pt}
\setlength{\parskip}{6pt plus 1.5pt minus 1.5pt}
\setlength{\emergencystretch}{2.5em}
\raggedbottom
\widowpenalty=10000
\clubpenalty=10000
\setlength{\tabcolsep}{7pt}
\renewcommand{\arraystretch}{1.3}
\allowdisplaybreaks[2]
\AtBeginDocument{%
  \setlength{\abovedisplayskip}{10pt plus 2pt minus 2pt}%
  \setlength{\belowdisplayskip}{10pt plus 2pt minus 2pt}%
  \setlength{\abovedisplayshortskip}{6pt}%
  \setlength{\belowdisplayshortskip}{6pt}%
}
\captionsetup{font=small,labelfont=bf,skip=6pt}

% ---------- 列表 ----------
\setlist{leftmargin=2em,labelsep=0.6em,itemsep=4pt,parsep=2pt,topsep=6pt,partopsep=0pt}
\setlist[itemize,2]{label=\textbullet,leftmargin=1.6em}

% ---------- 标题层级 ----------
\ifdefined\YHBbooktrue
  \titleformat{\chapter}[display]
    {\sffamily\bfseries\color{yhbAccent}}
    {\normalsize\sffamily 第\,\thechapter\,章}
    {6pt}
    {\LARGE}
    [{\vspace{5pt}{\color{yhbAccent!55}\titlerule[1pt]}}]
  \titlespacing*{\chapter}{0pt}{14pt}{24pt}
  \titleformat{\section}{\Large\sffamily\bfseries\color{yhbAccent}}{\thesection}{0.7em}{}
  \titlespacing*{\section}{0pt}{18pt plus 3pt minus 2pt}{8pt}
\else
  \titleformat{\section}{\Large\sffamily\bfseries\color{yhbAccent}}{\thesection}{0.7em}{}[{\vspace{3pt}{\color{yhbAccent!45}\titlerule[0.7pt]}}]
  \titlespacing*{\section}{0pt}{18pt plus 3pt minus 2pt}{9pt}
\fi
\titleformat{\subsection}{\large\sffamily\bfseries\color{yhbAccent!85!black}}{\thesubsection}{0.7em}{}
\titleformat{\subsubsection}{\normalsize\bfseries\color{black!85}}{\thesubsubsection}{0.7em}{}
\titlespacing*{\subsection}{0pt}{12pt plus 2pt minus 1pt}{5pt}
\titlespacing*{\subsubsection}{0pt}{9pt plus 2pt minus 1pt}{4pt}
\pretocmd{\section}{\Needspace{6\baselineskip}}{}{}
\pretocmd{\subsection}{\Needspace{5\baselineskip}}{}{}
\pretocmd{\subsubsection}{\Needspace{4\baselineskip}}{}{}

\ifdefined\YHBbooktrue
  \setcounter{secnumdepth}{2}
  \setcounter{tocdepth}{1}
\else
  \setcounter{secnumdepth}{3}
  \setcounter{tocdepth}{2}
\fi
\renewcommand{\contentsname}{目录}
% 目录编号颜色由 titletoc 的 titlecontents 控制，不再重定义 \numberline

% ---------- 目录样式 ----------
\ifdefined\YHBbooktrue
  \titlecontents{chapter}[0pt]
    {\addvspace{6pt}\sffamily\bfseries\color{yhbAccent}}
    {\contentslabel{2.2em}}{}{\titlerule*[0.6em]{.}\contentspage}
  \titlecontents{section}[2.2em]
    {\addvspace{2pt}}
    {\contentslabel{2.2em}}{}{\titlerule*[0.6em]{.}\contentspage}
\else
  \titlecontents{section}[0pt]
    {\addvspace{5pt}\sffamily\bfseries\color{yhbAccent}}
    {\contentslabel{2.2em}}{}{\titlerule*[0.6em]{.}\contentspage}
  \titlecontents{subsection}[2.2em]
    {\addvspace{2pt}}
    {\contentslabel{2.4em}}{}{\titlerule*[0.6em]{.}\contentspage}
\fi

% ---------- 页眉页脚 ----------
\pagestyle{fancy}
\fancyhf{}
\fancyhead[L]{\footnotesize\sffamily\color{textgray}\YHBcourse\ \YHBdoclabel}
\ifdefined\YHBbooktrue
  \fancyhead[R]{\footnotesize\sffamily\color{textgray}\nouppercase{\leftmark}}
  \renewcommand{\chaptermark}[1]{\markboth{#1}{}}
\else
  \fancyhead[R]{\footnotesize\sffamily\color{textgray}\nouppercase{\rightmark}}
  \renewcommand{\sectionmark}[1]{\markright{#1}}
\fi
\fancyfoot[C]{\footnotesize\color{textgray}\thepage}
\renewcommand{\headrulewidth}{0.35pt}
\renewcommand{\headrule}{{\color{yhbAccent!45}\hrule width\headwidth height\headrulewidth}}
\fancypagestyle{plain}{%
  \fancyhf{}%
  \fancyfoot[C]{\footnotesize\color{textgray}\thepage}%
  \renewcommand{\headrulewidth}{0pt}%
}

% ---------- 公式与考点框 ----------
\newcommand{\keyformula}[1]{%
  \begin{tcolorbox}[enhanced,breakable,colback=highlight,colframe=yhbAccent!30!white,boxrule=0.5pt,arc=2.5pt,left=11pt,right=11pt,top=9pt,bottom=9pt,boxsep=0pt,before skip=11pt,after skip=11pt,halign=center,width=\linewidth]
  \adjustbox{max width=\linewidth-2em}{$\displaystyle #1$}
  \end{tcolorbox}}
\newcommand{\corepoint}[1]{%
  \begin{tcolorbox}[enhanced,breakable,colback=highlight,colframe=yhbAccent2,boxrule=0pt,borderline west={2.4pt}{0pt}{yhbAccent2},arc=0pt,left=11pt,right=11pt,top=8pt,bottom=8pt,before skip=11pt,after skip=11pt]
  \textbf{\color{yhbAccent}核心考点}\quad #1
  \end{tcolorbox}}
\newcommand{\symb}[2]{\textbf{#1}：#2\par}

% ---------- 超链接（务必最后加载） ----------
\usepackage[unicode=true,colorlinks=true,linkcolor=yhbAccent,citecolor=yhbAccent,urlcolor=yhbAccent!85!black,pdfborder={0 0 0},bookmarksnumbered=true,pdfstartview=FitH]{hyperref}
\usepackage{hypcap}
\hypersetup{pdftitle={\YHBcourse 课程笔记},pdfauthor={浮萍},pdfsubject={课程学习与复习笔记},pdfcreator={XeLaTeX}}

\renewcommand{\tablename}{表}
\renewcommand{\figurename}{图}
\usepackage{tikz}
\usepackage{circuitikz}
\begin{document}
% ============================================================
% 2.cover.tex —— 统一封面（信息全部来自主文件定义）
% 需要的宏：\YHBcourse、\YHBversion、\YHBquote、yhbAccent、yhbAccent2
% 可选：\YHBcoverSize / \YHBcoverLead（课程名字号与行距）、\YHBcoversub（副标题）
% ============================================================
\providecommand{\YHBcoverSize}{38pt}
\providecommand{\YHBcoverLead}{46pt}
\providecommand{\YHBcoversub}{课程笔记}
\begin{titlepage}
\thispagestyle{empty}
\pdfbookmark[0]{封面}{cover}

\begin{tikzpicture}[remember picture,overlay]
  % 顶部深色块
  \fill[yhbAccent] (current page.north west) rectangle ([yshift=-132mm]current page.north east);
  % 分隔细条（辅色）
  \fill[yhbAccent2] ([yshift=-132mm]current page.north west) rectangle ([yshift=-135mm]current page.north east);
  % 极细横线，弱装饰，不压字
  \foreach \y in {20,34,48,62,76,90,104,118}{%
    \draw[white,opacity=0.08,line width=0.5pt]
      ([xshift=16mm,yshift=-\y mm]current page.north west) --
      ([xshift=-16mm,yshift=-\y mm]current page.north east);}
\end{tikzpicture}

\vspace*{32mm}
\noindent
{\xeCJKsetup{CJKecglue={}}\color{white}\sffamily\bfseries\fontsize{\YHBcoverSize}{\YHBcoverLead}\selectfont \YHBcourse\par}
\vspace{6mm}
{\color{white}\sffamily\fontsize{25}{32}\selectfont \YHBcoversub\par}
\vspace{9mm}
{\color{white}\fontsize{12.5}{18}\selectfont 作者\quad 浮萍\qquad 版本\quad \YHBversion\par}

\vspace{20mm}
{\sffamily\bfseries\fontsize{13}{20}\selectfont\color{yhbAccent}\YHBquote\par}

\vfill
{\color{yhbAccent}\rule{\textwidth}{1.4pt}}\par
\vspace{7mm}
{\fontsize{11.5}{17}\selectfont\sffamily
\textbf{QQ 邮箱}\quad\href{mailto:3071974740@qq.com}{3071974740@qq.com}\par
\vspace{2mm}
\textbf{GitHub}\quad\href{https://github.com/Duckweed-yhb}{github.com/Duckweed-yhb}\par}
\vspace{7mm}
{\fontsize{10.5}{16}\selectfont\color{black!78}
本笔记为个人学习整理，内容仍在持续完善。如发现内容、公式、文字或排版问题，欢迎通过上述 QQ 邮箱反馈；请尽量注明版本、页码与问题描述，便于核对和修订。\par}
\vspace*{12mm}
\end{titlepage}

\newpage
\tableofcontents
\newpage
% ===== 章节目录：在此按教材章节顺序 \input =====
\chapter{绪论}

电力电子技术是电气工程学科中最活跃的分支之一。它研究的是“如何用半导体开关把电能从一种形态搬运成另一种形态”，因此它既要有电路理论作骨架，又要有控制理论作神经。本章回答四个问题：什么是电力电子技术、它是怎么发展起来的、它用在哪里、这本教材怎么读。

% ============================================================
\section{什么是电力电子技术}

\subsection{定义与学科定位}

\begin{tcolorbox}[enhanced,breakable,colback=highlight,colframe=yhbAccent!30!white,boxrule=0.5pt,arc=2.5pt,left=11pt,right=11pt,top=9pt,bottom=9pt,boxsep=0pt,before skip=11pt,after skip=11pt]
\textbf{\color{yhbAccent}定义（电力电子技术）}\quad 电力电子技术是\textbf{使用电力电子器件对电能进行变换和控制的技术}。目前所用的电力电子器件均由半导体制成，故也称电力半导体器件。
\end{tcolorbox}

几点必须分清：

\begin{itemize}
  \item 电子技术分为\textbf{信息电子技术}与\textbf{电力电子技术}两大分支。通常所说的模拟电子技术、数字电子技术都属于信息电子技术。
  \item 信息电子技术主要用于\textbf{信息处理}，电力电子技术主要用于\textbf{电力变换}——这是二者本质上的不同。
  \item 电力电子技术中所变换的“电力”功率可以大到数百兆瓦甚至吉瓦，也可以小到数瓦甚至毫瓦级，跨度极大。
  \item 广义而言，电子电路中的功率放大和功率输出部分也可算作电力电子电路；各种电子装置（电视机、计算机）的电源部分都是电力电子电路。
\end{itemize}

\subsection{电力变换的四大类}

常用的电力有交流和直流两种。从公用电网直接得到的是交流，从蓄电池和干电池得到的是直流。这些电源往往不能直接满足要求，需要进行电力变换。电力变换通常分为四大类，见表~\ref{tab:1-1}。

\begin{table}[htbp]
\centering
\caption{电力变换的种类}
\label{tab:1-1}
\begin{tabular}{llll}
\toprule
\textbf{变换类型} & \textbf{英文缩写} & \textbf{习惯名称} & \textbf{实现电路} \\
\midrule
交流 $\rightarrow$ 直流 & AC--DC & 整流（Rectifier） & 整流电路 \\
直流 $\rightarrow$ 交流 & DC--AC & 逆变（Inverter） & 逆变电路 \\
直流 $\rightarrow$ 直流 & DC--DC & 直流斩波 & 直流斩波电路 \\
交流 $\rightarrow$ 交流 & AC--AC & 交流电力控制 / 变频 / 变相 & 交交变频、交流调压 \\
\bottomrule
\end{tabular}
\end{table}

\corepoint{直流变直流\textbf{并非电力种类的变换}，而是电压（或电流）的变换，即一种直流电压（或电流）变为另一种直流电压（或电流）；交流变交流除了电压或电流的变换外，还多了\textbf{频率或相数}的变换。进行上述电力变换的技术统称\textbf{变流技术}。}

\subsection{两个分支：器件制造技术与变流技术}

通常把电力电子技术分为\textbf{电力电子器件制造技术}和\textbf{变流技术}两个分支。变流技术也称为电力电子器件的应用技术，包括用电力电子器件构成各种电力变换电路、对这些电路进行控制、以及由这些电路构成电力电子装置和系统的技术。

\begin{itemize}
  \item 器件制造技术的理论基础是\textbf{半导体物理}；
  \item 变流技术的理论基础是\textbf{电路理论}；
  \item 没有晶闸管、电力晶体管、IGBT 等器件，就没有电力电子技术，所以\textbf{器件制造技术是基础}；而电力电子技术主要用于电力变换，所以\textbf{变流技术是核心}。
\end{itemize}

\subsection{“倒三角形”：电力电子学的学科交叉}

电力电子学（Power Electronics）这一名称出现在 20 世纪 60 年代（比晶闸管的出现晚）。1974 年，美国学者 W. Newell 用倒三角形描述了这一学科：它由电力学、电子学和控制理论三个学科交叉形成，这一观点被全世界普遍接受。

\begin{figure}[htbp]
\centering
\begin{tikzpicture}[font=\small]
  \draw[yhbAccent,thick] (0,0) -- (6,0) -- (3,-3.2) -- cycle;
  \node[text=yhbAccent] at (3,-1.55) {\textbf{电力电子技术}};
  \node at (1.55,-0.45) {电力学};
  \node at (4.45,-0.45) {电子学};
  \node at (3,-2.85) {控制理论};
  \node[text=textgray,font=\footnotesize] at (3,0.35) {连续、离散控制理论};
\end{tikzpicture}
\caption{描述电力电子学的倒三角形}
\label{fig:1-1}
\end{figure}

“电力电子学”和“电力电子技术”是分别从学术和工程技术两个不同角度来称呼的，其实际内容并没有很大的不同。

\subsection{与信息电子技术的根本区别}

电力电子技术与电子学的关系是显见的：信息电子学可分为电子器件和电子电路两大分支，这分别与电力电子器件和电力电子电路相对应；二者的制造理论基础（半导体理论）相同，大多数工艺也相同，现代电力电子器件的制造大都使用集成电路制造工艺。

差别在应用目的与工作状态：

\corepoint{信息电子技术中，半导体器件既可处于放大状态，也可处于开关状态；而电力电子技术中，\textbf{为避免功率损耗过大，电力电子器件总是工作在开关状态}。这是电力电子技术区别于信息电子技术的一个重要特征。}

电力电子技术广泛用于电气工程中，这是电力电子学和电力学的主要关系。电力电子装置广泛应用于高压直流输电、静止无功补偿、电力机车牵引、交直流电力传动、电解、励磁、电加热、高性能交直流电源之中，因此国内外通常把电力电子技术归属于电气工程学科。在我国的学科分类中，电气工程是一级学科，包含五个二级学科：电力系统及其自动化、电机与电器、高电压与绝缘技术、电力电子与电力传动、电工理论与新技术。

% ============================================================
\section{电力电子技术的发展史}

电力电子器件的发展对电力电子技术的发展起着决定性的作用，因此电力电子技术的发展史是以电力电子器件的发展史为纲的。

\subsection{史前期（黎明期）}

\begin{itemize}
  \item 1904 年出现\textbf{电子管}，它能在真空中对电子流进行控制，并应用于通信和无线电，从而开启了电子技术用于电力领域的先河。
  \item 后来出现\textbf{水银整流器}（汞弧整流器），把水银封于管内，利用对其蒸气的点弧可对大电流进行控制，其性能和晶闸管已经非常相似。缺点是水银对人体有害、电压降落也很高。
  \item 20 世纪 30--50 年代是水银整流器发展迅速并大量应用的时期，广泛用于电化学工业、电气铁道直流变电所、轧钢用直流电动机传动，甚至用于直流输电。这一时期，各种整流电路、逆变电路、周波变流电路的理论已经发展成熟并广为应用。
  \item 把交流变为直流的方法，除水银整流器外，还有发展更早的\textbf{电动机—直流发电机组}，即\textbf{变流机组}（旋转变流机组）。与旋转变流机组相对应，静止变流器的称呼从水银整流器开始而沿用。
\end{itemize}

\subsection{晶闸管时代}

1947 年，美国贝尔实验室发明了\textbf{晶体管}，引发了电子技术的一场革命；最先用于电力领域的半导体器件是硅二极管。1957 年，美国通用电气公司研制出第一个\textbf{晶闸管}，一般以此作为电力电子技术诞生的标志。

\corepoint{晶闸管是\textbf{通过对门极的控制能够使其导通而不能使其关断}的器件，属于\textbf{半控型器件}；对晶闸管电路的控制方式主要是\textbf{相位控制方式}（简称相控方式）；晶闸管的关断通常依靠电网电压等外部条件来实现。这使晶闸管的应用受到很大局限。}

晶闸管出现后，由于其优越的电气性能和控制性能，很快取代了水银整流器和旋转变流机组，应用范围迅速扩大。电化学工业、铁道电气机车、钢铁工业（轧钢电气传动、感应加热等）、电力工业（直流输电、无功补偿等）的迅速发展也给晶闸管提供了用武之地。\textbf{电力电子技术的概念和基础，就是由于晶闸管及晶闸管变流技术的发展而确立的}。

\subsection{全控型器件与复合型器件}

20 世纪 70 年代后期，以\textbf{门极可关断晶闸管（GTO）}、\textbf{电力双极型晶体管（BJT）}和\textbf{电力场效应晶体管（Power MOSFET）}为代表的全控型器件迅速发展。

\begin{itemize}
  \item 全控型器件的特点是：通过对门极（基极、栅极）的控制\textbf{既可使它开通又可使它关断}；开关速度普遍高于晶闸管，可用于开关频率较高的电路。
  \item 与晶闸管电路的相控方式相对应，采用全控型器件的电路主要控制方式为\textbf{脉冲宽度调制（PWM）}方式，相对于相位控制方式可称为\textbf{斩波控制方式}（斩控方式）。
\end{itemize}

20 世纪 80 年代后期，以\textbf{绝缘栅极双极型晶体管（IGBT）}为代表的\textbf{复合型器件}异军突起。IGBT 是 MOSFET 和 BJT 的复合，它把 MOSFET 驱动功率小、开关速度快的优点，和 BJT 通态压降小、载流能力大、可承受电压高的优点集于一身，性能十分优越，成为现代电力电子技术的主导器件。与之对应，\textbf{MOS 控制晶闸管（MCT）}和\textbf{集成门极换流晶闸管（IGCT）}都是 MOSFET 和 GTO 的复合，也综合了两种器件的优点，其中 IGCT 已经获得大量应用。

为使装置结构紧凑、体积减小，常把若干个电力电子器件及必要的辅助元件做成\textbf{模块}形式；后来又进一步把驱动、控制、保护电路和器件集成在一起，构成\textbf{电力电子集成电路（PIC）}。目前 PIC 的功率还较小、电压也较低，它面临着电压隔离（主电路为高压而控制电路为低压）、热隔离（主电路发热严重）、电磁干扰（开关器件通断高压大电流且与控制电路同处一个芯片）等几大难题，但这代表了电力电子技术发展的一个重要方向。电力电子集成技术的焦点是\textbf{混合集成技术}，即把不同的单个芯片集成封装在一起；此外\textbf{系统集成}也是重要方面，对超大功率集成技术尤其如此。

随着全控型器件不断进步，电力电子电路的工作频率不断提高，开关损耗也随之增大。为了减小开关损耗，\textbf{软开关技术}应运而生——\textbf{零电压开关（ZVS）}和\textbf{零电流开关（ZCS）}就是软开关的最基本形式。理论上讲采用软开关技术可使开关损耗降为零，可以提高效率，也使开关频率得以进一步提高，从而提高装置的功率密度。

\begin{table}[htbp]
\centering
\caption{电力电子技术发展阶段的粗略对照}
\label{tab:1-2}
\begin{tabular}{p{2.6cm}p{4.2cm}p{6.4cm}}
\toprule
\textbf{阶段} & \textbf{代表器件} & \textbf{技术特征} \\
\midrule
史前期（$\sim$1957 前） & 电子管、水银整流器、变流机组 & 旋转变流、静止汞弧整流；整流/逆变/周波变流理论已成熟 \\
晶闸管时代（1957 起） & 晶闸管（半控型） & 相控方式；关断依赖电网等外部条件；装置容量大、频率低 \\
全控型时代（70 年代后期起） & GTO、BJT、Power MOSFET & 门极可开通亦可关断；斩控（PWM）方式；开关频率提高 \\
复合型时代（80 年代后期起） & IGBT、MCT、IGCT & MOS 栅极 + 双极导电，驱动功率小且通态压降小；成为主导器件 \\
集成化与软开关 & 功率模块、PIC、ZVS/ZCS & 结构紧凑、功率密度高；开关损耗趋近于零 \\
\bottomrule
\end{tabular}
\end{table}

% ============================================================
\section{电力电子技术的应用}

电力电子技术的应用范围十分广泛。它不仅用于一般工业，也广泛用于交通运输、电力系统、通信系统、计算机系统、新能源系统等，在照明、空调等家用电器中也有广泛应用。

\subsection{一般工业}

工业中大量应用各种交直流电动机。直流电动机有良好的调速性能，为其供电的可控整流电源或直流斩波电源都是电力电子装置。近年来由于电力电子变频技术的迅速发展，交流电动机的调速性能可与直流电动机相媲美，交流调速技术逐渐大量应用并占据了主导地位。大至数千千瓦的轧钢机，小到几百瓦的数控机床伺服电动机，以及矿山牵引等场合，都广泛采用电力电子交流调速技术。一些调速性能要求不高的大型鼓风机也采用了变频装置以达到节能目的。还有些并不特别要求调速的电动机，为避免起动时的电流冲击而采用\textbf{软起动装置}，这也是电力电子装置。

\corepoint{电动机所消耗的电力甚至达到发电厂所发电力的 60\% 以上，以至于有人认为\textbf{电力传动是电力电子技术的“主战场”}。}

电化学工业大量使用直流电源，电解铝、电解食盐水等都需要大容量整流电源，电镀装置也需要整流电源。电力电子技术还大量用于冶金工业中的高频或中频感应加热电源、淬火电源及直流电弧炉电源等场合。

\subsection{交通运输}

电气化铁道中广泛采用电力电子技术：电气机车中的直流机车采用整流装置，交流机车采用变频装置；直流斩波器也广泛用于铁道车辆。在磁悬浮列车中，电力电子技术更是一项关键技术。除牵引电动机传动外，车辆中的各种辅助电源也都离不开电力电子技术。电动汽车的电动机依靠电力电子装置进行电力变换和驱动控制，蓄电池的充电也离不开电力电子装置。一台高级汽车中需要许多控制电动机，它们也要靠变频器和斩波器驱动并控制。航空和航海都离不开电力电子技术（电源和驱动）；电梯若算作交通运输，以前的电梯大都采用直流调速系统，近年来交流变频调速已成为主流。

\subsection{电力系统}

电力电子技术在电力系统中有着非常广泛的应用。据估计，发达国家在用户最终使用的电能中，有 60\% 以上的电能至少经过一次以上电力电子变流装置的处理。可以毫不夸张地说，如果离开电力电子技术，电力系统的现代化就是不可想象的。

\begin{itemize}
  \item \textbf{直流输电}：长距离、大容量输电时有很大优势，其送电端的整流阀和受电端的逆变阀都采用晶闸管变流装置；轻型直流输电则主要采用全控型的 IGBT 器件。
  \item \textbf{柔性交流输电（FACTS）}：依靠电力电子装置才得以实现。
  \item \textbf{无功补偿与谐波抑制}：晶闸管控制电抗器（TCR）、晶闸管投切电容器（TSC）都是重要的无功补偿装置；采用全控型器件（如 IGBT）的静止无功发生器（SVG）、有源电力滤波器（APF）等新型装置性能更为优越。
  \item \textbf{电能质量控制}：在配电网中，电力电子装置还可用于防止电网瞬时停电、瞬时电压跌落、闪变等，以改善供电质量。
  \item \textbf{变电所}：给操作系统提供可靠的交直流操作电源、给蓄电池充电等，都需要电力电子装置。
\end{itemize}

\subsection{电子装置用电源}

各种电子装置一般都需要不同电压等级的直流电源供电。通信设备中的程控交换机所用直流电源以前用晶闸管整流电源，现在已改为采用全控型器件的高频开关电源；大型计算机所需的工作电源、微型计算机内部的电源现在也都采用高频开关电源。以前大量采用的线性稳压电源，由于高频开关电源体积小、重量轻、效率高，现已逐渐被取代。可以说，信息技术装置都需要电力电子装置提供电源，\textbf{信息电子技术离不开电力电子技术}。在有大型计算机等场合，常常需要\textbf{不间断电源（UPS）}供电。

\subsection{家用电器与其他}

照明在家用电器中占有十分突出的地位。电力电子照明电源体积小、发光效率高、可节省大量能源，通常采用电力电子装置的光源被称为“节能灯”，正在逐步取代传统的白炽灯和荧光灯。\textbf{变频空调器}是家用电器中应用电力电子技术的典型例子；电视机、音响设备、家用计算机等电子设备的电源部分也都需要电力电子技术；不少洗衣机、电冰箱、微波炉中也应用了电力电子技术。

新能源及新型发电方式越来越受到重视，其中风力发电、太阳能发电发展最为迅速，燃料电池备受关注。太阳能发电和风力发电受环境的制约，发出的电力质量较差，常需要储能装置缓冲、需要改善电能质量，这就需要电力电子技术；当需要和电力系统联网时，更离不开电力电子技术。抽水储能发电站中大型电动机的起动和调速、超导储能的强大直流电源、核聚变反应堆的大容量脉冲电源，以及科学实验中的各种特种电源，都是电力电子技术的用武之地。

\corepoint{电力电子装置提供给负载的是各种不同的直流电源、恒频交流电源以及变频交流电源，因此也可以说，\textbf{电力电子技术研究的就是电源技术}。它在大型风机、水泵变频调速和照明电源方面的节能效果十分显著，因此也被称为\textbf{节能技术}。以前电力电子技术的应用偏重于中、大功率，现在在 1 kW 以下甚至几十瓦以下的功率范围内应用越来越广，这已成为一个重要的发展趋势。}

% ============================================================
\section{本教材的内容简介和使用说明}

本教材除第 1 章绪论外，可分为四大部分。

\begin{table}[htbp]
\centering
\caption{本教材的内容结构}
\label{tab:1-3}
\begin{tabular}{p{3.4cm}p{2.2cm}p{8.0cm}}
\toprule
\textbf{部分} & \textbf{章} & \textbf{内容与地位} \\
\midrule
器件 & 第 2 章、第 9 章 & 第 2 章是全书的基础，介绍各种器件的基本结构、工作原理、主要参数、应用特性，以晶闸管、IGBT、电力 MOSFET 三种应用最广的器件为重点；第 9 章是器件应用的共性问题，包括驱动、控制、保护以及串并联 \\
电路 & 第 3--6 章 & 各种基本电力变换电路：整流、逆变、直流—直流变流、交流—交流变流 \\
控制技术 & 第 7 章、第 8 章 & PWM 控制技术（没有新的电路拓扑，但使电路性能大为改善，适用于表~\ref{tab:1-1} 所列各类变换电路）；软开关技术（一般并不改变原来的基本电路，而是附加一些电路以实现软开关） \\
应用 & 第 10 章 & 电力电子装置的应用，体现第 5 版与第 4 版的重要差异 \\
\bottomrule
\end{tabular}
\end{table}

与前三版教材相比，第 4 版完成了以晶闸管电路为主体内容向以采用全控型器件的电路为主体内容的转移；第 5 版则在第 4 版问世 8 年来的新进展基础上做了充分反映。

学习建议：

\begin{itemize}
  \item 本课程要求先修“电路”和“电子技术基础”，最好也学过“自动控制理论”，并能较熟练地掌握示波器等电子仪器的使用方法。
  \item 电力电子技术实践性很强，实验在教学中有重要位置；教材正文后附“教学实验”部分，精选了五个最基本、有较高实用价值的实验。
  \item 术语在一门课程的学习中起着“纲”的作用，书末附有中英文“术语索引”，遇到拿不准的概念应回头查定义。
  \item 各章前有前言说明本章的地位及概要，大部分章末有小结，对全章要点和重点进行总结——复习时应以“小结核对 + 公式自推”为主，而不是背结论。
\end{itemize}

\corepoint{本章的考点集中在\textbf{概念辨析}：四大类电力变换的名称与英文缩写、半控型与全控型器件的分界、相控方式与斩控（PWM）方式的对应关系、“器件制造技术是基础、变流技术是核心”这一论断，以及电力电子器件\textbf{总工作在开关状态}这一与信息电子技术的根本区别。}

\chapter{电力电子器件}

就像晶体管和集成电路是模拟、数字电子电路的基础一样，\textbf{电力电子器件是电力电子电路的基础}。本章先概述器件的概念、特点与分类，再按“不可控器件$\rightarrow$半控型器件$\rightarrow$典型全控型器件$\rightarrow$其他新型器件”的顺序，逐个介绍其工作原理、基本特性、主要参数以及选择和使用中应注意的问题。器件应用中共性的驱动、保护、串并联问题留到第 9 章。

% ============================================================
\section{电力电子器件概述}

\subsection{电力电子器件的概念和特征}

在电气设备或电力系统中，\textbf{直接承担电能的变换或控制任务的电路称为主电路}（Main Power Circuit）。\textbf{电力电子器件是指可直接用于处理电能的主电路中，实现电能的变换或控制的电子器件}。

广义上电力电子器件也可分为电真空器件和半导体器件两类。但自 20 世纪 50 年代以来，除了在频率很高（如微波）的大功率高频电源中还在使用真空管外，基于半导体材料的电力电子器件已逐步取代汞弧整流器、闸流管等电真空器件，成为电能变换和控制领域的绝对主力，因此电力电子器件也往往\textbf{专指电力半导体器件}，主要材料仍是硅。

同处理信息的电子器件相比，电力电子器件具有如下特征：

\begin{enumerate}
  \item \textbf{所能处理电功率的大小（即承受电压和电流的能力）是其最重要的参数}。其处理电功率的能力小至毫瓦级、大至兆瓦级，一般都远大于处理信息的电子器件。
  \item 因为处理的电功率较大，为了减小本身的损耗、提高效率，电力电子器件一般都\textbf{工作在开关状态}。导通时（通态）阻抗很小、接近于短路，管压降接近于零，而电流由外电路决定；阻断时（断态）阻抗很大、接近于断路，电流几乎为零，而管子两端电压由外电路决定——就像普通晶体管的饱和与截止状态一样。因而器件的\textbf{动态特性（开关特性）和参数}也是其特性很重要的方面，有时甚至上升为第一位。
  \item 实际应用中电力电子器件往往需要由信息电子电路来控制。由于器件处理的电功率较大，普通信息电子电路的信号一般不能直接控制其导通或关断，需要中间电路对信号加以放大，这就是\textbf{驱动电路}。
  \item 尽管工作在开关状态，器件自身的功率损耗通常仍远大于信息电子器件，因而封装上讲究散热设计，工作时一般还需安装散热器。损耗包括\textbf{通态损耗}（通态压降 $\times$ 通态电流）、\textbf{断态损耗}（断态漏电流 $\times$ 断态电压）、\textbf{开关损耗}（开通损耗与关断损耗）。断态漏电流极其微小，故\textbf{通态损耗是主要成因}；当开关频率较高时，开关损耗会增大而可能成为主要因素。某些器件驱动电路注入的功率也是发热原因之一。
\end{enumerate}

\corepoint{把一个电力电子器件、或者外特性像一个开关的几个电力电子器件的组合，称为\textbf{电力电子开关}（电力半导体开关）；做电路分析时为简单起见往往用\textbf{理想开关}代替。}

\subsection{应用电力电子器件的系统组成}

电力电子器件在实际应用中，一般由\textbf{控制电路}、\textbf{驱动电路}和以电力电子器件为核心的\textbf{主电路}组成一个系统：由信息电子电路组成的控制电路按系统工作要求形成控制信号，通过驱动电路去控制主电路中器件的导通或关断，完成整个系统的功能。因此从宏观角度讲，电力电子电路也被称为\textbf{电力电子系统}。有的系统还需要\textbf{检测电路}采样主电路或现场信号；广义上把检测电路和驱动电路这些主电路以外的电路都归为控制电路，从而粗略地说，\textbf{电力电子系统由主电路和控制电路组成}。

主电路中的电压电流一般较大，而控制电路的元器件只能承受较小的电压电流，因此在主电路与控制电路连接的路径上（如驱动电路与主电路的连接处、驱动电路与控制信号的连接处、主电路与检测电路的连接处）一般需要\textbf{电气隔离}，通过光、磁等手段传递信号。此外，主电路中往往有电压电流过冲，而电力电子器件一般比主电路中普通元件昂贵、承受过电压过电流的能力却差一些，因此在主电路和控制电路中附加\textbf{保护电路}也往往非常必要。

电力电子器件一般都有三个端子：其中两个连接在主电路中、流通主电路电流，第三个称为\textbf{控制端（控制极）}。器件的导通或关断通过在控制端和一个主电路端子之间施加一定信号来控制，这个主电路端子是驱动电路和主电路的公共端，一般是主电路电流\textbf{流出}器件的那个端子。

\subsection{电力电子器件的分类}

\begin{table}[htbp]
\centering
\caption{电力电子器件的分类（按控制程度、驱动信号性质、载流子）}
\label{tab:2-1}
\begin{tabular}{p{2.7cm}p{4.8cm}p{6.5cm}}
\toprule
\textbf{分类依据} & \textbf{类别} & \textbf{代表器件与说明} \\
\midrule
\multirow{3}{2.4cm}{控制信号对通断的控制程度} & 不可控器件 & 电力二极管：不能用控制信号控制通断，因而也不需要驱动电路，只有两个端子；导通和关断完全由其在主电路中承受的电压电流决定 \\
 & 半控型器件 & 晶闸管及其大部分派生器件：控制信号可以控制其导通而\textbf{不能}控制其关断，关断完全由主电路承受的电压电流决定 \\
 & 全控型器件（自关断器件） & 通过控制信号\textbf{既可控制导通又可控制关断}；最常用的是 IGBT 和电力 MOSFET，此外还有 GTO、GTR 等 \\
\midrule
\multirow{2}{2.4cm}{驱动信号的性质} & 电流驱动型 & 通过从控制端注入或抽出电流来实现通断控制（晶闸管、GTO、GTR） \\
 & 电压驱动型（场控器件） & 仅通过在控制端和公共端之间施加一定电压信号即可控制通断；实质是用控制端电压在两个主电路端子间产生可控电场来改变电流，故又称场效应器件（电力 MOSFET、IGBT） \\
\midrule
\multirow{2}{2.4cm}{有效信号的波形} & 脉冲触发型 & 施加一个电压或电流\textbf{脉冲}即可开通或关断，此后只要主电路条件不变就能维持状态，不必继续施加控制信号（晶闸管、GTO） \\
 & 电平控制型 & 必须\textbf{持续}施加一定电平的信号才能开通并维持导通（或关断并维持阻断）（电力 MOSFET、IGBT） \\
\midrule
\multirow{3}{2.4cm}{导电载流子} & 单极型（多子器件） & 只有一种载流子参与导电：电力 MOSFET、SIT \\
 & 双极型（少子器件） & 电子和空穴两种载流子都参与导电：电力二极管、晶闸管、GTO、GTR \\
 & 复合型（混合型） & 由单极型与双极型器件集成混合而成：IGBT、MCT \\
\bottomrule
\end{tabular}
\end{table}

\subsection{本章内容和学习要点}

学习和选用电力电子器件时，最重要的是\textbf{掌握其基本特性}；此外应了解各国、各厂家对各种器件具体型号的命名方法，特别是要了解每种器件各个主要参数和特性曲线的意义。\textbf{掌握器件的型号命名法，以及其参数和特性曲线的使用方法，是在实际中正确应用电力电子器件的两个基本要求}。

许多电力电子器件都有其相对应的用于处理信息的电子器件（电力二极管、GTR、电力 MOSFET 分别与信息电子中的二极管、双极型晶体管、场效应晶体管对应），从半导体物理结构和工作原理上讲基本相同；但为了承受高电压和大电流，电力电子器件又具有与其信息电子对应者不同之处——这些不同之处，正是本章要抓的重点。

% ============================================================
\section{不可控器件——电力二极管}

电力二极管（Power Diode）自 20 世纪 50 年代初期就获得应用，当时也称为半导体整流器（SR），并已开始逐步取代汞弧整流器。虽然是不可控器件，但其结构和原理简单、工作可靠，所以直到现在仍大量应用于许多电气设备中；在采用全控型器件的电路中电力二极管往往不可缺少，特别是开通和关断速度很快的\textbf{快恢复二极管}和\textbf{肖特基二极管}具有不可替代的地位。

\subsection{PN 结与电力二极管的工作原理}

电力二极管实际上是由一个面积较大的 PN 结、两端引线及封装组成的。N 型半导体和 P 型半导体结合后构成 PN 结：交界处电子和空穴的浓度差别造成多子向另一区移动的\textbf{扩散运动}，在界面两侧留下带正、负电荷但不能任意移动的杂质离子，称为\textbf{空间电荷}；空间电荷建立的电场称为\textbf{内电场（自建电场）}，其方向阻止扩散运动，另一方面又吸引对方区内的少子向本区运动，这就是\textbf{漂移运动}。扩散与漂移相互联系又相互矛盾，最终达到动态平衡，形成稳定的\textbf{空间电荷区}（也称耗尽层、阻挡层或势垒区）。

\begin{itemize}
  \item \textbf{正向偏置}（外加电压正端接 P 区、负端接 N 区）：外加电场与自建电场方向\textbf{相反}，多子的扩散运动大于少子的漂移运动，空间电荷区变窄，外电路形成自 P 区流入、从 N 区流出的\textbf{正向电流} $I_F$——这是 PN 结的正向导通状态。
  \item \textbf{反向偏置}：外加电场与自建电场方向\textbf{相同}，少子的漂移运动大于多子的扩散运动，空间电荷区变宽，外电路形成自 N 区流入、从 P 区流出的\textbf{反向电流} $I_R$。少子浓度很小，温度一定时其数值趋于恒定，称为\textbf{反向饱和电流}，一般仅微安数量级，故反向偏置的 PN 结表现为高阻态——反向截止状态。
\end{itemize}

\textbf{PN 结的单向导电性}就是二极管的基本原理。为了建立承受高电压和大电流的能力，电力二极管的物理结构与信息电子二极管有如下不同：

\begin{enumerate}
  \item \textbf{垂直导电结构}：电流在硅片内流动的总体方向与硅片表面\textbf{垂直}（信息电子二极管一般是横向导电结构），使硅片中通过电流的有效面积增大，可显著提高通流能力。
  \item \textbf{P-i-N 结构}：在 P 区和 N 区之间多了一层\textbf{低掺杂 N 区}（记为 N$^-$），也称为\textbf{漂移区}（i 区）。低掺杂 N 区掺杂浓度低而接近本征半导体，可以承受很高的电压而不致被击穿——\textbf{低掺杂 N 区越厚，能承受的反向电压越高}。
  \item \textbf{电导调制效应}：低掺杂 N 区高电阻率对正向导通不利，这个矛盾通过电导调制效应解决。正向电流较小时，二极管电阻主要是作为基片的低掺杂 N 区的欧姆电阻，阻值较高且为常量，管压降随正向电流上升而增加；正向电流较大时，由 P 区注入并积累在低掺杂 N 区的少子空穴浓度将很大，为维持电中性，其多子浓度也相应大幅增加，使其电阻率明显下降（电导率大大增加）——这就是\textbf{电导调制效应}。它使电力二极管在正向电流较大时压降仍然很低，所以正向偏置的电力二极管表现为\textbf{低阻态}。
\end{enumerate}

PN 结具有一定的反向耐压能力，但当反向电压过大时反向电流将急剧增大，破坏反向偏置为截止的工作状态，这称为\textbf{反向击穿}，按机理分为\textbf{雪崩击穿}和\textbf{齐纳击穿}两种形式。反向击穿发生时，只要外电路采取措施将反向电流限制在一定范围内，则反向电压降低后 PN 结仍可恢复原来的状态；若反向电流未被限制住，使反向电流与反向电压的乘积超过 PN 结允许的耗散功率，就会因热量散发不出去导致结温上升直至过热烧毁，这就是\textbf{热击穿}。

PN 结中的电荷量随外加电压而变化，呈现电容效应，称为\textbf{结电容} $C_J$（微分电容），分为\textbf{势垒电容} $C_B$ 与\textbf{扩散电容} $C_D$：势垒电容只在外加电压变化时才起作用，其大小与 PN 结截面积成正比、与阻挡层厚度成反比；扩散电容仅在正向偏置时起作用。正向电压较低时以势垒电容为主，较高时扩散电容成为主要成分。结电容影响 PN 结的工作频率，高速开关状态下可能使其单向导电性变差甚至不能工作。

\subsection{电力二极管的基本特性}

\subsubsection*{1. 静态特性}

静态特性主要指\textbf{伏安特性}（图~\ref{fig:2-1}）。当承受的正向电压大到一定值（\textbf{门槛电压} $U_{TO}$）时，正向电流才开始明显增加，处于稳定导通状态；与正向电流对应的器件两端电压 $U_F$ 即其\textbf{正向电压降}。承受反向电压时，只有少子引起的微小而数值恒定的反向漏电流。

\begin{figure}[htbp]
\centering
\begin{tikzpicture}[font=\footnotesize,scale=0.95]
  \draw[->] (-4.2,0) -- (4.4,0) node[below] {$u$};
  \draw[->] (0,-3.1) -- (0,3.4) node[left] {$i$};
  \draw[yhbAccent,thick] (-4.0,-0.12) -- (-2.6,-0.12);
  \draw[yhbAccent,thick] (-2.6,-0.12) .. controls (-2.3,-0.9) and (-2.1,-2.6) .. (-1.9,-3.0);
  \draw[yhbAccent,thick] (0,0) -- (1.9,0.02);
  \draw[yhbAccent,thick] (1.9,0.02) .. controls (2.05,0.5) and (2.3,1.6) .. (2.6,3.0);
  \node[anchor=west] at (2.65,3.0) {$U_{TO}$};
  \node[anchor=west,text=textgray] at (-4.05,-0.45) {反向漏电流};
  \node[anchor=south,text=textgray] at (2.0,0.15) {门槛电压};
  \node[anchor=north] at (-2.4,0.1) {反向击穿};
  \draw[->,thin,textgray] (-1.2,-0.6) -- (-1.85,-0.9);
\end{tikzpicture}
\caption{电力二极管的伏安特性}
\label{fig:2-1}
\end{figure}

\subsubsection*{2. 动态特性}

因为结电容的存在，电力二极管在零偏置、正向偏置和反向偏置三种状态之间转换时必然经历一个过渡过程；过渡过程中 PN 结的一些区域需要一定时间调整带电状态，其电压—电流特性不能用伏安特性描述而是随时间变化的——这就是\textbf{动态特性}，并且往往专指反映通态与断态之间转换过程的\textbf{开关特性}。

\corepoint{{\bfseries 反向恢复过程}：原处于正向导通状态的电力二极管，外加电压突然由正向变为反向时，并\textbf{不能立即关断}，而需经过一段短暂时间才能重新获得反向阻断能力。关断之前有较大的反向电流出现，并伴随明显的反向电压过冲。原因是正向导通时在 PN 结两侧储存的大量少子需要被清除掉才能达到反向偏置稳态。}

设 $t=0$ 时刻外加电压突然由正向变为反向，正向电流在此反向电压作用下开始下降，下降的速率由主电路电感决定；到 $t=t_d$（延迟时间）时电流过零并开始反向增大，到 $t=t_f$ 时刻反向电流达到最大值 $I_{RP}$ 后开始减小，$t_f$ 称为电流下降时间。$t_d$ 与 $t_f$ 之和即\textbf{反向恢复时间} $t_{rr}$。其中 $t_d$ 内电荷存储效应起主要作用，$t_f$ 内则主要取决于结电容的充放电。

\keyformula{t_{rr}=t_d+t_f,\qquad Q_{rr}\approx \tfrac{1}{2}I_{RP}\,t_{rr}}

按反向恢复特性，二极管分为\textbf{软恢复特性}（反向电流下降较缓慢）与\textbf{硬恢复特性}（下降很陡、易引起较大电压过冲）两类。由零偏置转为正向偏置时也有开通过程（正向恢复过程），伴随电压过冲，但一般比反向恢复过程短得多、影响小。

\subsection{电力二极管的主要参数}

\begin{table}[htbp]
\centering
\caption{电力二极管的主要参数}
\label{tab:2-2}
\begin{tabular}{p{3.6cm}p{11.4cm}}
\toprule
\textbf{参数} & \textbf{含义与用法} \\
\midrule
正向平均电流 $I_{F(AV)}$ & 在规定的管壳温度和散热条件下，允许流过的\textbf{最大工频正弦半波电流的平均值}。定额按发热效应（即有效值相等）等效，正弦半波电流的\textbf{平均值:有效值 $=1:1.57$}，故实际波形电流有效值 $I$ 对应 $I_{F(AV)}=I/1.57$，选用时再留 1.5$\sim$2 倍裕量 \\
正向压降 $U_F$ & 在规定温度下流过规定电流时器件两端的正向电压降 \\
反向重复峰值电压 $U_{RRM}$ & 允许重复施加的反向最高峰值电压，通常为反向不重复峰值电压 $U_{RSM}$ 的 $90\%$；使用时取正常工作峰值电压的 2$\sim$3 倍 \\
反向漏电流 $I_R$ & 在额定结温、施加反向重复峰值电压时的反向电流 \\
最高工作结温 $T_{jm}$ & 通常为 $125\sim175\,^\circ$C，取决于器件封装与散热条件 \\
浪涌电流 $I_{FSM}$ & 所能承受的最大的连续一个或几个工频周期的过电流 \\
反向恢复时间 $t_{rr}$ & 关断过程中重新获得反向阻断能力所需时间 \\
\bottomrule
\end{tabular}
\end{table}

\subsection{电力二极管的主要类型}

\begin{itemize}
  \item \textbf{普通二极管}（Standard Recovery Diode）：又称整流二极管，多用于开关频率不高（1 kHz 以下）的整流电路，其反向恢复时间较长、在 $5\,\mu$s 以上，正向电流定额和反向电压定额却可以达到很高。
  \item \textbf{快恢复二极管}（Fast Recovery Diode，FRD）：反向恢复时间很短（$5\,\mu$s 以下），正向压降较低（约 0.9 V），多用于开关频率较高的电路，如与全控型器件配合的续流二极管。从性能上分为\textbf{快速恢复}和\textbf{超快恢复}两种。
  \item \textbf{肖特基二极管}（Schottky Barrier Diode，SBD）：以金属—半导体结代替 PN 结，反向恢复时间极短（几十纳秒），正向压降很低（0.4$\sim$0.6 V）；但反向耐压较低（一般不超 100 V 左右）、反向漏电流较大，多用于低电压、高频场合。
\end{itemize}

% ============================================================
\section{半控型器件——晶闸管}

晶闸管（Thyristor）是晶体闸流管的简称，又称可控硅整流器（SCR），是目前\textbf{容量最大、应用最广泛}的电力电子器件之一。

\subsection{晶闸管的结构与工作原理}

晶闸管是大功率四层（P1-N1-P2-N2）三端（阳极 A、阴极 K、门极 G）器件，内部有三个 PN 结 J$_1$、J$_2$、J$_3$。从外形看主要有螺栓形、平板形等封装。

\begin{center}
\begin{circuitikz}[font=\footnotesize,scale=0.9]
  \draw (0,0) node[left] {A} to[short,o-] (1,0) to[thyristor] (3.4,0) to[short,-o] (4.4,0) node[right] {K};
  \draw (2.55,-0.95) -- (1.95,-1.55) node[below] {G};
\end{circuitikz}
\end{center}

\textbf{工作原理（双晶体管等效）}：把晶闸管看成 P1-N1-P2 构成的 PNP 晶体管 V$_1$ 与 N1-P2-N2 构成的 NPN 晶体管 V$_2$ 的互联——V$_1$ 的集电极即 V$_2$ 的基极，V$_2$ 的集电极即 V$_1$ 的基极。设两个晶体管的共基极电流放大系数分别为 $\alpha_1$、$\alpha_2$，则阳极电流

\keyformula{I_A=\frac{\alpha_2 I_G+\left(\alpha_1+\alpha_2\right)I_{CBO}}{1-\left(\alpha_1+\alpha_2\right)}}

\corepoint{晶体管的特性是：\textbf{低发射极电流下 $\alpha$ 很小，发射极电流增大时 $\alpha$ 迅速增大}。因此当门极注入电流 $I_G$ 后，两个晶体管相互正反馈，$\alpha_1+\alpha_2$ 迅速趋近于 1，$I_A$ 急剧增大，晶闸管\textbf{饱和导通}。导通后 $\alpha_1+\alpha_2\approx1$，即使门极电流消失，晶闸管仍维持导通——\textbf{门极失去控制作用}。要使晶闸管关断，只有设法使阳极电流减小到维持电流以下，使 $\alpha_1+\alpha_2<1$。}

由此得到晶闸管导通与关断的条件：

\begin{table}[htbp]
\centering
\caption{晶闸管的导通与关断条件}
\label{tab:2-3}
\begin{tabular}{p{2.2cm}p{12.8cm}}
\toprule
\textbf{状态} & \textbf{条件} \\
\midrule
由断态转通态 & （1） 阳极电位高于阴极（\textbf{承受正向阳极电压}）；（2） 门极承受\textbf{正向门极电压并有足够的门极电流} $I_G$。两者\textbf{缺一不可} \\
维持通态 & 阳极电流大于\textbf{维持电流} $I_H$；此时门极信号已不起作用 \\
由通态转断态 & 阳极电流减小到维持电流以下（或阳极电压反向）。方法是：主电路电流过零、外加反向电压、或电路阻抗增大使电流降到 $I_H$ 以下 \\
\bottomrule
\end{tabular}
\end{table}

\subsection{晶闸管的基本特性}

\subsubsection*{1. 静态特性}

\begin{itemize}
  \item \textbf{正向特性}：门极断路（$I_G=0$）时，逐渐增大正向阳极电压，只有很小的正向漏电流，晶闸管处于\textbf{正向阻断状态}；当正向电压增大到\textbf{正向转折电压} $U_{BO}$ 时，漏电流急剧增大，由断态转入通态（这种靠过电压导通的方式叫“硬开通”，会损坏器件，正常不用）。门极施加电流 $I_G$ 后，正向转折电压随 $I_G$ 增大而降低；晶闸管导通后，通态压降很小（约 1 V 左右），流过很大的阳极电流。
  \item \textbf{反向特性}：晶闸管承受反向电压时，只有极小的反向漏电流，处于\textbf{反向阻断状态}；反向电压增大到一定程度时反向电流急剧增大，即\textbf{反向击穿}，器件损坏。
\end{itemize}

\begin{center}
\begin{tikzpicture}[font=\footnotesize,scale=0.95]
  \draw[->] (-4.0,0) -- (4.6,0) node[below] {$u$};
  \draw[->] (0,-2.9) -- (0,3.5) node[left] {$i$};
  % 正向阻断段
  \draw[yhbAccent,thick] (0,0) -- (2.75,0.05);
  \draw[yhbAccent,thick] (2.75,0.05) .. controls (2.95,0.7) and (3.0,1.6) .. (3.05,3.1);
  % 通态段
  \draw[yhbAccent,thick] (2.75,0.05) .. controls (2.2,1.1) and (1.6,2.2) .. (1.45,3.1);
  % 反向
  \draw[yhbAccent,thick] (-3.6,-0.08) -- (-2.9,-0.1);
  \draw[yhbAccent,thick] (-2.9,-0.1) .. controls (-2.75,-0.7) and (-2.65,-1.9) .. (-2.6,-2.7);
  \node[anchor=west] at (3.1,3.05) {$U_{BO}$};
  \node[anchor=east] at (1.4,3.05) {通态};
  \node[anchor=west] at (2.78,1.9) {\small $I_G$ 增大};
  \node[anchor=north,textgray] at (-2.55,-0.15) {\small 反向击穿};
  \node[anchor=south,textgray] at (2.2,0.2) {\small 正向阻断};
\end{tikzpicture}
\end{center}

\subsubsection*{2. 动态特性}

\begin{itemize}
  \item \textbf{开通过程}：门极加触发脉冲后，阳极电流滞后一段\textbf{延迟时间} $t_d$（约零点几到 1.5 $\mu$s）才明显上升，随后是\textbf{上升时间} $t_r$（约零点几到 3 $\mu$s）；\textbf{开通时间} $t_{gt}=t_d+t_r$。开通时须注意\textbf{通态电流临界上升率 di/dt}：上升太快则刚导通时电流集中在门极附近的小区域，局部过热而损坏器件（须串电感或用强触发）。
  \item \textbf{关断过程}：先把阳极电流降到零并施加反压，器件内部存储电荷需要清除，存在\textbf{反向阻断恢复时间} $t_{rr}$；之后还要经过\textbf{正向阻断恢复时间} $t_{gr}$ 才恢复正向阻断能力。两者之和即\textbf{关断时间}

\keyformula{t_q=t_{rr}+t_{gr}}

  在正向阻断恢复时间内若重新施加正向电压，晶闸管会重新正向导通而不是由门极控制导通，故实际应用中须施加足够长时间的反压使其充分恢复阻断能力。
\end{itemize}

\corepoint{{\bfseries 断态电压临界上升率 du/dt}：在额定结温和门极开路情况下，不导致晶闸管从断态到通态转换的外加电压最大上升率。正向电压上升过快时，处于阻断状态的 J$_2$ 结（相当于一个电容）会有充电\textbf{位移电流}流过，流经 J$_3$ 时起类似门极触发电流的作用，\textbf{电压上升率过大会使晶闸管误导通}。使用中实际电压上升率必须低于此临界值（可加 RC 吸收电路）。}

\subsection{晶闸管的主要参数}

普通晶闸管在反向稳态下一定是阻断状态；而正向工作时既可能导通也可能阻断，因此提到参数时，\textbf{断态和通态都是为了区分正向的不同状态，“正向”二字可省去}。各项主要参数往往与结温相联系。

\subsubsection*{1. 电压定额}

\begin{itemize}
  \item \textbf{断态重复峰值电压} $U_{DRM}$：门极断路、结温为额定值时，允许重复加在器件上的正向峰值电压（国标规定重复频率 50 Hz、每次持续时间不超过 10 ms）。规定 $U_{DRM}$ 为\textbf{断态不重复峰值电压}（断态最大瞬时电压）的 $90\%$，后者应低于正向转折电压。
  \item \textbf{反向重复峰值电压} $U_{RRM}$：同样条件下允许重复加在器件上的反向峰值电压，为反向不重复峰值电压的 $90\%$，后者应低于反向击穿电压。
  \item 通常取 $U_{DRM}$ 与 $U_{RRM}$ 中\textbf{较小}的标称值作为该器件的\textbf{额定电压}；选用时额定电压要留裕量，一般取为正常工作时器件承受峰值电压的\textbf{2$\sim$3 倍}。
  \item \textbf{通态（峰值）电压} $U_{TM}$：通以某一规定倍数额定通态平均电流时的通态峰值电压。
\end{itemize}

\subsubsection*{2. 电流定额}

\begin{itemize}
  \item \textbf{通态平均电流} $I_{T(AV)}$：在环境温度 40 $^\circ$C 和规定冷却状态下，稳定结温不超过额定结温时所允许流过的\textbf{最大工频正弦半波电流的平均值}，也是标称其额定电流的参数。与电力二极管一样，该参数按发热效应（有效值相等）等效，正弦半波\textbf{平均值:有效值 $=1:1.57$}，故先由实际波形有效值算出 $I_{T(AV)}=I/1.57$，再取\textbf{1.5$\sim$2 倍}裕量。
  \item \textbf{维持电流} $I_H$：使晶闸管维持导通所必需的最小电流，一般为几十到几百毫安；结温越高 $I_H$ 越小。
  \item \textbf{擎住电流} $I_L$：晶闸管刚从断态转入通态并\textbf{移除触发信号后}能维持导通所需的最小电流。对同一器件通常 $I_L\approx(2\sim4)I_H$。
  \item \textbf{浪涌电流} $I_{FSM}$：电路异常情况引起的、使结温超过额定结温的不重复性最大正向过载电流，可作为设计保护电路的依据。
\end{itemize}

\corepoint{选型步骤记牢三步：（1） 由实际波形求电流\textbf{有效值} $I$；（2） 折算到正弦半波平均值 $I_{T(AV)}=I/1.57$；（3） 乘 1.5$\sim$2 倍裕量后取标准档。电压则按承受的\textbf{峰值}乘 2$\sim$3 倍。}

\subsubsection*{3. 动态参数}

除开通时间 $t_{gt}$ 和关断时间 $t_q$ 外，还有\textbf{断态电压临界上升率 du/dt} 与\textbf{通态电流临界上升率 di/dt}。普通晶闸管 $t_q$ 约为数百微秒，快速晶闸管为数十微秒，高频晶闸管为 10 $\mu$s 左右。

\subsection{晶闸管的派生器件}

\begin{itemize}
  \item \textbf{快速晶闸管}（FST）：包括所有专为快速应用而设计的晶闸管，有常规快速晶闸管和高频晶闸管，可分别应用于 400 Hz 和 10 kHz 以上的斩波或逆变电路。其开关时间以及 du/dt、di/dt 耐量都明显改善；不足是电压和电流定额都不易做高，且因工作频率较高，选择通态平均电流时\textbf{不能忽略开关损耗的发热效应}。
  \item \textbf{双向晶闸管}（TRIAC）：可认为是一对反并联普通晶闸管的集成，有两个主电极 T$_1$、T$_2$ 和一个门极 G。门极使器件在主电极的正反两个方向均可触发导通，因此在第 I、第 III 象限有对称的伏安特性。与一对反并联晶闸管相比更经济、控制电路简单，广泛用于\textbf{交流调压电路}、固态继电器（SSR）和交流电动机调速。
  \item \textbf{逆导晶闸管}（RCT）：将晶闸管和二极管反并联集成在同一管芯上，其正向压降小、关断时间短、高温特性好、额定结温高；用于不需要阻断反向电压的电路中（如斩波、逆变中的续流）。
  \item \textbf{光控晶闸管}（LTT）：用光信号代替电信号触发，门极触发改为光触发，具有主电路与控制电路间天然电气隔离、抗干扰能力强的优点，多用于高压直流输电等超高压大容量场合。
\end{itemize}

% ============================================================
\section{典型全控型器件}

20 世纪 70 年代后期，以 GTO、GTR 和电力 MOSFET 为代表的全控型器件迅速发展，使电力电子技术的面貌焕然一新。本节按 GTO $\rightarrow$ GTR $\rightarrow$ 电力 MOSFET $\rightarrow$ IGBT 的顺序介绍。

\subsection{门极可关断晶闸管（GTO）}

GTO（Gate Turn-Off Thyristor）是晶闸管的一种派生器件，但可以通过在门极施加负的脉冲电流使其关断，因而属于\textbf{全控型器件}。其电压、电流容量较大，与普通晶闸管接近，在兆瓦级以上的大功率场合仍有较多应用。

\begin{itemize}
  \item \textbf{结构特点}：与普通晶闸管的 PNPN 四层结构相同，但\textbf{多元集成结构}——一个 GTO 由数百甚至上千个小 GTO 单元并联而成，阴极条做得细长，以便门极电流能迅速均匀地流过整个阴极面积。门极为横向结构、阴极条窄而长，是它能实现门极关断的结构基础。
  \item \textbf{工作原理}：GTO 导通时 $\alpha_1+\alpha_2$ 比 1 略大但很接近 1（普通晶闸管则远大于 1 进入饱和），处于\textbf{临界饱和}状态，这为门极关断创造了条件。关断时门极加负电流 $I_{GM}$ 抽出存储电荷，使 $\alpha_1+\alpha_2<1$，阳极电流下降并关断。
  \item \textbf{电流关断增益} $\beta_{off}$：最大可关断阳极电流 $I_{ATO}$ 与门极负电流最大值 $I_{GM}$ 之比（注：此处 $I_{ATO}$ 即手册中的 $I_{ATM}$，习惯写作 $I_{ATO}$）
\end{itemize}

\keyformula{\beta_{off}=\frac{I_{ATO}}{I_{GM}}}

$\beta_{off}$ 一般很小，只有 5 左右，这是 GTO 的一个主要缺点——需要很大的门极负脉冲电流才能关断，驱动功率大。

\subsection{电力晶体管（GTR）}

电力晶体管（Giant Transistor，GTR）是一种双极型大功率晶体管，通常采用\textbf{共发射极接法}，由基极电流控制集电极电流。它是双极型电流驱动器件，有电导调制效应，通态压降较低。

\begin{itemize}
  \item \textbf{输出特性}：明显分为截止区、放大区和饱和区。作为开关使用时，工作在\textbf{截止区（断态）}和\textbf{饱和区（通态）}，放大区只是过渡。饱和时管压降很低（$U_{CE(sat)}$），截止时只有极小的漏电流。
  \item \textbf{主要参数}：$U_{CEO}$（$I_B=0$ 时集电极—发射极击穿电压）等一组击穿电压、集电极最大电流 $I_{CM}$、集电极最大耗散功率 $P_{CM}$、共射极电流放大系数 $h_{FE}$（开通时是 $h_{FE}$，关断时须注意存储时间）。
  \item \textbf{二次击穿}：GTR 特有的失效机理。当集电极电压较高、电流较大时，局部热点使电流更集中于该点，形成正反馈，导致器件在远低于额定功耗的条件下瞬间损坏。因此 GTR 有\textbf{安全工作区（SOA）}，分正向偏置 SOA 与反向偏置 SOA，使用时必须保证工作点落在 SOA 之内，并加\textbf{缓冲电路}把开关轨迹限制在安全区内。
  \item GTR 的开关时间在几微秒量级，比晶闸管短得多，但仍远长于电力 MOSFET；驱动电流大、驱动功率大是其主要缺点，现已逐步被 IGBT 取代。
\end{itemize}

\subsection{电力场效应晶体管（Power MOSFET）}

电力 MOSFET 是小功率 MOSFET 的功率化，属于\textbf{单极型电压驱动器件}，只靠\textbf{多子}导电，因而开关速度极快。

\begin{itemize}
  \item \textbf{结构}：采用\textbf{垂直导电结构}（VVMOS、VDMOS），有 N 沟道和 P 沟道两种，以 N 沟道（N-MOSFET）为主；内部\textbf{源极与衬底短接}，并内含反并联的体二极管。由于没有少子存储效应，其开关过程只是对输入电容（$C_{iss}$）的充放电。
  \item \textbf{静态特性}：输出特性分为\textbf{截止区、饱和区（恒流区）、可变电阻区}。在开关应用中工作在截止区与可变电阻区之间切换；导通时表现为一个\textbf{通态电阻} $R_{on}$（$R_{DS(on)}$），因此通态压降随电流线性增大，这是它在高压大电流场合通态损耗偏大的原因。
  \item \textbf{转移特性}：漏极电流 $I_D$ 与栅源电压 $U_{GS}$ 的关系，$U_{GS}$ 超过\textbf{开启电压} $U_{GS(th)}$ 后才出现 $I_D$；跨导 $g_{fs}=\Delta I_D/\Delta U_{GS}$ 表征其放大能力（作开关用时关心的是 $R_{on}$ 与栅极电荷）。
  \item \textbf{动态特性}：开通、关断过程即栅极电容充放电过程，开关时间在几十到几百纳秒量级；开关损耗与驱动电阻、栅极电荷有关。
\end{itemize}

\corepoint{电力 MOSFET 的三大优点：（1） \textbf{驱动功率小}（电压驱动、输入阻抗高、静态几乎不耗驱动功率）；（2） \textbf{开关速度极快}（无少子存储效应）；（3） \textbf{无二次击穿}，且 $R_{on}$ 具有\textbf{正温度系数}，多管并联时能自动均流。主要缺点：\textbf{通态压降大}（耐压越高 $R_{on}$ 越大，故耐压不易做高），一般用于 1000 V 以下、高频场合。}

\subsection{绝缘栅双极晶体管（IGBT）}

GTR 和 GTO 是双极型电流驱动器件，电流驱动、驱动功率大、开关速度较低；电力 MOSFET 是单极型电压驱动器件，驱动功率小、速度快但通态压降大。IGBT 把二者的优点结合起来，成为\textbf{现代电力电子技术的主导器件}。

\begin{itemize}
  \item \textbf{结构}：由沟道型 VDMOSFET 与双极型晶体管组合而成的复合器件，相当于在 MOSFET 的漏极侧再加一层 P$^+$ 区，形成四层结构。对外仍是三个电极：\textbf{栅极 G、集电极 C、发射极 E}（对应 MOSFET 的 G、D、S）。
  \item \textbf{特性}：\textbf{电压驱动}（栅极输入阻抗高、驱动功率小），\textbf{通态压降小}（因为 P$^+$ 区向漂移区注入少子，产生电导调制效应）；转移特性与输出特性与 MOSFET 相似，输出特性同样有截止区、饱和区、可变电阻区，但饱和区（恒流区）较窄。
  \item \textbf{参数}：$U_{CES}$（集电极—发射极最高电压）、$I_C$（额定集电极电流）、$U_{CE(sat)}$（通态饱和压降）、$U_{GE(th)}$（开启电压，一般 3$\sim$6 V）、开关时间等。
  \item \textbf{掣住效应（latch-up）}：IGBT 内部天然存在一个寄生晶闸管（PNPN），当集电极电流过大或关断过快时寄生晶闸管可能被触发导通，栅极失去控制作用，导致器件损坏。这是 IGBT 使用中必须避免的失效模式，须限制最大集电极电流、必要时加负栅压关断。
\end{itemize}

\begin{table}[htbp]
\centering
\caption{几种典型电力电子器件的对比}
\label{tab:2-4}
\begin{tabular}{p{2.5cm}p{2.5cm}p{2.4cm}p{2.4cm}p{3.6cm}}
\toprule
\textbf{器件} & \textbf{驱动方式} & \textbf{控制程度} & \textbf{开关速度} & \textbf{特点} \\
\midrule
电力二极管 & 无需驱动 & 不可控 & 取决于类型 & 结构简单、可靠；快恢复/肖特基用于高频 \\
晶闸管 & 电流（脉冲） & 半控 & 慢（$t_q$ 数百 $\mu$s） & 容量最大、耐压最高；需强迫换流关断 \\
GTO & 电流（脉冲） & 全控 & 较慢 & 容量大、可关断；$\beta_{off}$ 小、驱动功率大 \\
GTR & 电流（持续） & 全控 & 中（几 $\mu$s） & 通态压降小；有二次击穿、驱动功率大 \\
电力 MOSFET & 电压 & 全控 & 快（ns 级） & 驱动功率小、$R_{on}$ 正温度系数；耐压不易高 \\
IGBT & 电压 & 全控 & 较快 & 兼具 MOSFET 驱动优点与 GTR 通态优点；须防掣住 \\
\bottomrule
\end{tabular}
\end{table}

% ============================================================
\section{其他新型电力电子器件}

\subsection{MOS 控制晶闸管（MCT）}

MCT（MOS Controlled Thyristor）是把 MOSFET 与晶闸管集成的复合器件，可在门极加信号控制其开通与关断，属于电压驱动型全控器件。其特点是：通态压降小、通流能力强、du/dt 与 di/dt 耐量高、开关速度快；缺点是制造工艺复杂、成本高。MCT 一度被看作很有前途的大功率器件，但实用化进展不如 IGBT。

\subsection{静电感应晶体管（SIT）}

SIT（Static Induction Transistor）是一种结型场效应器件，为\textbf{多子导电的单极型器件}，具有输出功率大、输入阻抗高、开关速度快、无二次击穿、抗辐射能力强等优点。缺点是通态电阻较大、驱动电路复杂，多用于高频大功率放大与感应加热电源。

\subsection{静电感应晶闸管（SITH）}

SITH（Static Induction Thyristor）又称场控晶闸管（FCTh），是在 SIT 基础上发展起来的，属\textbf{双极型}器件。其通态压降小、通流能力强、开关速度快、du/dt 与 di/dt 耐量高；缺点是需要较大的驱动功率、关断时存在尾部电流。

\subsection{集成门极换流晶闸管（IGCT）}

IGCT（Integrated Gate Commutated Thyristor）是 GTO 与 MOS 驱动电路集成而成的新型器件，实质是把门极驱动电路与 GTO 管芯集成在一起，实现“门极换流”。其特点是：容量大（可达数千安、数千伏）、通态压降小、开关频率比 GTO 高、无需缓冲电路、可靠性高，在大功率变频器、电力系统（如静止同步补偿器、直流输电）中已获得大量应用。

\subsection{基于宽禁带半导体材料的电力电子器件}

到目前为止，硅材料一直是电力电子器件的主要材料。硅的禁带宽度为 1.12 eV，而\textbf{宽禁带半导体材料}是指禁带宽度在 3.0 eV 及以上的半导体材料，典型的是\textbf{碳化硅（SiC）、氮化镓（GaN）}等。宽禁带材料具有临界击穿场强高、热导率高、饱和电子漂移速度高等优点，用其制成的器件\textbf{耐压更高、通态电阻更小、可在更高温度和更高频率下工作}，是电力电子器件的重要发展方向。目前 SiC 肖特基二极管、SiC MOSFET 已商品化，但成本、可靠性与驱动技术仍有很长的路要走。

% ============================================================
\section{功率集成电路与集成电力电子模块}

自 20 世纪 80 年代中后期开始，在电力电子器件研制和开发中的一个共同趋势是\textbf{集成化}，这就是电力电子集成电路（PIC）的概念。

\begin{itemize}
  \item \textbf{功率模块}（Power Module）：把若干个电力电子器件及必要的辅助元件封装成一个模块，给应用带来很大方便——结构紧凑、体积减小、可靠性提高。
  \item \textbf{集成电力电子模块}（IPEM）/ \textbf{智能功率模块}（IPM）：进一步把驱动、控制、保护（含过流、过热、欠压保护）电路与电力电子器件集成在一起，构成“功率 + 驱动 + 保护”的模块。用户只需给输入信号，模块自己完成保护动作，这在实际工程中极为常用。
  \item \textbf{单片集成}：把主电路器件与控制电路做在同一芯片上，难度最大，受制于\textbf{电压隔离}（主电路高压、控制电路低压）、\textbf{热隔离}（主电路发热严重）、\textbf{电磁干扰}（开关器件通断高压大电流且与控制电路同处一芯）等难题，目前功率都还较小。
  \item \textbf{混合集成}：把不同的单个芯片集成封装在一起，功率密度不如单片集成，但为解决上述难题提供了很大方便，\textbf{封装技术成为关键技术}，是当前电力电子集成技术发展的焦点。此外\textbf{系统集成}也是重要方面，对超大功率集成技术尤其如此。
\end{itemize}

\corepoint{本章复习主线：\textbf{“控制程度—驱动性质—载流子”三把尺子}给器件定位；\textbf{晶闸管}抓“门极只开不关 + 双晶体管正反馈 + 选型三步”；\textbf{全控型器件}抓“驱动功率与通态压降的折中”——GTR 通态好但驱动费劲、MOSFET 驱动省但通态差、IGBT 取二者之长，这就是第 5—8 章各种电路选用器件的依据。}

\chapter{整流电路}

整流电路（Rectifier）是电力电子电路中出现最早、应用最广泛的一种。本章讲的是\textbf{相控整流}：以晶闸管为开关，通过控制触发脉冲的相位（触发延迟角 $\alpha$）来调节输出电压平均值。分析任何整流电路都按同一套步骤走：\textbf{（1） 画电路，认清哪两个管子构成一条导电回路；（2） 定自然换相点（$\alpha$ 的计量起点）；（3） 分段看器件通断、写回路方程；（4） 积分求 $U_d$、求电流有效值；（5） 核对器件的电压电流定额}。

% ============================================================
\section{单相可控整流电路}

典型的单相可控整流电路包括单相半波、单相桥式全控、单相全波（双半波）、单相桥式半控等。本节先以最简单的半波电路建立“触发延迟角 $\alpha$、导通角 $\theta$、移相范围”等基本概念，再讲工程上最常用的桥式电路。

\subsection{单相半波可控整流电路}

\begin{center}
\begin{circuitikz}[font=\footnotesize,scale=0.9]
  \draw (0,0) to[sV, l=$u_2$] (0,-2.2);
  \draw (0,0) -- (1.6,0) to[thyristor] (3.4,0) -- (4.6,0) to[R, l=$R$] (4.6,-2.2) -- (0,-2.2);
  \node[above] at (1.0,0.15) {\scriptsize $+$};
  \node[above] at (2.55,0.35) {$VT$};
\end{circuitikz}
\end{center}

\subsubsection*{1. 带电阻负载}

变压器二次电压 $u_2=\sqrt{2}U_2\sin\omega t$。在 $u_2$ 正半周，若门极未加触发脉冲，晶闸管处于正向阻断状态，负载电压为零；在 $\omega t=\alpha$ 时刻给门极加触发脉冲，晶闸管导通，$u_d=u_2$；到 $\omega t=\pi$（$u_2$ 过零）时，流过晶闸管的电流降至零而自然关断。负半周晶闸管承受反压而阻断，$u_d=0$。

\begin{itemize}
  \item \textbf{触发延迟角} $\alpha$：从晶闸管开始承受正向电压起到施加触发脉冲止的电角度，也称触发角或控制角。
  \item \textbf{导通角} $\theta$：晶闸管在一个周期内导通的电角度。电阻负载时 $\theta=\pi-\alpha$。
  \item \textbf{移相范围}：改变 $\alpha$ 使输出电压从一个极端变到另一个极端时 $\alpha$ 的可调范围。电阻负载半波电路为 $0^\circ\sim180^\circ$。
\end{itemize}

\keyformula{U_d=\frac{1}{2\pi}\int_{\alpha}^{\pi}\sqrt{2}U_2\sin\omega t\,\mathrm{d}(\omega t)=\frac{\sqrt{2}U_2}{2\pi}\left(1+\cos\alpha\right)=0.45U_2\cdot\frac{1+\cos\alpha}{2}}

$\alpha=0$ 时 $U_d$ 最大，记为 $U_{d0}=0.45U_2$；$\alpha=180^\circ$ 时 $U_d=0$。负载电流平均值 $I_d=U_d/R$。晶闸管承受的\textbf{最大正向电压为 $\sqrt{2}U_2$、最大反向电压也为 $\sqrt{2}U_2$}；变压器二次电流即负载电流，含有直流分量。

\subsubsection*{2. 带阻感负载}

电感 $L$ 的存在使 $u_2$ 由正变负时晶闸管不会立即关断——$L$ 的感应电动势试图维持电流，只要 $u_L=u_2-u_d$ 为负使电流下降，晶闸管就继续导通，直到电感储能释放完毕、电流降到零才关断。因此 $u_d$ 波形出现\textbf{负值}，平均值下降。

回路由 $L$、$R$ 串联，晶闸管导通期间有
\keyformula{L\frac{\mathrm{d}i_d}{\mathrm{d}t}+Ri_d=\sqrt{2}U_2\sin\omega t}

解得 $i_d=\dfrac{\sqrt{2}U_2}{Z}\left[\sin(\omega t-\varphi)-\sin(\alpha-\varphi)\mathrm{e}^{-\frac{\omega t-\alpha}{\omega\tau}}\right]$，其中 $Z=\sqrt{R^2+(\omega L)^2}$、$\varphi=\arctan(\omega L/R)$、$\tau=L/R$。由 $i_d(\alpha+\theta)=0$ 得导通角与 $\alpha$、$\varphi$ 的关系式（超越方程，只能迭代或查曲线）：

\keyformula{\sin(\alpha+\theta-\varphi)=\sin(\alpha-\varphi)\mathrm{e}^{-\frac{\theta}{\omega\tau}}}

$\varphi$ 越大（负载电感越重），导通角 $\theta$ 越大，$u_d$ 负面积越大、$U_d$ 越小；$\alpha$ 增大时 $U_d$ 减小，当 $\alpha$ 大到一定程度 $U_d$ 甚至会降到零（正负面积相等）。

\subsubsection*{3. 带续流二极管的阻感负载}

在负载两端反并联一个二极管 $VD_R$（续流二极管）。$u_2$ 过零变负时 $VD_R$ 导通，$u_d=0$；此时为负的 $u_2$ 通过 $VD_R$ 向晶闸管施加反压使其关断，$L$ 储存的能量保证电流在 $L$—$R$—$VD_R$ 回路中继续流通，这一过程称为\textbf{续流}。

\begin{itemize}
  \item 有续流二极管后 $u_d$ 波形与电阻负载时基本相同（不再出现负值），故\textbf{计算公式仍为} $U_d=0.45U_2(1+\cos\alpha)/2$。
  \item 晶闸管不承受反压；晶闸管电流平均值 $I_{dT}=\dfrac{\pi-\alpha}{2\pi}I_d$，续流二极管电流平均值 $I_{dD}=\dfrac{\pi+\alpha}{2\pi}I_d$。
  \item 电感越大，$i_d$ 越平直，上述按平均电流分配的式子越准确。
\end{itemize}

\corepoint{单相半波电路元件少、线路简单，但\textbf{变压器二次电流含直流分量}，会造成铁心直流磁化、铁心饱和，需增大铁心截面积、增大设备容量，因此\textbf{实际很少应用}。分析它的主要目的是利用其简单易学的特点，建立整流电路的基本概念。}

\subsection{单相桥式全控整流电路}

把四只晶闸管接成桥式（$VT_1$、$VT_4$ 为一对，$VT_2$、$VT_3$ 为一对），$VT_1$ 与 $VT_4$ 同时触发导通、$VT_2$ 与 $VT_3$ 同时触发导通，两组轮流导电，变压器二次电流正负半周都有，因而不含直流分量。

\begin{center}
\begin{circuitikz}[font=\footnotesize,scale=0.85]
  % 左桥臂：VT1 上、VT3 下；右桥臂：VT2 上、VT4 下
  \draw (0,2.9) to[thyristor] (0,-1) to[thyristor] (0,-4.9);
  \draw (4,2.9) to[thyristor] (4,-1) to[thyristor] (4,-4.9);
  % 交流侧：接两桥臂中点
  \draw (0,-1) -- (1.3,-1) to[sV, l=$u_2$] (2.7,-1) -- (4,-1);
  % 直流侧：上、下汇流排
  \draw (0,2.9) -- (5.4,2.9);
  \draw (0,-4.9) -- (5.4,-4.9);
  % 负载
  \draw (5.4,2.9) to[R, l=$R$] (5.4,0.2) to[L, l=$L$] (5.4,-2.4) -- (5.4,-4.9);
  % 标注
  \node[left] at (-0.45,2.0) {\scriptsize $VT_1$};
  \node[left] at (-0.45,-3.7) {\scriptsize $VT_3$};
  \node[right] at (4.45,2.0) {\scriptsize $VT_2$};
  \node[right] at (4.45,-3.7) {\scriptsize $VT_4$};
  \node[right] at (0.1,-1.45) {\scriptsize $a$};
  \node[left] at (3.9,-1.45) {\scriptsize $b$};
  \node[right] at (6.15,0.6) {\scriptsize $u_d$};
\end{circuitikz}
\end{center}

\subsubsection*{1. 带电阻负载}

$\alpha$ 的起点是 $u_2$ 由负变正的自然换相点。在 $u_2$ 正半周 $\alpha$ 处触发 $VT_1$、$VT_4$，电流路径 $a\to VT_1\to R\to VT_4\to b$；在负半周 $\alpha+\pi$ 处触发 $VT_2$、$VT_3$，电流路径 $b\to VT_2\to R\to VT_3\to a$。两组交替，$u_d$ 为两段相同的脉动波。

\keyformula{U_d=\frac{1}{\pi}\int_{\alpha}^{\pi}\sqrt{2}U_2\sin\omega t\,\mathrm{d}(\omega t)=\frac{2\sqrt{2}U_2}{\pi}\cdot\frac{1+\cos\alpha}{2}=0.9U_2\cdot\frac{1+\cos\alpha}{2}}

$\alpha=0$ 时 $U_{d0}=0.9U_2$；$\alpha=180^\circ$ 时 $U_d=0$，故\textbf{移相范围为 $0^\circ\sim180^\circ$}。$\alpha$ 角相同时，桥式电路的 $U_d$ 是半波电路的\textbf{两倍}。

流过晶闸管的电流平均值 $I_{dT}=I_d/2$，有效值 $I_T=I_d/\sqrt{2}$；变压器二次电流有效值 $I_2=I_d$。晶闸管承受的\textbf{最大反向电压为 $\sqrt{2}U_2$}。

\subsubsection*{2. 带阻感负载（电流连续）}

电感足够大使 $i_d$ 平直（近似水平线）。$u_2$ 过零变负时，因 $L$ 的作用 $VT_1$、$VT_4$ 不关断，直到 $\alpha+\pi$ 处触发 $VT_2$、$VT_3$ 才被迫关断（换相）。因此 $u_d$ 出现\textbf{负值}，正负面积相抵。

\keyformula{U_d=\frac{1}{\pi}\int_{\alpha}^{\pi+\alpha}\sqrt{2}U_2\sin\omega t\,\mathrm{d}(\omega t)=\frac{2\sqrt{2}U_2}{\pi}\cos\alpha=0.9U_2\cos\alpha}

$\alpha=0$ 时 $U_d=0.9U_2$；$\alpha=90^\circ$ 时 $U_d=0$，故\textbf{阻感负载时移相范围为 $0^\circ\sim90^\circ$}。$\alpha>90^\circ$ 时 $U_d<0$，电路进入\textbf{有源逆变状态}（见 3.7 节）。

电流关系：$I_{dT}=I_d/2$，$I_T=I_d/\sqrt{2}=0.707I_d$；变压器二次电流为正负各 $180^\circ$ 的矩形波，有效值 $I_2=I_d$；晶闸管承受的最大正反向电压均为 $\sqrt{2}U_2$。

\subsubsection*{3. 带反电动势—电阻负载（电流断续）}

直流侧为电阻 $R$ 与反电动势 $E$ 串联，且 $E$ 的极性对晶闸管为正（阻碍电流）。只有 $u_2>E$ 时晶闸管才可能导通，故存在\textbf{停止导电角} $\delta=\arcsin\dfrac{E}{\sqrt{2}U_2}$：在 $\omega t=\pi-\delta$ 处 $u_2$ 降到 $E$，电流降为零而关断。

\keyformula{U_d=\frac{\sqrt{2}U_2}{\pi}\left(\cos\alpha+\cos\delta\right)=0.45U_2\left(\cos\alpha+\cos\delta\right)}

\corepoint{整流输出电压 $U_d$ 波形在一个周期内有部分时间为零的情况称为\textbf{电流断续}，不出现为零的情况称为\textbf{电流连续}。$\alpha$ 角相同时，电流断续时的整流输出电压\textbf{比电阻负载时大}（因导通段集中在电压较高处）；$\alpha$ 小于 $\delta$ 时，触发脉冲到来时 $u_2<E$、晶闸管不能导通，须把脉冲留到 $\omega t=\delta$ 之后（“宽脉冲”或“脉冲列”）。这也是判断“电流是否断续”的常用判据：\textbf{先按连续公式算 $I_d$，若 $I_d$ 小到不足以维持连续，则改用断续公式}。}

\subsection{单相全波可控整流电路}

单相全波（又称双半波）可控整流电路用带中心抽头的变压器与两只晶闸管：$u_2$ 正半周 $VT_1$ 导通，负半周 $VT_2$ 导通，两只管子共用一个负载回路。

\begin{itemize}
  \item 输出电压、电流波形与单相桥式全控电路在\textbf{电阻负载时完全相同}，故 $U_d=0.9U_2(1+\cos\alpha)/2$，阻感负载连续时 $U_d=0.9U_2\cos\alpha$。
  \item 与桥式相比：\textbf{少一个管压降}（每支路只有一只晶闸管），通态损耗小；但变压器二次绕组需中心抽头，\textbf{变压器利用率低}，且每只晶闸管承受的最大反向电压为 $2\sqrt{2}U_2$（是桥式的两倍）。
\end{itemize}

\subsection{单相桥式半控整流电路}

在单相桥式全控电路中，每一个导电回路里有两只晶闸管，实际上对每个导电回路进行控制只需一只晶闸管，另一只可以用二极管代替，从而简化电路——把 $VT_2$、$VT_4$ 换成二极管 $VD_2$、$VD_4$，即为\textbf{单相桥式半控整流电路}。

\begin{itemize}
  \item \textbf{电阻负载}：工作情况与全控桥完全相同，$U_d=0.9U_2(1+\cos\alpha)/2$，移相范围 $0^\circ\sim180^\circ$。
  \item \textbf{阻感负载（无续流二极管）}：$u_2$ 过零变负时，因电感作用负载电流经 $VT_1$ 和 $VD_2$ 继续流通（$VD_2$ 起续流作用），$u_d=0$ 而不出现负值——$U_d$ 与电阻负载时相同。但此时若 $\alpha$ 突然增大到 $180^\circ$ 或触发脉冲丢失，正在导通的晶闸管会一直导通到下一个周期，电路失去控制，称为\textbf{失控}。
  \item \textbf{加续流二极管 $VD_R$}：把续流任务交给 $VD_R$，$u_d$ 波形与电阻负载时一样，且续流期间导电回路中只有一只管子，\textbf{少了一个管压降}，有利于降低损耗、避免失控。
\end{itemize}

\corepoint{半控桥的优点是\textbf{器件少、触发器简单、成本低}，缺点是\textbf{不能用于有源逆变}（$u_d$ 不能为负）、电压调节范围受限。凡“直流侧要出现负电压”的场合，一律必须用全控电路。}

% ============================================================
\section{三相可控整流电路}

当负载容量较大，或要求直流电压脉动较小、容易滤波时，应采用三相可控整流电路。最基本的\textbf{三相半波可控整流电路}，其整流电压在一个周期内脉动三次（脉波数 $m=3$）——三相桥式全控、双反星形、十二脉波等电路，均可在三相半波的基础上进行分析。

\subsection{三相半波可控整流电路}

三相半波（三相零式）电路把三只晶闸管的阴极连在一起（共阴极）接负载，阳极分别接三相电源。$\alpha$ 的计量起点（自然换相点）是\textbf{相邻两相电压的交点}，即相电压正半周交点，距相电压过零点 $30^\circ$（$\pi/6$）。

\begin{center}
\begin{circuitikz}[font=\footnotesize,scale=0.85]
  % 三相电源（下端为中性点 n）
  \draw (0,0) to[sV, l=$u_a$] (0,-2.0);
  \draw (2,0) to[sV, l=$u_b$] (2,-2.0);
  \draw (4,0) to[sV, l=$u_c$] (4,-2.0);
  \draw (0,-2.0) -- (5.9,-2.0);
  \node[below] at (2,-2.05) {\scriptsize $n$};
  % 三只晶闸管共阴极
  \draw (0,0) to[thyristor] (0,1.6) -- (5.9,1.6);
  \draw (2,0) to[thyristor] (2,1.6);
  \draw (4,0) to[thyristor] (4,1.6);
  % 负载接在共阴极与中性点之间
  \draw (5.6,1.6) to[R, l=$R$] (5.6,-2.0);
  \node[right] at (6.7,-0.3) {\scriptsize $u_d$};
  \node[left] at (-0.45,0.8) {\scriptsize $VT_1$};
  \node[left] at (1.55,0.8) {\scriptsize $VT_2$};
  \node[left] at (3.55,0.8) {\scriptsize $VT_3$};
  \node[left] at (-0.2,2.15) {\scriptsize $+$};
\end{circuitikz}
\end{center}

\subsubsection*{1. 电阻负载}

\begin{itemize}
  \item \textbf{$\alpha\le30^\circ$：电流连续}。负载电流处于连续状态，$u_d$ 为三相相电压正半周的包络线，每只晶闸管导通 $120^\circ$。
  \item \textbf{$\alpha>30^\circ$：电流断续}。$u_d$ 波形不再出现负值（电阻负载 $u_d$ 不能为负），导通角减小。
\end{itemize}

\keyformula{U_d=\begin{cases}\dfrac{3\sqrt{6}U_2}{2\pi}\cos\alpha=1.17U_2\cos\alpha, & \alpha\le30^\circ\\[6pt] 0.675U_2\left[1+\cos\left(\dfrac{\pi}{6}+\alpha\right)\right], & \alpha>30^\circ\end{cases}}

$\alpha=0$ 时 $U_d$ 最大，$U_{d0}=1.17U_2$；$\alpha=150^\circ$ 时 $U_d=0$，故\textbf{电阻负载移相范围为 $0^\circ\sim150^\circ$}。负载电流 $I_d=U_d/R$。

晶闸管承受的\textbf{最大反向电压为 $\sqrt{2}\cdot\sqrt{3}U_2=\sqrt{6}U_2=2.45U_2$}（线电压峰值），最大正向电压为 $\sqrt{2}U_2$（相电压峰值）。流过晶闸管的电流平均值 $I_{dT}=I_d/3$。

\subsubsection*{2. 阻感负载}

电感足够大时 $i_d$ 平直，$u_d$ 波形出现负值部分（每周期三段，每段 $-30^\circ+\alpha$ 的负面积）。

\keyformula{U_d=1.17U_2\cos\alpha\qquad(0^\circ\le\alpha\le90^\circ)}

$\alpha=90^\circ$ 时 $U_d=0$；$\alpha>90^\circ$ 时 $U_d<0$，进入有源逆变状态。变压器二次电流为\textbf{正负各 $120^\circ$} 的矩形波，有效值 $I_2=\sqrt{\dfrac{120^\circ}{360^\circ}}I_d=\dfrac{I_d}{\sqrt{3}}=0.577I_d$。

\corepoint{三相半波电路的主要缺点是\textbf{变压器二次电流中含有直流分量}，因此其应用较少。为克服此缺点可采用\textbf{双反星形}电路（3.6 节）。}

\subsection{三相桥式全控整流电路}

三相桥式全控整流电路共 6 只晶闸管，\textbf{共阴极组}（$VT_1$、$VT_3$、$VT_5$）与\textbf{共阳极组}（$VT_4$、$VT_6$、$VT_2$）接成桥。共阴极组的二极管在相电压最高时导通，共阳极组的在相电压最低时导通，因此 $u_d$ 是三相线电压包络线，一周期脉动 6 次（$m=6$）。

\begin{table}[htbp]
\centering
\caption{三相桥式全控整流电路的触发顺序（$\alpha=0$ 时）}
\label{tab:3-1}
\begin{tabular}{cccccc}
\toprule
\textbf{时段} & \textbf{导通管} & \textbf{输出电压 $u_d$} & \textbf{触发时刻（电角度）} & \textbf{换相} & \textbf{备注} \\
\midrule
1 & $VT_6$、$VT_1$ & $u_{ab}$ & $0^\circ\sim60^\circ$ & $VT_5\to VT_1$ & 共阴组换相 \\
2 & $VT_1$、$VT_2$ & $u_{ac}$ & $60^\circ\sim120^\circ$ & $VT_6\to VT_2$ & 共阳组换相 \\
3 & $VT_2$、$VT_3$ & $u_{bc}$ & $120^\circ\sim180^\circ$ & $VT_1\to VT_3$ & 共阴组换相 \\
4 & $VT_3$、$VT_4$ & $u_{ba}$ & $180^\circ\sim240^\circ$ & $VT_2\to VT_4$ & 共阳组换相 \\
5 & $VT_4$、$VT_5$ & $u_{ca}$ & $240^\circ\sim300^\circ$ & $VT_3\to VT_5$ & 共阴组换相 \\
6 & $VT_5$、$VT_6$ & $u_{cb}$ & $300^\circ\sim360^\circ$ & $VT_4\to VT_6$ & 共阳组换相 \\
\bottomrule
\end{tabular}
\end{table}

触发脉冲必须按 $VT_1\to VT_2\to VT_3\to VT_4\to VT_5\to VT_6$ 的顺序、依次相差 $60^\circ$ 送入。为保证合闸后或电流断续后能形成导通回路，须用\textbf{双窄脉冲}（每个触发单元在一个周期内发两个相隔 $60^\circ$ 的窄脉冲）或\textbf{宽脉冲}（脉宽大于 $60^\circ$）。

\subsubsection*{1. 阻感负载（电流连续）}

$i_d$ 平直，$u_d$ 波形出现负值部分，$u_d$ 在一个周期内由 6 段相同形状组成：

\keyformula{U_d=\frac{1}{\pi/3}\int_{\frac{\pi}{3}+\alpha}^{\frac{2\pi}{3}+\alpha}\sqrt{6}U_2\sin\omega t\,\mathrm{d}(\omega t)=\frac{3\sqrt{6}U_2}{\pi}\cos\alpha=2.34U_2\cos\alpha}

$\alpha=0$ 时 $U_{d0}=2.34U_2$；$\alpha=90^\circ$ 时 $U_d=0$，故\textbf{阻感负载移相范围为 $0^\circ\sim90^\circ$}。

数量关系（电流连续、忽略换相）：
\begin{itemize}
  \item 流过每只晶闸管的电流平均值 $I_{dT}=I_d/3$，有效值 $I_T=I_d/\sqrt{3}=0.577I_d$；
  \item 变压器二次电流为正负各 $120^\circ$、前沿相差 $180^\circ$ 的矩形波，有效值 $I_2=\sqrt{2}I_{dT}$，按星形接法 $I_2=\sqrt{\dfrac{2}{3}}I_d=0.816I_d$；
  \item 晶闸管承受的最大正反向电压均为线电压峰值 $\sqrt{6}U_2=2.45U_2$。
\end{itemize}

\subsubsection*{2. 电阻负载}

\begin{itemize}
  \item $\alpha\le60^\circ$：$u_d$ 连续，公式同上 $U_d=2.34U_2\cos\alpha$。
  \item $\alpha>60^\circ$：$u_d$ 出现为零的段（电流断续），
\end{itemize}

\keyformula{U_d=2.34U_2\left[1+\cos\left(\frac{\pi}{3}+\alpha\right)\right],\qquad \alpha>60^\circ}

$\alpha=120^\circ$ 时 $U_d=0$，故\textbf{电阻负载移相范围为 $0^\circ\sim120^\circ$}。

\corepoint{三相桥式全控电路输出脉波数 $m=6$，与三相半波相比：输出电压提高一倍（$2.34U_2$ 对 $1.17U_2$）、脉动大大减小（纹波因数 $4.18\%$ 对 $18.27\%$）、变压器二次电流无直流分量、变压器利用率高。它是\textbf{大功率整流装置应用最广泛的一种形式}。}

% ============================================================
\section{变压器漏感对整流电路的影响}

前面分析都假设整流变压器是理想的，即绕组无漏感、换相瞬间完成。实际上变压器绕组总有漏感，折算到二次侧的每相漏感为 $L_B$（对应漏抗 $X_B=\omega L_B$）。由于电感电流不能突变，换相\textbf{不能瞬间完成}，而需要一段时间——从一相换到另一相的过程称为\textbf{换相过程}，所对应的电角度称为\textbf{换相重叠角} $\gamma$。

以三相半波带大电感负载、从 $VT_1$ 换相到 $VT_2$ 为例。在触发 $VT_2$ 之后的一段时间内，$VT_1$、$VT_2$ 同时导通，两相通过漏感短路，回路方程为

\keyformula{u_a-u_b=2L_B\frac{\mathrm{d}i_k}{\mathrm{d}t}}

换相期间\textbf{整流电压为同时导通的两个晶闸管所对应的两个相电压的平均值}，即 $u_d=(u_a+u_b)/2$。与不考虑漏感时相比，每次换相 $u_d$ 波形都少了一块（阴影面积），导致平均值降低，降低量称为\textbf{换相压降} $\Delta U_d$。

\begin{table}[htbp]
\centering
\caption{各种整流电路的换相压降与换相重叠角}
\label{tab:3-2}
\begin{tabular}{lll}
\toprule
\textbf{电路形式} & \textbf{换相压降 $\Delta U_d$} & \textbf{换相重叠角关系} \\
\midrule
单相全波、单相全控桥 & $\Delta U_d=\dfrac{2X_BI_d}{\pi}$ & $\cos\alpha-\cos(\alpha+\gamma)=\dfrac{\sqrt{2}X_BI_d}{U_2}$ \\
三相半波 & $\Delta U_d=\dfrac{3X_BI_d}{2\pi}$ & $\cos\alpha-\cos(\alpha+\gamma)=\dfrac{2X_BI_d}{\sqrt{6}U_2}$ \\
三相全控桥 & $\Delta U_d=\dfrac{3X_BI_d}{\pi}$ & 同上式 \\
\bottomrule
\end{tabular}
\end{table}

由 $U_d=U_{d0}\cos\alpha-\Delta U_d$ 可算出计及漏感后的输出电压。换相重叠角 $\gamma$ 的变化规律（由上式分析得出）：

\begin{enumerate}
  \item $I_d$ 越大，$\gamma$ 越大；
  \item $X_B$ 越大（变压器漏抗越大），$\gamma$ 越大；
  \item 当 $\alpha<90^\circ$ 时，$\alpha$ 越小，$\gamma$ 越大。
\end{enumerate}

\corepoint{漏感的存在带来两方面后果：\textbf{坏处}是输出电压平均值降低（$\Delta U_d$）、实际移相范围扩大、功率因数略有下降；\textbf{好处}是使电流变化平缓，相当于在每个桥臂串了一个电感，从而\textbf{限制了晶闸管的 di/dt 和短路电流上升率}，对器件有保护作用——所以实际装置中有时反而人为串入进线电抗器（“换相电抗器”）。}

% ============================================================
\section{电容滤波的不可控整流电路}

3.1、3.3 节介绍的可控整流电路重点针对阻感负载；而在交—直—交变频器、不间断电源、开关电源等场合，大都采用\textbf{不可控整流电路经电容滤波}后提供直流电源，供后级的逆变器、斩波器使用。把前述各种全控整流电路中的晶闸管换成整流二极管，就是不可控整流电路（二极管整流电路）。

\subsection{电容滤波的单相不可控整流电路}

本电路常用于小功率单相交流输入的场合（微机、电视机等家电产品开关电源的整流部分就是单相桥式电容滤波电路）。

工作原理：$u_2$ 正半周且 $u_2>u_d$ 时，$VD_1$、$VD_4$ 导通，电源向电容 $C$ 充电并向负载供电，$u_d=u_2$；到 $u_2$ 峰值后 $u_2$ 下降，当 $u_2<u_d$ 时二极管关断，$C$ 向 $R$ 放电，$u_d$ 按 $u_d=u_{d0}\mathrm{e}^{-t/(RC)}$ 指数下降；到负半周 $|u_2|>u_d$ 时 $VD_2$、$VD_3$ 导通，又向 $C$ 充电。

\begin{itemize}
  \item 由于电容储能，\textbf{二极管的导通角小于 $180^\circ$}，且 $RC$ 越大导通角越小、$u_d$ 越平直。二极管中流过的电流是\textbf{窄而高的脉冲}，有效值大、谐波含量高。
  \item 输出电压平均值与 $RC$ 和负载有关：\textbf{空载时 $U_d=\sqrt{2}U_2$}；负载加重（$R$ 减小）时 $U_d$ 下降，重载时趋近电阻负载的 $0.9U_2$。
  \item 工程估算常取 $U_d\approx1.2U_2$（$RC$ 足够大时）。$RC$ 的确定还要兼顾：$RC$ 过大则充电冲击电流大、二极管电流峰值高。
\end{itemize}

\keyformula{U_{d0}=\sqrt{2}U_2\ \ (\text{空载}),\qquad U_d\approx1.2U_2\ \ (\text{工程估算}),\qquad U_d\to0.9U_2\ \ (\text{重载})}

\subsection{电容滤波的三相不可控整流电路}

最常用的是三相桥式结构。某一对二极管导通时，输出直流电压等于交流侧线电压中最大的一个；该线电压既向电容充电，也向负载供电；当线电压最大值过去、电容电压高于线电压时，所有二极管关断，电容向负载放电。

\begin{itemize}
  \item \textbf{轻载}：直流侧充电电流\textbf{断续}（电容放电时间占主导）；
  \item \textbf{重载}：$VD$ 一直有导通，换相在相邻二极管间交替进行，电流\textbf{连续}；
  \item 二者的\textbf{临界情况}：$VD_6$ 和 $VD_1$ 同时导通的阶段与 $VD_1$ 和 $VD_2$ 同时导通的阶段恰好衔接，$i_d$ 恰好连续。临界条件由 $RC$ 与电源频率决定（分界点 $R=\dfrac{1}{\omega C}$ 量级）。
  \item 输出电压：\textbf{空载时 $U_d=\sqrt{2}U_{LL}=\sqrt{6}U_2$}，满载时趋近 $1.35U_{LL}=2.34U_2$。
\end{itemize}

\keyformula{U_d\approx1.35U_{LL}=2.34U_2\ \ (\text{工程估算，重载}),\qquad U_{d0}=\sqrt{2}U_{LL}\ \ (\text{空载})}

\corepoint{电容滤波电路的共性问题：\textbf{（1） 二极管导通角小、电流脉冲窄，电流有效值大、谐波大、功率因数低（约 0.6）；（2） 合闸瞬间有很大的充电冲击电流}，实际电路中须串入滤波电感或限流电阻抑制；（3） 输出电压随负载变化大，不宜用于要求稳压的场合（需再加 DC-DC 变换器）。}

% ============================================================
\section{整流电路的谐波和功率因数}

各种电力电子装置在电力系统、工业、交通、家庭等领域应用日益广泛，由此带来的\textbf{谐波}（Harmonics）和\textbf{无功}（Reactive Power）问题也日益严重。我国已制定《电能质量 公用电网谐波》标准并实施。

\subsection{谐波和无功功率分析基础}

正弦波电压可表示为 $u(t)=\sqrt{2}U\sin(\omega t+\varphi_u)$。当电压为正弦、电流为非正弦（含谐波）时，将电流分解为傅里叶级数

\keyformula{i(t)=I_{dc}+\sum_{n=1}^{\infty}\sqrt{2}I_n\sin(n\omega t+\varphi_n)}

则有功功率只由同频率的电压电流分量产生：$P=UI_1\cos\varphi_1$；视在功率 $S=UI$（$I$ 为电流总有效值）；功率因数

\keyformula{\lambda=\frac{P}{S}=\frac{I_1}{I}\cos\varphi_1=\nu\cos\varphi_1}

式中 $\nu=I_1/I$ 称为\textbf{基波因数}，$\cos\varphi_1$ 称为\textbf{位移因数}（基波功率因数）。可见\textbf{功率因数由基波电流相移和电流波形畸变这两个因素共同决定}。非正弦情况下的无功功率定义很多，一种被广泛接受的定义是 $Q=S\sin\varphi_1$；也可仿照正弦情况定义 $Q_f=UI_1\sin\varphi_1$，二者之差即为\textbf{畸变功率} $D=\sqrt{S^2-P^2-Q_f^2}$。

电流谐波总畸变率定义为 $\mathrm{THD}_i=\dfrac{\sqrt{\sum_{n=2}^{\infty}I_n^2}}{I_1}\times100\%$。

\subsection{带阻感负载时可控整流电路交流侧谐波和功率因数分析}

\subsubsection*{1. 单相桥式全控整流电路}

忽略换相过程和电流脉动，阻感负载、电流连续时，变压器二次电流为\textbf{正负各 $180^\circ$ 的矩形波}，将其分解为傅里叶级数：

\keyformula{i_2=\frac{4I_d}{\pi}\left(\sin\omega t+\frac{1}{3}\sin3\omega t+\frac{1}{5}\sin5\omega t+\cdots\right)=\frac{4I_d}{\pi}\sum_{n=1,3,5,\cdots}\frac{1}{n}\sin n\omega t}

\begin{itemize}
  \item 电流中仅含\textbf{奇次谐波}（$2k\pm1$，$k$ 为正整数），各次谐波有效值与谐波次数成反比，且与基波有效值的比值为谐波次数的倒数；
  \item 基波有效值 $I_1=\dfrac{2\sqrt{2}I_d}{\pi}=0.9I_d$，而 $I=I_d$，故\textbf{基波因数 $\nu=0.9$}；
  \item 由图 3-6 可明显看出，电流基波与电压的相位差就等于控制角 $\alpha$，故\textbf{位移因数 $\cos\varphi_1=\cos\alpha$}。
\end{itemize}

\keyformula{\lambda=\nu\cos\varphi_1=0.9\cos\alpha}

\subsubsection*{2. 三相桥式全控整流电路}

阻感负载、电流连续时，交流侧电流为\textbf{正负半周各 $120^\circ$ 的方波}，三相电流波形相同且依次相差 $120^\circ$，其有效值与直流电流的关系为 $I=\sqrt{2/3}I_d=0.816I_d$。分解为傅里叶级数（以相电流为例，将负、正两半波的中点作为时间零点）：

\keyformula{i_a=\frac{2\sqrt{3}I_d}{\pi}\left(\sin\omega t-\frac{1}{5}\sin5\omega t-\frac{1}{7}\sin7\omega t+\frac{1}{11}\sin11\omega t+\cdots\right)}

\begin{itemize}
  \item 电流中仅含 $6k\pm1$（$k$ 为正整数）次谐波，各次谐波有效值与谐波次数成反比，与基波有效值的比值为谐波次数的倒数；
  \item 基波有效值 $I_1=\dfrac{\sqrt{6}I_d}{\pi}=0.78I_d$，基波因数 $\nu=\dfrac{I_1}{I}=0.955$；
  \item 位移因数仍为 $\cos\varphi_1=\cos\alpha$。
\end{itemize}

\keyformula{\lambda=0.955\cos\alpha}

\subsection{电容滤波的不可控整流电路交流侧谐波和功率因数分析}

实用的单相不可控整流电路带电容滤波时，通常串联滤波电感抑制冲击电流，或因电网侧有电感而具有换相过程。此时交流侧电流为\textbf{窄脉冲}，谐波含量大。其特点是：

\begin{enumerate}
  \item 功率因数低（一般只有 0.6 左右），且\textbf{与输出电压、负载大小有关}；
  \item 随负载加重（$R$ 减小），总的功率因数提高；同时随滤波电感加大，总功率因数也提高。
\end{enumerate}

\subsection{整流输出电压和电流的谐波分析}

整流电路的输出电压是周期性非正弦函数，主要成分为直流，同时包含各种频率的谐波。设 $\alpha=0$ 时 $m$ 脉波整流电路，将纵坐标选在整流电压峰值处，在一个周期内整流电压为 $u_d=\sqrt{2}U_2\cos\omega t$，傅里叶分解后得

\keyformula{u_d=U_{d0}\left[1-\sum_{n=mk}\frac{2\cos\frac{n\pi}{m}}{n^2-1}\cos n\omega t\right],\qquad n=mk\ (k=1,2,3,\cdots)}

\begin{itemize}
  \item $m$ 脉波整流电压的谐波次数为 $mk$ 次，即 $m$ 的倍数次；整流电流的谐波由整流电压的谐波决定，也为 $mk$ 次。
  \item 当 $m$ 一定时，随谐波次数增大，谐波幅值迅速减小，说明\textbf{低次谐波是最主要的}；$m$ 增大时最低次谐波次数增大、且幅值减小，故\textbf{增加脉波数是减小输出电压脉动的主要手段}。
  \item 电压纹波因数 $\gamma_u=U_R/U_d$（$U_R$ 为谐波分量有效值），不同脉波数时的值见表~\ref{tab:3-3}。
\end{itemize}

\begin{table}[htbp]
\centering
\caption{不同脉波数时的电压纹波因数}
\label{tab:3-3}
\begin{tabular}{lcccc}
\toprule
\textbf{脉波数 $m$} & 2 & 3 & 6 & 12 \\
\midrule
\textbf{纹波因数 $\gamma_u$} & $48.2\%$ & $18.27\%$ & $4.18\%$ & $0.994\%$ \\
\bottomrule
\end{tabular}
\end{table}

\corepoint{记牢两组“谐波次数”：\textbf{交流侧电流}——单相桥为 $2k\pm1$（奇次），三相桥为 $6k\pm1$；\textbf{直流侧电压}——$m$ 脉波为 $mk$ 次。功率因数两个公式 $\lambda=0.9\cos\alpha$（单相桥）与 $\lambda=0.955\cos\alpha$（三相桥）必须能默写，并能说明“$\nu$ 反映波形畸变、$\cos\varphi_1$ 反映相位滞后”。}

% ============================================================
\section{大功率可控整流电路}

\subsection{双反星形可控整流电路}

双反星形（带平衡电抗器）电路由\textbf{两套三相半波电路}并联组成：变压器二次侧有两组星形绕组，同名端相反（一组接成星形、另一组也接星形但相序相反），两组的三相半波整流电路通过\textbf{平衡电抗器 $L_p$} 并联，共同向负载供电。

\begin{itemize}
  \item \textbf{平衡电抗器的作用}：使两组整流电路能同时导电。两组输出电压的瞬时值不相等，其差值由 $L_p$ 承担，$L_p$ 的中点即为负载的一端。由于 $L_p$ 的均压作用，两组电流能平均分配。
  \item \textbf{输出电压}：与三相半波相同，$U_d=1.17U_2\cos\alpha$（$U_2$ 为二次相电压有效值）。
  \item \textbf{输出电流}：负载电流为两组电流之和 $I_d=2I_{d1}$，每只晶闸管（每组）承担 $I_d/2$，故\textbf{器件电流定额可比三相桥小一半}。
\end{itemize}

与三相桥式电路的比较（教材结论）：

\begin{enumerate}
  \item 三相桥式电路谐波次数少（$6k\pm1$），双反星形为 $3k\pm1$，但双反星形的\textbf{变压器磁路平衡、无直流磁化}，变压器利用率高；
  \item 双反星形在同样的直流输出电压下，\textbf{变压器二次电压可以低一些}（因两组并联，每组的电压即为输出电压），适用于\textbf{低电压、大电流}场合（如电解、电镀）；
  \item 双反星形需要平衡电抗器，结构较复杂。
\end{enumerate}

\subsection{多重化整流电路}

随着整流装置功率的进一步加大，它所产生的谐波、无功功率等对电网的干扰也随之加大。为减轻干扰，可采用\textbf{多重化整流电路}：将两个或多个相同结构的整流电路（如三相桥）按一定方式组合，使它们的交流侧电流谐波相互抵消。

\begin{itemize}
  \item \textbf{串联二重联结}：两台三相桥的变压器二次绕组分别接成星形和三角形（移相 $30^\circ$），直流侧串联。交流侧电流中 \textbf{5、7 次谐波相互抵消}，只含 $12k\pm1$ 次谐波，构成\textbf{12 脉波}整流电路。
  \item \textbf{串联四重联结}：二次绕组移相 $15^\circ$，构成 \textbf{24 脉波}整流电路，交流侧谐波次数为 $24k\pm1$。
  \item 多重化的效果：位移因数与功率因数提高。12 脉波时 $\lambda\approx0.9886\cos\alpha$，24 脉波时 $\lambda\approx0.9949\cos\alpha$（单桥时为 $0.955\cos\alpha$）。
\end{itemize}

\keyformula{\lambda_{12}\approx0.9886\cos\alpha,\qquad \lambda_{24}\approx0.9949\cos\alpha}

% ============================================================
\section{整流电路的有源逆变工作状态}

\subsection{逆变的概念}

\textbf{逆变}（Inverter）是把直流电变成交流电。按交流侧的去向分为两类：

\begin{itemize}
  \item \textbf{无源逆变}：交流侧直接接负载（如变频器带动电动机），电能直接送给负载，不送回电网——第 4 章讲。
  \item \textbf{有源逆变}：交流侧接电网，即把直流电逆变成与电网同频率的交流电送回电网。
\end{itemize}

同一套晶闸管变流装置，\textbf{整流与逆变是同一装置的两个工作状态}：$\alpha<90^\circ$ 时 $U_d>0$，装置从交流侧吸收有功功率（整流）；$\alpha>90^\circ$ 时 $U_d<0$，装置向交流侧送出有功功率（逆变）。

\begin{table}[htbp]
\centering
\caption{实现有源逆变的条件}
\label{tab:3-4}
\begin{tabular}{p{1.2cm}p{13.4cm}}
\toprule
\textbf{条件} & \textbf{内容} \\
\midrule
（1） & 直流侧必须有一个直流电动势，其\textbf{极性需和晶闸管的导通方向一致}，其值应\textbf{大于变流器直流侧的平均电压} \\
（2） & 要求晶闸管的控制角 $\alpha>\pi/2$，使 $U_d$ 为负值 \\
\bottomrule
\end{tabular}
\end{table}

两者必须同时具备才能实现有源逆变。还要注意：\textbf{半控桥或有续流二极管的电路，因其整流电压不能出现负值，也不允许直流侧出现负极性的电动势，故不能实现有源逆变}——欲实现有源逆变，只能采用全控电路。

\textbf{逆变角} $\beta$ 与控制角 $\alpha$ 的计量方向相反，两者关系为

\keyformula{\alpha+\beta=\pi\quad\text{或}\quad\beta=\pi-\alpha}

\subsection{三相桥整流电路的有源逆变工作状态}

以三相桥式全控电路为例，$\alpha>90^\circ$（即 $\beta<90^\circ$）时，$u_d$ 波形负面积大于正面积，$U_d=-2.34U_2\cos\beta<0$，直流侧电动势 $E$ 与 $U_d$ 同极性（$E$ 驱动电流流入变流器），功率由直流侧流向交流电网。

\begin{itemize}
  \item 逆变状态下，晶闸管承受的电压波形与整流时相同，只是导通区间对应的电压为负；换相仍在自然换相点之后（对逆变而言是之前）发生。
  \item 每个周期内流经电源的线电流导通角仍为 $120^\circ$，波形与整流时形状相同而相位相反。
  \item 逆变时交流侧电流的谐波与功率因数分析方法和整流时一样，但因电流基波与电压的相位差大于 $90^\circ$，\textbf{功率因数为负值}（表示功率倒送）。
\end{itemize}

\subsection{逆变失败与最小逆变角的限制}

\textbf{逆变失败（逆变颠覆）}：逆变运行时一旦发生换相失败，外接的直流电源就会通过晶闸管电路形成短路，或使变流器的输出平均电压和直流电动势变成顺向串联，形成很大的短路电流。

逆变失败的主要原因：

\begin{enumerate}
  \item \textbf{触发电路工作不可靠}，不能适时、准确地给各晶闸管分配脉冲（脉冲丢失、脉冲延迟）；
  \item \textbf{晶闸管发生故障}，失去阻断或导通能力；
  \item \textbf{交流电源缺相或突然消失}；
  \item \textbf{换相的裕量角不足}，引起换相失败（逆变角 $\beta$ 过小）。
\end{enumerate}

为防止逆变失败，必须限制\textbf{最小逆变角} $\beta_{\min}$：

\keyformula{\beta_{\min}=\gamma+\delta+\theta'}

式中 $\gamma$ 为换相重叠角（由 $I_d$、$X_B$、$\alpha$ 决定），$\delta$ 为晶闸管关断时间 $t_q$ 所对应的电角度（$\delta=\omega t_q$，普通晶闸管 $t_q$ 约数百微秒，50 Hz 时 $\delta$ 约为 $4^\circ\sim5^\circ$），$\theta'$ 为安全裕量角（考虑电网电压波动、脉冲不对称等，一般取 $10^\circ$ 左右）。

\corepoint{工程上一般取 $\beta_{\min}=30^\circ\sim35^\circ$。触发电路必须保证 $\beta\ge\beta_{\min}$，并在 $\beta_{\min}$ 处设置\textbf{保护（如“最小逆变角限制”）}。判断题常问“为什么 $\beta$ 不能太小”——因为 $\beta$ 太小则换相尚未完成就反向电压到来，正在导通的晶闸管来不及恢复阻断能力，造成直通短路。}

% ============================================================
\section{整流电路相位控制的实现}

晶闸管可控整流电路是通过控制触发角 $\alpha$ 的大小、即控制触发脉冲起始相位来控制输出电压大小的，属于\textbf{相控电路}。对触发电路的基本要求是：

\begin{enumerate}
  \item 触发脉冲应具有足够的\textbf{幅值和功率}，前沿要陡（上升时间短）；
  \item 脉冲的\textbf{宽度}应保证晶闸管可靠导通（阻感负载需宽脉冲或双窄脉冲）；
  \item 触发脉冲应具有足够的\textbf{移相范围}，并与主电路电源\textbf{同步}；
  \item 各相触发脉冲应\textbf{对称}（间隔均匀）。
\end{enumerate}

\subsection{同步信号为锯齿波的触发电路}

在由模拟电子电路构成的整流装置触发电路中，\textbf{以同步信号为锯齿波的触发电路应用最多}。其基本环节为：

\begin{enumerate}
  \item \textbf{同步环节}：由同步变压器引入与主电路同频率的同步电压，保证锯齿波与主电路电压保持固定相位关系；
  \item \textbf{锯齿波形成}：利用电容恒流充电（或积分电路）在同步电压作用下形成锯齿波，作为移相的比较波形；
  \item \textbf{移相控制}：控制电压 $U_c$ 与锯齿波比较，交点即决定脉冲产生的时刻，改变 $U_c$ 即可连续移相；
  \item \textbf{脉冲形成与放大}：比较后经整形、脉冲变压器隔离与功率放大输出触发脉冲；
  \item \textbf{双窄脉冲环节}：为实现三相全控桥的“双脉冲触发”，各触发单元之间需按一定规律互连线（如 $VT_1$ 单元的输出端接 $VT_2$ 单元的输入端），使每个单元在一个周期内发出两个相隔 $60^\circ$ 的脉冲。
\end{enumerate}

\subsection{集成触发器}

集成电路可靠性高、技术性能好、体积小、功耗低、调试方便。随着集成电路制造技术的提高，晶闸管触发电路的集成化已逐渐普及，现已逐步取代分立式电路。

\begin{itemize}
  \item 国内常用的是 \textbf{KJ004}（旧型号 KC04）移相触发器，一片可输出两路相位互差 $180^\circ$ 的脉冲；\textbf{三片 KJ004 加一片 KJ041（或 KJ042）}即可构成三相全控桥的触发电路。
  \item 集成触发器的移相范围宽、对称度好，但分立元件与模拟集成方案仍有\textbf{脉冲不对称度较高（可达 $3^\circ\sim4^\circ$）、精度低}的缺点。
  \item 在对精度要求高的大容量变流装置中，越来越多地采用\textbf{数字触发电路}（如基于单片机的数字触发器），可获得很好的触发脉冲对称度，精度可达 $0.7^\circ$ 甚至更高。
\end{itemize}

\corepoint{本章的“主线公式”必须能默写并说清适用条件：半波 $0.45U_2\frac{1+\cos\alpha}{2}$（电阻）／$0.9U_2\frac{1+\cos\alpha}{2}$（单相桥电阻）／$0.9U_2\cos\alpha$（单相桥阻感连续）／$1.17U_2\cos\alpha$（三相半波阻感）／$2.34U_2\cos\alpha$（三相桥阻感）。再记住三句判据：\textbf{电阻负载 $u_d$ 不能为负、移相范围到电压自然过零为止；阻感负载电流连续、$u_d$ 可负、移相到 $90^\circ$；要有源逆变必须全控且 $\alpha>90^\circ$}。}

\chapter{逆变电路}

把直流电变成交流电称为\textbf{逆变}。逆变电路与整流电路相对应，是电力电子技术中最重要的变换之一：整流是把交流“整理”成直流，逆变则是把直流“还原”成交流。按交流侧的去向，逆变分为两类——\textbf{有源逆变}（交流侧接电网，已见 3.7 节）与\textbf{无源逆变}（交流侧直接接负载，如变频器带电动机、不间断电源带负载）。本章讲的是\textbf{无源逆变}。

逆变电路的核心问题是\textbf{换流}：电流要从一个桥臂转移到另一个桥臂，能不能可靠转移，决定了电路能不能工作。

% ============================================================
\section{换流方式}

\textbf{换流}：电流从一个桥臂向另一个桥臂转移的过程。按促使晶闸管（或全控器件）关断的机理不同，换流方式分为四种。

\subsection{器件换流}

利用\textbf{全控型器件}的自关断能力进行换流。由于 IGBT、电力 MOSFET、GTO、GTR 等器件都能由控制信号关断，因此由它们构成的电路全部采用器件换流。

\begin{itemize}
  \item 优点：不需要附加换流电路，主电路简单、效率高；
  \item 这是现代逆变电路的主流方式（第 4 章后面各节，除特别说明外均为器件换流）。
\end{itemize}

\subsection{电网换流}

由\textbf{电网}提供换流电压——只要把负的电网电压加在欲关断的晶闸管上即可使其关断。这种方式不需要器件具有自关断能力，也不需要附加换流电路，经济方便；但\textbf{只能用于由交流电网供电的电路}（如相控整流电路、有源逆变电路、交交变频电路），不能用于无源逆变。

\subsection{负载换流}

由\textbf{负载}提供换流电压。条件是\textbf{负载电流超前于负载电压}（容性负载），这样在电流过零之前电压已经反向，使欲关断的晶闸管承受反压而关断。

典型应用是\textbf{并联谐振式逆变电路}（用于感应加热、中频电源）：负载为感应线圈（$L$、$r$ 串联）与补偿电容 $C$ 并联，电路工作在接近谐振的状态，负载对基波呈高阻抗、对谐波呈低阻抗，因此负载电压接近正弦波。为保证可靠换流，触发脉冲必须\textbf{引前}于负载电压一个角度 $\beta$（逆变角／触发引前角），

\keyformula{\beta>\gamma+\delta}

式中 $\gamma$ 为换相重叠角，$\delta$ 为晶闸管关断时间所对应的电角度。$\beta$ 太小则换流失败（逆变颠覆）；$\beta$ 太大则输出功率因数下降、输出电压降低。实际取 $\beta$ 略大于 $\gamma+\delta$ 的最小值。

\subsection{强迫换流}

设置\textbf{附加的换流电路}，由换流电路内的电容（或电感）的储能向欲关断的晶闸管施加反向电压或反向电流使其关断。由于换流电压／电流是“强行”加上的，故称强迫换流，也称脉冲换流。

\begin{itemize}
  \item 直接耦合式强迫换流：由换流电路内电容直接提供换流电压；
  \item 电感耦合式强迫换流：通过换流电路内电容和电感耦合提供换流电压或换流电流。
\end{itemize}

强迫换流需要附加电路，增加了成本与损耗，但对晶闸管电路而言是唯一可用于无源逆变的换流方式（晶闸管本身不能自关断，又没有电网或超前负载电压可用）。

\begin{table}[htbp]
\centering
\caption{四种换流方式的比较}
\label{tab:4-1}
\begin{tabular}{p{2.4cm}p{4.4cm}p{7.6cm}}
\toprule
\textbf{换流方式} & \textbf{换流能量来源} & \textbf{适用场合} \\
\midrule
器件换流 & 全控器件的控制信号 & 由 IGBT、MOSFET、GTO、GTR 构成的各类电路 \\
电网换流 & 交流电网电压 & 相控整流、有源逆变、交交变频（必须接电网） \\
负载换流 & 负载的超前电压 & 并联谐振式逆变、感应加热中频电源（负载须容性） \\
强迫换流 & 附加电路的储能（电容/电感） & 晶闸管无源逆变电路（需附加换流电路） \\
\bottomrule
\end{tabular}
\end{table}

% ============================================================
\section{电压型逆变电路}

\textbf{电压型逆变电路}（Voltage Source Inverter，VSI）的直流侧是\textbf{电压源}，用大电容储能并使直流电压基本无脉动，因此交流侧输出电压为\textbf{矩形波}（或阶梯波），而电流波形取决于负载。

\begin{itemize}
  \item 直流侧并联大电容，相当于电压源，直流侧电压基本无脉动，直流回路呈现\textbf{低阻抗}；
  \item 交流侧输出电压波形为矩形波，与负载阻抗角无关；输出电流波形和相位因负载不同而不同；
  \item 需要器件换流（全控器件）或强迫换流（晶闸管）；
  \item \textbf{当交流侧为阻感负载时需要提供无功功率}，直流侧电容起缓冲无功能量的作用；为给交流侧向直流侧反馈的无功能量提供通道，各桥臂的开关器件\textbf{都须反并联二极管（续流二极管）}。
\end{itemize}

\subsection{单相电压型逆变电路}

\subsubsection*{1. 半桥逆变电路}

由两只开关器件（$S_1$、$S_2$）与两只大电容串联组成，负载接在两电容中点和两器件中点之间。$S_1$、$S_2$ 交替导通、各导通 $180^\circ$，输出方波电压的幅值为 $U_d/2$。

\keyformula{U_o=\frac{U_d}{2},\qquad U_{o1m}=\frac{2U_d}{\pi},\qquad U_{o1}=\frac{\sqrt{2}U_d}{\pi}=0.45U_d}

其中 $U_{o1m}$ 为基波幅值、$U_{o1}$ 为基波有效值。半桥电路使用器件少，但直流侧需两只电容分压，且实际中两电容电压可能不均衡（需均压措施），故只用于几千瓦以下的小功率场合。

\subsubsection*{2. 全桥逆变电路}

由四个开关器件组成两个桥臂，可看成两个半桥电路的组合。$S_1$、$S_4$ 为一对（同时导通），$S_2$、$S_3$ 为另一对，两对交替各导通 $180^\circ$，输出方波电压幅值为 $U_d$（是半桥的两倍）。

\keyformula{U_o=U_d,\qquad U_{o1m}=\frac{4U_d}{\pi},\qquad U_{o1}=\frac{2\sqrt{2}U_d}{\pi}=0.9U_d}

\begin{figure}[htbp]
\centering
\begin{tikzpicture}[font=\footnotesize,scale=0.95]
  \draw[->] (0,0) -- (8.4,0) node[below] {$\omega t$};
  \draw[->] (0,-2.2) -- (0,2.6) node[left] {$u_o$};
  \draw[yhbAccent,thick] (0,2.0) -- (4,2.0) -- (4,-2.0) -- (8,-2.0) -- (8,2.0) -- (8.3,2.0);
  \draw[dashed,textgray] (0,2.0) node[left] {\scriptsize $U_d$} -- (0,2.0);
  \draw[dashed,textgray] (0,-2.0) node[left] {\scriptsize $-U_d$};
  \draw[yhbAccent2,thick] plot[domain=0:8,samples=120] (\x,{2*4/3.1416*sin(57.3*\x)});
  \node[text=yhbAccent2] at (6.8,2.5) {\scriptsize 基波};
  \node[text=yhbAccent] at (0.75,2.35) {\scriptsize 方波};
\end{tikzpicture}
\caption{单相全桥电压型逆变电路的输出电压波形}
\label{fig:4-1}
\end{figure}

方波电压含有的谐波为 $n=2k\pm1$（奇次），各次谐波幅值与次数成反比：$U_{onm}=\dfrac{4U_d}{n\pi}$。因此\textbf{要改变输出电压有效值，只能改变直流电压 $U_d$}；要减少谐波，则须采用 PWM 技术（第 7 章）或多重化。

\subsection{三相电压型逆变电路}

三个桥臂、六只开关器件组成三相桥，是应用最广的逆变电路。常用\textbf{180$^\circ$ 导电方式}：每个桥臂的上下两管互补导通各 $180^\circ$，三个桥臂的导通依次相差 $120^\circ$，因此\textbf{任何时刻都有三个桥臂同时导通}（同一桥臂上下管绝不能同时导通，否则直流侧短路，实际须留“死区”）。

输出波形：

\begin{itemize}
  \item \textbf{线电压}为 $120^\circ$ 矩形波（幅值 $U_d$），因为线电压只取决于两相桥臂的状态；
  \item \textbf{相电压}（星形负载）为\textbf{六阶梯波}，幅值为 $\pm2U_d/3$ 与 $\pm U_d/3$ 交替。
\end{itemize}

\keyformula{U_{UN}=\frac{\sqrt{2}}{3}U_d=0.471U_d,\quad U_{UN1m}=\frac{2U_d}{\pi}=0.637U_d,\quad U_{UN1}=\frac{\sqrt{2}U_d}{\pi}=0.45U_d}

\keyformula{U_{UV}=\sqrt{\frac{2}{3}}U_d=0.816U_d,\quad U_{UV1m}=\frac{2\sqrt{3}U_d}{\pi}=1.1U_d,\quad U_{UV1}=\frac{\sqrt{6}U_d}{\pi}=0.78U_d}

\corepoint{三相电压型桥式逆变的三个“0.471 / 0.816 / 0.45 / 0.78”必须记准：\textbf{相电压有效值 0.471$U_d$、线电压有效值 0.816$U_d$、相电压基波有效值 0.45$U_d$、线电压基波有效值 0.78$U_d$}。还要记住两点：（1） 同一桥臂上下管\textbf{互补导通}，换流是“器件换流”，但必须设\textbf{死区}防止直通；（2） 负载为阻感时，无功能量经\textbf{反并联二极管}回馈到直流侧电容。}

% ============================================================
\section{电流型逆变电路}

\textbf{电流型逆变电路}（Current Source Inverter，CSI）的直流侧是\textbf{电流源}，用大电感储能并使直流电流基本无脉动，因此交流侧输出电流为\textbf{矩形波}，而电压波形取决于负载。

\begin{itemize}
  \item 直流侧串联大电感，相当于电流源，直流回路呈现\textbf{高阻抗}；
  \item 交流侧输出电流波形为矩形波，与负载阻抗角无关；输出电压波形和相位因负载不同而不同；
  \item 因为直流侧电感储能，\textbf{无功能量由电感缓冲}，不需要反并联二极管（这一点与电压型相反）；
  \item 换流时须采取\textbf{“先断后通”}的方法：先给应关断的器件关断信号，待其关断并留一定时间裕量后，再给应导通的器件发开通信号，否则将造成直流侧电源短路。
\end{itemize}

\subsection{单相电流型逆变电路}

典型电路是\textbf{并联谐振式逆变电路}（用于感应加热、中频电源）：四个桥臂的晶闸管各串联一个电抗器 $L_T$（限制开通时的 di/dt），负载为感应线圈（$L$、$r$ 串联）与补偿电容 $C$ 并联，构成并联谐振电路。

\begin{itemize}
  \item 工作频率接近负载电路的谐振频率，负载对基波呈\textbf{高阻抗}、对谐波呈低阻抗，故\textbf{负载电压接近正弦波}；
  \item 换流方式为\textbf{负载换流}：负载电流超前于负载电压（容性），电流过零前电压已反向，使欲关断的晶闸管承受反压而关断，因此必须使触发脉冲\textbf{引前}负载电压一个角度 $\beta$；
  \item 由于 $L_T$ 的存在，换相需要时间，形成\textbf{换相重叠角} $\gamma$；负载的功率因数角 $\varphi$ 与 $\beta$、$\gamma$ 的关系决定输出电压大小，$\beta$ 越大输出电压越低（功率因数越低）；
  \item 起动问题：需先给电容充电（预充电），起动时把电容能量释放到负载形成衰减振荡，检测出振荡信号实现自励。
\end{itemize}

\subsection{三相电流型逆变电路}

典型电路是三相桥式电流型逆变电路，采用\textbf{120$^\circ$ 导电方式}：每个桥臂导通 $120^\circ$，任何时刻只有\textbf{两个桥臂同时导通}（一个共阴组、一个共阳组），输出电流为 $120^\circ$ 矩形波。

\keyformula{I=\sqrt{\frac{2}{3}}I_d=0.816I_d,\qquad I_1=\frac{\sqrt{6}I_d}{\pi}=0.78I_d}

其中 $I$ 为输出电流有效值，$I_1$ 为基波有效值。各桥臂之间换流采用\textbf{强迫换流}，连接于各桥臂之间的电容 $C_1\sim C_6$ 即为换流电容；换流过程分“恒流放电阶段”（放电电流恒为 $I_d$）与“谐振放电阶段”等。由于是强迫换流，换流电容的容量必须足以关断晶闸管，故工作频率不能太高。

\begin{table}[htbp]
\centering
\caption{电压型与电流型逆变电路的比较}
\label{tab:4-2}
\begin{tabular}{p{2.6cm}p{6.0cm}p{6.0cm}}
\toprule
\textbf{比较项} & \textbf{电压型（VSI）} & \textbf{电流型（CSI）} \\
\midrule
直流侧储能元件 & 大电容（并联），电压源性质、低阻抗 & 大电感（串联），电流源性质、高阻抗 \\
输出波形 & 输出电压为矩形波，电流取决于负载 & 输出电流为矩形波，电压取决于负载 \\
反并联二极管 & 必须有（为无功能量提供通路） & 不需要 \\
换流方式 & 器件换流（或强迫换流） & 负载换流、强迫换流；须“先断后通” \\
短路保护 & 直流侧短路电流大，保护较难 & 电感限流，短路电流上升慢，保护容易 \\
动态响应 & 快 & 较慢（电感电流不能突变） \\
典型应用 & 通用变频器、UPS、开关电源 & 感应加热、中频电源、大容量电动机调速 \\
\bottomrule
\end{tabular}
\end{table}

% ============================================================
\section{多重逆变电路和多电平逆变电路}

前面介绍的逆变电路输出都是矩形波，含大量低次谐波。改善输出波形的途径有两条：一是\textbf{把几个矩形波按一定相位差组合起来}，使某些主要谐波相互抵消——\textbf{多重逆变电路}；二是\textbf{改变电路结构，输出更多的电平}，使波形向正弦波靠近——\textbf{多电平逆变电路}。

\subsection{多重逆变电路}

把若干个逆变电路的输出按一定的相位差组合起来，使它们所含的某些主要谐波分量相互抵消，就可以得到较为接近正弦波的波形。

\begin{itemize}
  \item \textbf{二重逆变电路}：两个单相全桥逆变电路（输出电压相位错开一定角度）经变压器合成，输出波形为阶梯波，最低次谐波次数提高、幅值减小；两个三相桥式逆变电路相位错开 $30^\circ$ 合成，可构成 \textbf{12 脉波}逆变电路。
  \item 一般规律：使 $m$ 个三相桥式逆变电路的相位依次错开 $\pi/(3m)$ 运行，连同使它们输出电压合成并抵消上述相位差的变压器，就可以构成\textbf{脉波数为 $6m$ 的逆变电路}（二重 12 脉波、三重 18 脉波）。
  \item 多重化的效果：输出波形中最低次谐波次数提高到 $6m\pm1$，谐波幅值随次数成反比减小，故波形明显改善；同时每个逆变单元承受的功率减小，便于用较低电压等级的器件构成大容量装置。
\end{itemize}

\subsection{多电平逆变电路}

\textbf{多电平逆变电路}能够输出较多的电平，从而使输出电压向正弦波靠近。常见结构有三种：

\begin{enumerate}
  \item \textbf{中点钳位型（NPC）三电平逆变电路}：每相桥臂由四个开关器件串联，两只钳位二极管把桥臂中点与直流侧电容中点相连。以 U 相为例，当 $V_{11}$、$V_{12}$（或 $VD_{11}$、$VD_{12}$）导通、$V_{41}$、$V_{42}$ 关断时，U 点与直流中点（$\varphi$ 点）之间的电压为 $+U_d/2$；反之则为 $-U_d/2$；上下各一只导通时输出为 0。因此每相可输出 $+U_d/2$、$0$、$-U_d/2$ 三种电平，线电压有五种电平。
  \item \textbf{飞跨电容型（FLC）三电平逆变电路}：用飞跨电容代替钳位二极管，结构稍简单，但需要控制电容电压平衡，使用较少；构成更多电平时所需电容数目急剧增加。
  \item \textbf{单元串联多电平逆变电路}：把多个独立的低压逆变单元（H 桥）串联叠加，用移相变压器分别供电，可获得很高的输出电压与很多电平数，谐波小、可靠性高（单元故障可旁路），在高压大容量变频器中应用较多。
\end{enumerate}

\corepoint{本章的考点集中在三处：\textbf{（1） 四种换流方式的判别}（看到“全控器件”想到器件换流，看到“接电网”想到电网换流，看到“负载电流超前”想到负载换流，看到“附加电容”想到强迫换流）；\textbf{（2） 电压型与电流型的对比}（电容／电感、是否要反并联二极管、是否“先断后通”）；\textbf{（3） 输出电压有效值与基波的四个系数}（单相全桥 $0.9U_d$、三相相电压 $0.45U_d$ 与 $0.471U_d$、线电压 $0.78U_d$ 与 $0.816U_d$）。要改输出电压有效值只能改 $U_d$，要改波形只能靠 PWM 或多重化。}

\chapter{直流—直流变流电路}

直流—直流变流电路（DC-DC Converter）把一种直流电压变换成另一种固定或可调的直流电压，也称\textbf{直流斩波电路}（DC Chopper）。它的基本手段是“开关 + 储能元件”：用开关把直流“切”成脉冲，再用 $L$、$C$ 滤波或被控对象的惯性把脉冲“抹平”，通过改变开关的\textbf{占空比}（导通时间与周期的比）来调节输出平均电压。

本章的分析套路固定为三步：\textbf{（1） 分阶段画等效电路（开关导通／关断）；（2） 对电感用伏秒平衡 $\int u_L\,\mathrm{d}t=0$（稳态时电感电压平均值为零）；（3） 由输入输出功率平衡或电容安秒平衡求电流关系}。

% ============================================================
\section{基本斩波电路}

\subsection{降压斩波电路}

降压斩波电路（Buck Chopper）是最基础的一种：输出电压平均值低于输入电压。

\begin{center}
\begin{circuitikz}[font=\footnotesize,scale=0.9]
  \draw (0,0) to[battery1, l=$E$] (0,-2.6);
  \draw (0,0) -- (1.0,0) to[switch, l=$S$] (2.8,0) -- (3.8,0) -- (5.6,0) to[L, l=$L$] (7.4,0) -- (8.8,0);
  \draw (3.8,-2.6) to[D] (3.8,0);
  \draw (6.6,0) to[C, l=$C$] (6.6,-2.6);
  \draw (8.4,0) to[R, l=$R$] (8.4,-2.6);
  \draw (0,-2.6) -- (8.8,-2.6);
  \draw (5.6,0) -- (6.6,0);
  \node[right] at (3.95,-1.9) {\scriptsize $VD$};
  \node[below] at (7.5,-2.75) {\scriptsize $u_o$};
\end{circuitikz}
\end{center}

\textbf{工作过程}（假设电感足够大、电流连续，$T$ 为开关周期，$t_{on}$ 为导通时间，占空比 $\alpha=t_{on}/T$）：

\begin{itemize}
  \item $S$ 导通（$0\sim t_{on}$）：电源 $E$ 经 $L$ 向负载供电，$u_o=E$，电感电流线性上升，$L$ 储能；二极管 $VD$ 反偏截止。
  \item $S$ 关断（$t_{on}\sim T$）：$L$ 的感应电动势使 $VD$ 正偏导通，电感经 $VD$—负载续流，$u_o=0$，电流线性下降。
\end{itemize}

对电感用伏秒平衡：$E\,t_{on}=(E-u_o)\,t_{off}$ 即得

\keyformula{U_o=\frac{t_{on}}{T}E=\alpha E}

\textbf{电流连续}时输出电流 $I_o=U_o/R$（或对电动机负载 $I_o=(U_o-E_M)/R$），输入电流平均值 $I_i=\alpha I_o$，输入输出功率相等 $EI_i=U_oI_o$。

\textbf{电流断续}时（负载很轻、$L$ 较小），$S$ 管再次导通前 $L$ 中电流已到零，$u_o$ 在断续期间等于反电动势 $E_M$（阻感反电动势负载），此时

\keyformula{U_o=\frac{t_{on}}{t_{on}+t_{off}'}E}

其中 $t_{off}'$ 为电流实际衰减到零的时间（$t_{off}'<t_{off}$），因此\textbf{断续时输出电压比连续时高}，且不再与 $\alpha$ 成线性关系——这是分析斩波电路时最容易出错的地方。

\corepoint{判断电流是否连续的方法：先按连续公式算 $U_o=\alpha E$ 与 $I_o=U_o/R$，再算电流脉动峰值 $\Delta I=\dfrac{E(1-\alpha)\alpha T}{L}$（或按 $U_o(1-\alpha)T/L$），若 $\dfrac{\Delta I}{2}\ge I_o$ 则电流断续。临界电感为 $L_{cr}=\dfrac{(1-\alpha)U_oT}{2I_o}=\dfrac{(1-\alpha)RT}{2}$，所以\textbf{负载越轻（$R$ 越大）、$L$ 越小、频率越低，越容易断续}。}

\subsection{升压斩波电路}

升压斩波电路（Boost Chopper）的输出电压高于输入电压，拓扑上把电感放在输入端、把开关放在并联位置。

\begin{center}
\begin{circuitikz}[font=\footnotesize,scale=0.9]
  \draw (0,0) to[battery1, l=$E$] (0,-2.6);
  \draw (0,0) -- (2.2,0) to[L, l=$L$] (4.0,0) -- (5.0,0) to[D] (6.4,0) -- (7.8,0);
  \draw (5.0,-2.6) to[switch, l=$S$] (5.0,0);
  \draw (6.4,0) -- (8.8,0);
  \draw (7.4,0) to[C, l=$C$] (7.4,-2.6);
  \draw (8.6,0) to[R, l=$R$] (8.6,-2.6);
  \draw (0,-2.6) -- (8.8,-2.6);
\end{circuitikz}
\end{center}

\begin{itemize}
  \item $S$ 导通（$t_{on}$）：电源经 $L$ 短路，$L$ 储能、电流上升，二极管截止，负载由电容 $C$ 供电。
  \item $S$ 关断（$t_{off}$）：$L$ 的感应电动势与 $E$ 串联向负载和电容供电，$u_o=E+u_L>E$；$L$ 释能，电流下降。
\end{itemize}

电感伏秒平衡：$E\,t_{on}=(U_o-E)t_{off}$，解得

\keyformula{U_o=\frac{T}{t_{off}}E=\frac{1}{1-\alpha}E}

由于 $\dfrac{1}{1-\alpha}\ge1$，输出电压\textbf{总是高于输入电压}，$\alpha\to1$ 时 $U_o\to\infty$（实际受寄生参数限制）。由功率平衡 $EI_i=U_oI_o$ 得 $I_i=\dfrac{I_o}{1-\alpha}$，即\textbf{输入电流大于输出电流}。

\textbf{电流断续}时 $U_o=\dfrac{t_{on}+t_{off}'}{t_{off}'}E$，输出电压比连续时更高。

\corepoint{Buck 与 Boost 的对偶关系要一次记清：\textbf{Buck 的 $U_o=\alpha E$、输入电流断续（$I_i$ 为脉冲而 $I_o$ 连续）；Boost 的 $U_o=E/(1-\alpha)$、输出电流断续（$I_o$ 为脉冲而 $I_i$ 连续）}。Boost 电路\textbf{不能空载运行}（空载时电容电压会一直充到很高而损坏器件），这是它与 Buck 的重要区别。}

\subsection{升降压斩波电路和 Cuk 斩波电路}

\subsubsection*{1. 升降压斩波电路（Buck-Boost）}

把 Buck 与 Boost 串接并共用一只开关和一只二极管，得到升降压斩波电路。设 $S$ 导通时电感从电源储能（此时负载由电容供电），$S$ 关断时电感经二极管向负载释能，则

\keyformula{U_o=-\frac{t_{on}}{t_{off}}E=-\frac{\alpha}{1-\alpha}E}

负号表示\textbf{输出极性与输入相反}（反极性）。$\alpha<0.5$ 时 $|U_o|<E$（降压），$\alpha>0.5$ 时 $|U_o|>E$（升压），故称升降压。由 $EI_i=|U_o|I_o$ 得 $I_i=\dfrac{\alpha}{1-\alpha}I_o$。

该电路的优点是\textbf{电压调节范围宽}，缺点是\textbf{输入、输出电流都是脉动的}，滤波电容 $C$ 需承受较大的脉动电流。

\subsubsection*{2. Cuk 斩波电路}

Cuk 斩波电路同样输出反极性电压、电压比也是 $-\alpha/(1-\alpha)$，但它把电感放在\textbf{输入和输出两端}，中间用电容 $C_1$ 传递能量：

\keyformula{U_o=-\frac{\alpha}{1-\alpha}E}

\begin{itemize}
  \item $S$ 导通时：$E$—$L_1$—$S$ 回路，$L_1$ 储能；同时 $C_1$ 经 $S$—$L_2$—负载放电（$L_2$ 储能）。
  \item $S$ 关断时：$E$—$L_1$—$C_1$—$VD$ 回路，$L_1$ 向 $C_1$ 充电；$L_2$ 经 $VD$—负载续流。
  \item 由于 $L_1$、$L_2$ 分列输入输出两侧，\textbf{输入电流和输出电流都是连续的}，脉动小，这是 Cuk 相对 Buck-Boost 的最大优点；代价是需一个承受高脉动电流的传递电容 $C_1$，且电路较复杂。
\end{itemize}

\subsection{Sepic 斩波电路和 Zeta 斩波电路}

这两种电路可以看成 Cuk 电路的变形，其电压比与 Buck-Boost 相同但\textbf{输出极性与输入相同}（非反相）。

\keyformula{U_o=\frac{\alpha}{1-\alpha}E\quad(\text{Sepic、Zeta})}

\begin{itemize}
  \item \textbf{Sepic 电路}：电感在输入端、二极管在输出端，中间经电容耦合；输入电流连续。
  \item \textbf{Zeta 电路}：电感在输出端、二极管在输入端；输出电流连续。
\end{itemize}

\begin{table}[htbp]
\centering
\caption{基本斩波电路的电压比与特点}
\label{tab:5-1}
\begin{tabular}{p{2.6cm}p{3.4cm}p{1.6cm}p{2.4cm}p{3.6cm}}
\toprule
\textbf{电路} & \textbf{电压比 $U_o/E$} & \textbf{极性} & \textbf{连续量} & \textbf{主要特点} \\
\midrule
Buck（降压） & $\alpha$ & 同相 & 输出电流连续 & 结构最简单，输出电流脉动大，输入电流断续 \\
Boost（升压） & $1/(1-\alpha)$ & 同相 & 输入电流连续 & 输出电流断续，不能空载 \\
Buck-Boost（升降压） & $-\alpha/(1-\alpha)$ & 反相 & 均断续 & 调节范围宽，输入输出电流都脉动 \\
Cuk & $-\alpha/(1-\alpha)$ & 反相 & 均连续 & 输入输出电流连续，需传递电容 \\
Sepic & $\alpha/(1-\alpha)$ & 同相 & 输入电流连续 & 非反相的升降压 \\
Zeta & $\alpha/(1-\alpha)$ & 同相 & 输出电流连续 & Sepic 的对偶 \\
\bottomrule
\end{tabular}
\end{table}

% ============================================================
\section{复合斩波电路和多相多重斩波电路}

\subsection{电流可逆斩波电路}

把降压斩波电路与升压斩波电路组合起来，就可构成\textbf{电流可逆斩波电路}：$V_1$、$VD_1$ 构成降压斩波（电动机工作状态，能量由电源流向电动机），$V_2$、$VD_2$ 构成升压斩波（再生制动状态，能量由电动机回馈到电源）。

\begin{itemize}
  \item 使 $V_1$ 工作、$V_2$ 不工作：电动机电动运行，$I_o>0$，$U_o=\alpha_1E$；
  \item 使 $V_2$ 工作、$V_1$ 不工作：电动机反接制动或再生制动，电流反向，能量回馈电源；
  \item $V_1$、$V_2$ 交替工作（两者不能同时导通，否则电源短路），可实现\textbf{电流的双向流动}——这是“单象限”变“两象限”的关键。
\end{itemize}

\subsection{桥式可逆斩波电路}

把两个电流可逆斩波电路组合，构成\textbf{桥式可逆斩波电路}（四象限斩波器）：四个桥臂全由可控器件与反并联二极管组成，电动机接在桥的中点。

\begin{itemize}
  \item 通过控制四个桥臂的导通组合，可使电动机的\textbf{电压和电流都能正负双向}，从而运行在\textbf{四个象限}：正转电动、正转制动、反转电动、反转制动。
  \item 必须保证\textbf{同一桥臂上下管不能同时导通}（否则电源直通短路），故驱动信号须留死区、并设互锁逻辑。
\end{itemize}

\subsection{多相多重斩波电路}

把 $m$ 个结构相同的斩波电路并联（相），每个斩波电路在一个周期内导通 $n$ 次（重），即构成 $m$ 相 $n$ 重斩波电路。

\begin{itemize}
  \item 各相开关的驱动脉冲依次错开 $T/m$，使\textbf{输入电流和输出电流的脉动频率提高到 $m$ 倍}，脉动幅值大大减小，因而可减小滤波元件、提高功率密度；
  \item 同时每个单元承担的电流减小，便于用较小容量的器件构成大电流装置；
  \item 常用于低压大电流场合（如电解电源、电镀电源）。
\end{itemize}

% ============================================================
\section{带隔离的直流—直流变流电路}

前述电路的输入输出共地、不隔离；在需要\textbf{电气隔离}的场合（安全、抗干扰），必须通过高频变压器耦合，这就是带隔离的直流—直流变流电路。按变压器励磁方式分为两类：

\begin{itemize}
  \item \textbf{单端电路}：变压器磁通只沿磁滞回线的一个分支变化（单向励磁），如正激、反激；需专门设置磁复位电路，否则磁通会逐周期累积而饱和。
  \item \textbf{双端电路}：变压器双向励磁，磁通在正负两个方向交替变化，利用率高，如半桥、全桥、推挽。
\end{itemize}

\subsection{正激电路}

正激（Forward）电路：开关 $S$ 导通时，电源经变压器向二次侧传输能量，二次侧经整流二极管向负载供电并给电感储能；$S$ 关断时，电感经续流二极管续流，变压器需\textbf{复位绕组}把励磁能量送回电源（否则磁通累积饱和）。

\begin{itemize}
  \item 输出电压 $U_o=\alpha\dfrac{N_2}{N_1}E$（$N_1$、$N_2$ 为一、二次匝数）；
  \item 因复位绕组的存在，\textbf{占空比必须小于 0.5}（复位时间需占一定比例），否则变压器饱和；
  \item 优点：结构简单、二次侧电流脉动小；缺点：变压器单向励磁、利用率低、需复位电路。
\end{itemize}

\subsection{反激电路}

反激（Flyback）电路：$S$ 导通时变压器一次绕组储能，二次侧二极管截止、负载由输出电容供电；$S$ 关断时一次绕组能量耦合到二次侧，经二极管向负载和电容释放。

\begin{itemize}
  \item 输出电压 $U_o=\dfrac{\alpha}{1-\alpha}\cdot\dfrac{N_2}{N_1}E$（连续模式）；
  \item 反激电路的变压器实际起\textbf{储能电感}的作用，磁通单向变化，\textbf{不需要复位绕组}，占空比可大于 0.5；
  \item 优点：结构最简单、成本最低，广泛用于小功率电源（手机充电器、适配器）；缺点：变压器利用率低、输出纹波较大。
\end{itemize}

\subsection{半桥电路}

半桥电路用两只开关器件与两只电容串联，变压器一次绕组接在两中点之间。

\begin{itemize}
  \item 变压器一次电压幅值为 $E/2$，\textbf{双向励磁}，无磁通累积问题；
  \item 同一桥臂上下管须留死区防直通；电容中点电压需均衡；
  \item 输出功率中等，用于几百瓦到几千瓦的开关电源。
\end{itemize}

\subsection{全桥电路}

全桥电路用四只开关器件组成两个桥臂，变压器一次绕组接在两桥臂中点之间。

\begin{itemize}
  \item 变压器一次电压幅值为 $E$（是半桥的两倍），同样双向励磁；
  \item 优点：功率大、变压器利用率高；缺点：器件多、驱动复杂、有直通风险；
  \item 是大功率开关电源、逆变焊机、电力操作电源的常用拓扑。
\end{itemize}

\subsection{推挽电路}

推挽（Push-Pull）电路用两只开关器件与\textbf{带中心抽头}的变压器一次绕组构成。

\begin{itemize}
  \item 两只开关交替导通，变压器一次两个半绕组轮流工作，磁通双向变化；
  \item 每个开关承受的电压为 $2E$（关断时另一半绕组的感应电动势叠加），故适合低压输入场合；
  \item 缺点是变压器需中心抽头、开关耐压要求高、易出现“偏磁”（两半绕组不对称导致直流磁化）。
\end{itemize}

\subsection{全波整流和全桥整流}

双端电路中常用的二次侧整流形式为\textbf{全波整流}（变压器二次带中心抽头，两只二极管）与\textbf{全桥整流}（四只二极管，无需中心抽头）。

\begin{table}[htbp]
\centering
\caption{带隔离直流—直流变流电路比较}
\label{tab:5-2}
\begin{tabular}{p{2.4cm}p{2.4cm}p{2.6cm}p{2.0cm}p{4.0cm}}
\toprule
\textbf{电路} & \textbf{变压器励磁} & \textbf{一次电压幅值} & \textbf{占空比限制} & \textbf{特点} \\
\midrule
正激 & 单向 & $E$ & $<0.5$ & 需复位绕组，结构简单，中小功率 \\
反激 & 单向（储能） & $E$ & 无 & 不需复位，最简单最便宜，小功率 \\
半桥 & 双向 & $E/2$ & $<0.5$（防直通） & 器件少，中功率 \\
全桥 & 双向 & $E$ & $<0.5$（防直通） & 功率大、利用率高，器件多 \\
推挽 & 双向 & $E$ & $<0.5$（防偏磁） & 需中心抽头，开关耐压 $2E$，低压输入 \\
\bottomrule
\end{tabular}
\end{table}

\corepoint{本章最容易失分的三点：\textbf{（1） 电压比记混}——Buck 是 $\alpha$、Boost 是 $1/(1-\alpha)$、Buck-Boost 与 Cuk 是 $-\alpha/(1-\alpha)$、Sepic 与 Zeta 是 $+\alpha/(1-\alpha)$；\textbf{（2） 忽略“电流断续”}——所有用 $\alpha$ 表示的线性关系只在电流连续时成立，断续时输出电压偏高，必须先判连续；\textbf{（3） 忘掉“极性”}——Buck-Boost 与 Cuk 输出反极性，Sepic 与 Zeta 输出同极性。另外记住隔离电路中正激需复位且 $D<0.5$、反激不需复位。}

\chapter{交流—交流变流电路}

把一种形式的交流电变成另一种形式的交流电（改变电压、电流，或改变频率、相数）称为\textbf{交流—交流变流}。只改变电压或电流而不改变频率的，称为\textbf{交流电力控制电路}，其中输出电压可连续调节的称为\textbf{交流调压电路}；既改变电压又改变频率的，称为\textbf{交—交变频电路}（周波变流器）。

% ============================================================
\section{交流调压电路}

交流调压电路用两只晶闸管反并联（或用一只双向晶闸管）串在负载回路中，在每个半周内控制其导通起始时刻，从而改变负载上电压的有效值。与相控整流不同，交流调压的负载直接接在交流电源上，\textbf{没有整流环节}。

\subsection{单相交流调压电路}

\subsubsection*{1. 电阻负载}

设电源电压 $u_1=\sqrt{2}U_1\sin\omega t$，晶闸管在 $\alpha$ 处触发导通，到 $\pi$ 处电流过零关断；负半周由反并联的管子重复。负载电压为缺角的正弦波。

\keyformula{U_o=U_1\sqrt{\frac{1}{\pi}\left(\pi-\alpha+\frac{\sin 2\alpha}{2}\right)}=U_1\sqrt{\frac{\pi-\alpha}{\pi}+\frac{\sin2\alpha}{2\pi}}}

$\alpha=0$ 时 $U_o=U_1$（全导通）；$\alpha=\pi$ 时 $U_o=0$。故\textbf{移相范围为 $0\sim\pi$}。负载电流 $I_o=U_o/R$，负载消耗的功率 $P=U_o^2/R$；由于负载为纯电阻，\textbf{功率因数就等于电压比}：

\keyformula{\lambda=\frac{P}{S}=\frac{U_oI_o}{U_1I_o}=\frac{U_o}{U_1}=\sqrt{\frac{\pi-\alpha}{\pi}+\frac{\sin2\alpha}{2\pi}}}

\subsubsection*{2. 阻感负载}

阻感负载时电流滞后于电压，晶闸管的导通角 $\theta$ 由 $\alpha$ 与负载阻抗角 $\varphi=\arctan(\omega L/R)$ 共同决定，满足超越方程

\keyformula{\sin(\alpha+\theta-\varphi)=\sin(\alpha-\varphi)\,\mathrm{e}^{-\frac{\theta}{\omega\tau}},\qquad \tau=\frac{L}{R}}

\corepoint{阻感负载的交流调压有一个\textbf{极重要的结论}：若 $\alpha<\varphi$，则触发脉冲到来时电流尚未过零（上一半周的电流还在流），晶闸管无法关断，电路实际上\textbf{失去控制作用}，输出电压等于电源电压（不可控）；因此\textbf{必须保证 $\alpha\ge\varphi$}，即最小触发角为 $\varphi$。若要既能调压又能用窄脉冲触发，须采用\textbf{宽脉冲}或\textbf{脉冲列}（因电流可能延续到脉冲消失之后）。}

\subsubsection*{3. 晶闸管电流定额}

单相交流调压中每只晶闸管流过半个周期的负载电流，其有效值 $I_T=I_o/\sqrt{2}$。选用器件时须按此有效值折算到通态平均电流（见 2.3.3 节的 1.57 系数），并注意\textbf{阻感负载时电流波形非正弦}，不能用简单的正弦公式。

\subsection{三相交流调压电路}

三相交流调压电路按其联结形式有多种，各有不同的适用场合与限制。

\begin{table}[htbp]
\centering
\caption{三相交流调压电路的联结形式}
\label{tab:6-1}
\begin{tabular}{p{3.6cm}p{4.6cm}p{6.2cm}}
\toprule
\textbf{联结形式} & \textbf{构成与特点} & \textbf{使用注意} \\
\midrule
星形联结带中性线 & 三个单相调压电路分别接在三相四线制上 & 三次谐波电流能经中性线流通，中性线电流大；只能用于小容量 \\
星形联结不带中性线 & 三只双向晶闸管（或三对反并联晶闸管）接成星形，负载星形接 & 要求 $\alpha\ge\varphi$；\textbf{须用宽脉冲或双窄脉冲}，否则某相可能因无触发而不能形成回路；三次谐波无通路 \\
三角形联结 & 三只双向晶闸管接在负载三角形内部 & 线电流中三次谐波不能流通，但相电流中有 \\
支路控制三角形联结 & 晶闸管控制三角形负载的支路电流（TCR，晶闸管控制电抗器） & 用于静止无功补偿，连续调节感性无功 \\
中点控制三角形联结 & 三相负载接成三角形，用三只晶闸管控制其中点 & 用得较少 \\
\bottomrule
\end{tabular}
\end{table}

\corepoint{三相交流调压最常考“星形不带中性线”：由于没有中性线，\textbf{三次及三的倍数次谐波电流没有通路}，故线电流中不含 $3k$ 次谐波（只含 $6k\pm1$ 次）。另外\textbf{窄脉冲不行}——因为每相电流要靠另两相构成回路，脉冲必须足够宽（或用双窄脉冲）才能保证成对导通。}

% ============================================================
\section{其他交流电力控制电路}

\subsection{交流调功电路}

交流调功电路与交流调压电路\textbf{电路形式完全相同}，区别只在控制方式：它不是在每个周期内控制导通相位，而是\textbf{以交流电的周期为单位}控制通断——接通几个周期、再断开几个周期，通过改变通断周期数的比值来调节输出功率。这种控制方式也称\textbf{周波数控制}或“通断控制”。

设控制周期为 $N$ 个电源周期，其中接通 $n$ 个周期，则输出功率

\keyformula{P=\frac{n}{N}P_{\max},\qquad P_{\max}=\frac{U_1^2}{R}}

\begin{itemize}
  \item 优点：晶闸管在电压过零时通断，\textbf{不产生相位控制的谐波}，对电网的干扰小；
  \item 缺点：输出功率是阶梯式调节，且通断周期较长时会引起负载温度脉动（对温度惯性大的负载如电炉、温控系统正合适）；
  \item 适用场合：电炉温度控制、空调压缩机、电热水器等\textbf{惯性大、对谐波敏感}的负载。
\end{itemize}

\subsection{交流电力电子开关}

把晶闸管（或双向晶闸管）反并联接在负载回路中，只作“通/断”用而不调压，就是\textbf{交流电力电子开关}（固态开关，SSR）。

\begin{itemize}
  \item 与机械式接触器相比：\textbf{无触点、无电弧、动作快、寿命长、无噪声}，可实现频繁通断；
  \item 与交流调压相比：只需判断“通”或“断”，控制简单，通态时晶闸管全导通、损耗小；
  \item 常与补偿电容配合构成\textbf{晶闸管投切电容器（TSC）}，用于无功补偿——在电容电压与电网电压相等的时刻投入，可避免冲击电流。
\end{itemize}

% ============================================================
\section{交—交变频电路}

\textbf{交—交变频电路}（周波变流器，Cycloconverter）不经过中间直流环节，直接把一种频率的交流电变成另一种频率的交流电，因此也称“直接变频”。

\subsection{单相交—交变频电路}

\subsubsection*{1. 电路构成和基本工作原理}

单相交交变频电路由\textbf{两组反并联的相控整流电路}（正组 P 与反组 N）构成：正组工作时输出正半周电流，反组工作时输出负半周电流。改变两组整流的控制角 $\alpha$，就可使输出电压平均值按正弦规律变化。

\begin{itemize}
  \item 输出电压的幅值由 $\alpha$ 决定（$\alpha$ 在 $0\sim\pi/2$ 之间变化时 $U_d$ 由正的最大值变到零）；
  \item 输出电压的频率由 $\alpha$ 变化的速率决定——$\alpha$ 按正弦规律从 $\pi/2$ 减小到某个最小值再回到 $\pi/2$，输出就完成一个正弦半周；
  \item 这种控制方法称为\textbf{余弦交点法}：把期望输出的正弦波与各段线电压的余弦曲线求交，交点时刻即为触发时刻，从而使输出平均电压逼近正弦。
\end{itemize}

\subsubsection*{2. 输出频率的上限}

由于交交变频靠“拼”电源电压的片段来合成正弦波，一个输出周期内必须包含足够多的电源电压片段才能保证波形质量，因此\textbf{输出频率不能高于输入频率}：

\keyformula{f_o\le\left(\frac{1}{3}\sim\frac{1}{2}\right)f_i}

即电网 50 Hz 时，交交变频的输出频率上限约为 20 Hz 左右——这使它特别适合\textbf{低速大功率}传动（如轧机、矿井提升机、水泥磨机）。

\subsubsection*{3. 无环流与有环流控制}

正、反两组反并联，如果两组同时导通，会在两组之间产生\textbf{环流}（不经过负载的短路电流）。处理方式有两种：

\begin{itemize}
  \item \textbf{无环流控制}：任何时刻只允许一组工作，另一组封锁脉冲；切换时须先封锁、等电流过零并留死区，再开放另一组，因此\textbf{切换有死区、输出波形有畸变}，但不需要环流电抗器，效率高——实际多用这种。
  \item \textbf{有环流控制}：两组同时加触发脉冲，用\textbf{环流电抗器}限制环流；切换平滑、无死区，但需笨重的电抗器、损耗大。
\end{itemize}

\subsection{三相交—交变频电路}

三相交交变频电路由三组（每相一组）单相交交变频电路构成，按输出绕组连接方式分为\textbf{公共交流母线进线方式}与\textbf{输出星形联结方式}。

\begin{itemize}
  \item \textbf{公共交流母线进线方式}：三组单相变频电路的输入端接在同一交流母线上，输出端分别接电动机的三相绕组（绕组须分开、各有 6 个端子），三组相互独立，控制简单。
  \item \textbf{输出星形联结方式}：三组输出接成星形，电动机三相绕组的一端连成中性点。由于中性点不引出，\textbf{三组的输出电流之和为零}，可实现三相电流的自动平衡；但每组输出端对地电压较高，绝缘要求高。
\end{itemize}

三相交交变频的特点：\textbf{输出频率低（$\le f_i/3$）、功率大、效率高、无需中间直流环节}，但\textbf{功率因数低}（因相控整流在深控时 $\cos\varphi$ 很小）、\textbf{输入电流谐波大}、输出波形含较多谐波。

% ============================================================
\section{矩阵式变频电路}

\textbf{矩阵式变频电路}（Matrix Converter）把三相交流输入直接变成三相交流输出，中间不经任何储能元件，只用 9 个\textbf{双向开关}（每个由两只反并联的全控器件构成）按一定规律通断，把输入的三相电压“切片”后拼成输出三相电压。

\begin{itemize}
  \item \textbf{理论上可输出任意频率}（不受 $f_i/3$ 的限制），且输入电流可控制为正弦、功率因数可调（可做到 $\lambda\approx1$）；
  \item 无中间储能元件（无大电容、大电感），\textbf{功率密度高、体积小、寿命长}；
  \item 约束条件：（1） 任一输出相不能与两个输入相同时接通（否则输入线电压短路）；（2） 任一输入相不能同时断开所有输出（否则感性负载电流无通路）；这就是“\textbf{双向开关的安全换流}要求”，实际须采用四步换流或两步换流逻辑；
  \item 缺点：器件数量多（9 个双向开关 = 18 个全控器件）、驱动与保护复杂、电压传输比有限（理论上最大约 0.866），故目前主要用于中小功率、要求高功率密度的场合。
\end{itemize}

\corepoint{本章复习抓“三条线”：\textbf{（1） 交流调压}——电阻负载 $U_o=U_1\sqrt{\frac{\pi-\alpha}{\pi}+\frac{\sin2\alpha}{2\pi}}$、$\lambda=U_o/U_1$，阻感负载必须 $\alpha\ge\varphi$ 否则失控；\textbf{（2） 交流调功}——电路同调压但按周期通断，$P=(n/N)P_{\max}$，谐波小；\textbf{（3） 交交变频}——两组反并联相控整流 + 余弦交点法，输出频率 $\le f_i/3$，适合低速大功率，功率因数低。矩阵式变频则是“无储能直接变换”的方向性电路。}

\chapter{PWM 控制技术}

PWM（Pulse Width Modulation，脉冲宽度调制）技术是电力电子技术中最重要的控制技术之一。它没有引入新的电路拓扑，却使电路的控制性能大为改善，凡是用全控型器件的电路（逆变、整流、斩波、交交变频）都可以用它。本章讲四件事：PWM 的理论依据、PWM 逆变电路及其控制方法、跟踪型 PWM、PWM 整流电路。

% ============================================================
\section{PWM 控制的基本原理}

PWM 控制的理论基础是\textbf{面积等效原理}（冲量等效原理）：

\begin{tcolorbox}[enhanced,breakable,colback=highlight,colframe=yhbAccent!30!white,boxrule=0.5pt,arc=2.5pt,left=11pt,right=11pt,top=9pt,bottom=9pt,boxsep=0pt,before skip=11pt,after skip=11pt]
\textbf{面积等效原理}\quad 冲量相等而形状不同的窄脉冲加在具有惯性的环节上时，其效果基本相同。“冲量”指窄脉冲的面积，“效果基本相同”指环节的输出响应波形基本相同。
\end{tcolorbox}

由此可以推得：把正弦半波分成 $N$ 等份，用 $N$ 个\textbf{等幅不等宽}的矩形脉冲代替，使每个矩形脉冲的\textbf{面积与对应正弦波段的面积相等}，则这串脉冲加到惯性环节上，其效果与正弦半波基本相同。这就是 \textbf{SPWM}（正弦波脉宽调制）的由来。

\begin{itemize}
  \item 脉冲的\textbf{宽度}按正弦规律变化、\textbf{幅值}恒定——称为 PWM 波形；
  \item 同样的道理，也可以把非正弦的期望波形（如梯形波、指定谐波消除波形）用等幅不等宽脉冲等效；
  \item 用 PWM 波形代替正弦波，关键在于\textbf{开关频率足够高}（远高于输出基波频率）时，负载（电动机等惯性环节）只“感受到”平均效果。
\end{itemize}

\corepoint{PWM 控制的三个基本量必须分清：\textbf{调制波}（期望输出的参考波，通常是正弦波）、\textbf{载波}（高频三角波或锯齿波，决定开关频率）、\textbf{调制比} $m=U_{rm}/U_{cm}$（决定输出基波幅值）、\textbf{载波比} $N=f_c/f_r$（决定每周期内的脉冲个数与谐波分布）。PWM 的输出电压基波幅值与 $m$ 成正比，这是它“调压”的原理；而开关频率决定了谐波往高频搬移，这是它“改善波形”的原理。}

% ============================================================
\section{PWM 逆变电路及其控制方法}

\subsection{计算法和调制法}

\textbf{计算法}：根据预期的输出电压波形，按面积等效原理计算出各脉冲的宽度与间隔，用数字方式（查表或实时计算）生成 PWM 波形。计算法复杂、不灵活，但可以\textbf{精确控制}波形（如指定谐波消除），多用于离线计算、查表运行的场合。

\textbf{调制法}：把期望的输出波形作为\textbf{调制信号}，把接受调制的信号作为\textbf{载波}，通过比较得到 PWM 波形。这是工程上应用最广的方法。

\begin{itemize}
  \item \textbf{单极性调制}：载波为单极性三角波（只在正或负半周变化），调制波为正弦，比较结果为\textbf{单极性}脉冲列（$0$ 或 $+U_d$，或 $0$ 或 $-U_d$），输出谐波少、但需两只管子按半周交替工作。
  \item \textbf{双极性调制}：载波为双极性三角波（在正负之间变化），比较结果为\textbf{双极性}脉冲列（$+U_d$ 或 $-U_d$），控制简单、同一桥臂上下管互补，是三相桥 SPWM 最常用的方式。
\end{itemize}

单相全桥双极性 SPWM 的输出电压基波幅值为

\keyformula{U_{o1m}=mU_d\qquad(0\le m\le1)}

三相桥式 SPWM 逆变电路（三个桥臂共用同一三角载波，三相调制波互差 $120^\circ$）输出相电压的基波幅值为

\keyformula{U_{UN1m}=\frac{mU_d}{2}\qquad(0\le m\le1)}

\subsection{异步调制和同步调制}

按载波比 $N=f_c/f_r$ 是否固定，分为两种调制方式：

\begin{table}[htbp]
\centering
\caption{异步调制与同步调制}
\label{tab:7-1}
\begin{tabular}{p{2.4cm}p{5.4cm}p{6.2cm}}
\toprule
\textbf{方式} & \textbf{做法} & \textbf{特点} \\
\midrule
异步调制 & 载波频率 $f_c$ 固定，$N$ 随 $f_r$ 变化 & 低频时 $N$ 很大、脉冲多、波形好；高频时 $N$ 变小、正负半周脉冲不对称，产生偶次谐波与波形畸变 \\
同步调制 & $N$ 固定（$f_c$ 随 $f_r$ 成正比变化） & 正负半周严格对称、无偶次谐波，但低频时 $N$ 太小、谐波次数低 \\
分段同步调制 & 把 $f_r$ 分成若干段，每段内取不同的固定 $N$ & 兼顾两者：低频段用大 $N$、高频段用小 $N$，是实际变频器常用方案 \\
\bottomrule
\end{tabular}
\end{table}

\subsection{规则采样法}

自然采样法的脉宽由正弦波与三角波的自然交点决定，需解超越方程，难以实时计算。\textbf{规则采样法}用固定的采样点近似：

\begin{itemize}
  \item \textbf{对称规则采样}：在每个载波周期内的固定时刻（如三角波顶点）对调制波采样，以该值决定脉宽，脉宽关于载波周期中点对称；
  \item \textbf{不对称规则采样}：在三角波的顶点和底点各采样一次，脉宽不对称，精度更高。
\end{itemize}

对称规则采样法的脉宽公式为

\keyformula{\delta=\frac{T_c}{2}\left(1+m\sin\omega_r t_s\right)}

式中 $T_c$ 为载波周期，$t_s$ 为采样时刻。规则采样法计算简单、便于数字实现，其输出与自然采样法的差别在开关频率足够高时很小。

\subsection{PWM 逆变电路的谐波分析}

PWM 逆变电路输出电压的谐波分布与载波比 $N$ 密切相关。以三相桥双极性 SPWM 为例（$m$ 一定、$N$ 为整数）：

\begin{itemize}
  \item 输出线电压中\textbf{不含低次谐波}，最低次谐波出现在\textbf{载波频率 $f_c$ 附近}（即 $N$ 次及其边频 $N\pm2$、$N\pm4$ 等）；
  \item 谐波幅值随 $m$ 增大而增大，但\textbf{谐波次数不随 $m$ 改变}；
  \item $N$ 越大，最低次谐波次数越高、越容易被滤除，故提高载波频率是改善波形的主要手段（代价是开关损耗增大）。
\end{itemize}

\subsection{提高直流电压利用率和减少开关次数}

SPWM 的\textbf{直流电压利用率低}是它的主要缺点：$m\le1$ 时三相桥输出线电压基波幅值最大为 $\dfrac{\sqrt{3}}{2}U_d=0.866U_d$，即\textbf{利用率只有 0.866}。提高的办法有：

\begin{enumerate}
  \item \textbf{注入三次谐波}：在正弦调制波中注入幅值为基波 $1/6$ 的三次谐波（三线制中三次谐波不产生线电流），可使调制波峰值降低而不改变基波，从而把直流电压利用率\textbf{提高约 15\%}；
  \item \textbf{梯形波调制}：用梯形波作调制波，利用率可提高到 1 左右，但输出含较多低次谐波；
  \item \textbf{空间矢量 PWM（SVPWM）}：从电动机磁链圆的角度出发选择电压矢量，直流电压利用率可达 \textbf{1.0}（比 SPWM 高约 15.5\%），且便于数字实现，是目前变频器的主流算法。
\end{enumerate}

\keyformula{\text{SPWM: }\frac{U_{UV1m}}{U_d}\le0.866;\qquad \text{SVPWM: }\frac{U_{UV1m}}{U_d}\le1.0}

\subsection{空间矢量 PWM 控制}

三相桥式逆变电路的 6 个开关共有 8 种合法开关状态，对应 \textbf{8 个基本电压矢量}：6 个幅值为 $\frac{2}{3}U_d$ 的非零矢量（$U_1\sim U_6$，互差 $60^\circ$，构成正六边形）与 2 个零矢量（$U_0$、$U_7$）。

\begin{itemize}
  \item 期望的电压矢量落在某个 $60^\circ$ 扇区内时，可用该扇区相邻的两个非零矢量与零矢量\textbf{在一个载波周期内分时合成}（矢量合成即“伏秒平衡”）；
  \item 设两相邻矢量作用时间为 $T_1$、$T_2$，零矢量作用时间为 $T_0$，则 $T_1+T_2+T_0=T_c$，由伏秒平衡可解出 $T_1$、$T_2$；
  \item 优点：直流电压利用率高、开关次数少（每次只切换一个桥臂）、便于微机实时计算。
\end{itemize}

\subsection{PWM 逆变电路的多重化}

把几组 PWM 逆变电路按一定相位差组合（多重化），可使载波频率附近的谐波相互抵消、等效开关频率成倍提高，从而在\textbf{不提高器件开关频率}的前提下改善波形——这与第 4 章的多重化思想一致，只是这里抵消的是 PWM 的高频谐波。

% ============================================================
\section{PWM 跟踪控制技术}

\textbf{跟踪控制}不预先计算脉冲宽度，而是把期望输出（电流或电压）作为指令，\textbf{实时比较实际输出与指令的偏差}，由偏差直接决定开关动作，使实际输出跟踪指令变化。它属于“闭环 PWM”。

\subsection{滞环比较方式}

把指令信号与实际信号送入\textbf{滞环比较器}（环宽为 $\pm\Delta I$），当偏差超出环宽时改变开关状态。以电流跟踪型 PWM 为例：

\begin{itemize}
  \item 实际电流上升超过上限 $I^*+\Delta I$ 时，关断使电流下降的开关组合；低于下限 $I^*-\Delta I$ 时，开通使电流上升的组合；
  \item \textbf{环宽越小}，跟踪误差越小，但开关频率越高、开关损耗越大；\textbf{环宽越大}，开关频率低而误差大；
  \item 优点：硬件简单、响应快、不需要载波，天然限流；缺点：\textbf{开关频率不固定}（随电流变化率而变），给滤波设计带来困难，且输出谐波频谱较宽。
\end{itemize}

\subsection{三角波比较方式}

把指令电流与实际电流的偏差经 PI 调节后，与\textbf{三角波载波}比较产生 PWM 信号。

\begin{itemize}
  \item \textbf{开关频率固定}（等于载波频率），便于滤波器设计，输出谐波集中在载波频率附近；
  \item 由于有 PI 调节器，稳态误差小、跟踪精度高；
  \item 缺点是响应速度略慢于滞环方式，且 PI 参数需整定。
\end{itemize}

% ============================================================
\section{PWM 整流电路及其控制方法}

\subsection{PWM 整流电路的工作原理}

PWM 整流电路与 PWM 逆变电路结构相同（电压型／电流型、单相／三相），只是\textbf{能量流向相反}：它把交流电变成直流电，但通过 PWM 控制可使\textbf{交流侧电流接近正弦}且\textbf{功率因数可调（可做到单位功率因数，甚至电流超前）}，还能实现\textbf{能量的双向流动}（可整流也可逆变）——这是相控整流做不到的。

\begin{itemize}
  \item \textbf{电压型 PWM 整流电路}：直流侧为大电容，交流侧经电感接入电网，桥臂由全控器件与反并联二极管构成；
  \item \textbf{电流型 PWM 整流电路}：直流侧为大电感，交流侧需电容滤波。
\end{itemize}

工作原理要点：通过控制桥臂开关，使桥臂中点电压（PWM 电压）中含有与电网电压同频同相的基波分量，控制该基波的\textbf{幅值与相位}，即可控制交流侧电流的大小与相位——本质上是用 PWM 电压“顶”着电网电压，靠电感上的电压差控制电流。

\subsection{PWM 整流电路的控制方法}

\begin{enumerate}
  \item \textbf{间接电流控制}（幅相控制）：直接控制 PWM 电压的幅值与相位，间接控制交流电流。结构简单、无需电流传感器，但\textbf{动态响应慢}、对参数敏感，电流波形受电网电压畸变影响大。
  \item \textbf{直接电流控制}：引入交流电流闭环，把电流指令与实际电流的偏差经调节器后生成 PWM 指令。又分为：
    \begin{itemize}
      \item \textbf{滞环电流控制}：环宽决定开关频率，响应快；
      \item \textbf{固定开关频率的电流控制}（如三角波比较、同步旋转坐标系 PI 控制）：开关频率固定、电流正弦度好，工业上最常用；
      \item \textbf{预测电流控制}：按下一拍期望电流反算所需电压，动态最快。
    \end{itemize}
  \item 工程上电压型 PWM 整流通常采用\textbf{电压外环 + 电流内环}的双闭环结构：电压外环稳定直流母线电压并给出有功电流指令，电流内环实现快速跟踪；无功电流指令置零即得单位功率因数。
\end{enumerate}

\corepoint{本章必须记牢的数字与概念：\textbf{调制比 $m=U_{rm}/U_{cm}$、载波比 $N=f_c/f_r$}；\textbf{单相全桥 $U_{o1m}=mU_d$、三相桥相电压 $U_{UN1m}=mU_d/2$}；\textbf{SPWM 直流电压利用率 0.866、SVPWM 为 1.0（提高 15.5\%）、三次谐波注入提高 15\%}；\textbf{异步调制低频好高频差、同步调制对称但低频谐波大、分段同步折中}；\textbf{滞环比较开关频率不固定、三角波比较频率固定}；\textbf{PWM 整流能实现单位功率因数与能量双向流动}。}

\chapter{软开关技术}

现代电力电子装置的发展方向是高频化、小型化、高功率密度。但开关频率越高，\textbf{开关损耗}越大，器件发热、效率下降、电磁干扰加剧——这个矛盾靠“硬开关”本身无法解决。\textbf{软开关技术}就是在开关过程前后引入谐振（或辅助电路），使开关在\textbf{电压为零}或\textbf{电流为零}的时刻动作，从而把开关损耗降到接近零。

% ============================================================
\section{软开关的基本概念}

\subsection{硬开关与软开关}

前面各章的分析都把开关理想化，认为器件瞬间完成通断。实际上器件的开通和关断都需要时间，在这段时间内\textbf{电压和电流同时不为零}，其乘积即为损耗。

\begin{itemize}
  \item \textbf{开通损耗}：开通时电流上升、电压下降，二者交叠产生损耗；若器件两端并联电容，开通瞬间电容储能经器件释放，还会产生很大的\textbf{电流尖峰}（容性开通）。
  \item \textbf{关断损耗}：关断时电流下降、电压上升，二者交叠产生损耗；若回路中有电感，关断瞬间电感储能产生很高的\textbf{电压尖峰}（感性关断），可能击穿器件。
  \item 此外，硬开关还伴随严峻的\textbf{电磁干扰}（du/dt、di/dt 很大）与噪声。
\end{itemize}

\begin{tcolorbox}[enhanced,breakable,colback=highlight,colframe=yhbAccent!30!white,boxrule=0.5pt,arc=2.5pt,left=11pt,right=11pt,top=9pt,bottom=9pt,boxsep=0pt,before skip=11pt,after skip=11pt]
\textbf{硬开关}\quad 开关过程中电压、电流均不为零，出现了\textbf{交叠}，因而有开关损耗。\\
\textbf{软开关}\quad 通过谐振等手段，使开关\textbf{开通前电压先降到零}，或\textbf{关断前电流先降到零}，从而消除电压电流的交叠，使开关损耗为零（理论上）。软开关同时还能抑制 du/dt、di/dt，降低电磁干扰。
\end{tcolorbox}

\subsection{零电压开关与零电流开关}

\begin{itemize}
  \item \textbf{零电压开关（ZVS）}：使开关开通前其两端电压为零，则开通时不会产生损耗（也就没有容性开通电流尖峰）。要实现 ZVS，须让开关两端的电压在开通前先谐振到零，并让反并联二极管先导通（把电压钳在零）再给开关开通信号。
  \item \textbf{零电流开关（ZCS）}：使开关关断前其电流为零，则关断时不会产生损耗（也就没有感性关断电压尖峰）。要实现 ZCS，须让流过开关的电流在关断前先谐振到零，并让反并联二极管先导通。
\end{itemize}

\corepoint{记住对应关系：\textbf{ZVS 对应“开通”（电压为零）}、\textbf{ZCS 对应“关断”（电流为零）}；实际电路中往往同时安排“零电压开通 + 零电流关断”（如移相全桥），即所谓“全软开关”。判断某电路属于哪种软开关，就看\textbf{谐振元件（$L_r$、$C_r$）与开关的连接方式}以及\textbf{开关动作时刻电压电流是否为零}。}

% ============================================================
\section{软开关电路的分类}

按实现软开关所用电抗元件与开关的连接方式，软开关电路分为三大类。

\subsection{准谐振电路}

在开关两端（或串联）接入 $L_r$、$C_r$ 构成谐振回路，使电压或电流在开关动作时刻形成\textbf{正弦半波}（故称“准谐振”）。分为：

\begin{itemize}
  \item \textbf{零电压开关准谐振电路（ZVS-QR）}：$L_r$、$C_r$ 串联后并联在开关两端，开关开通前 $C_r$ 上的电压先谐振到零；
  \item \textbf{零电流开关准谐振电路（ZCS-QR）}：$L_r$、$C_r$ 与开关串联，开关关断前 $L_r$ 中的电流先谐振到零；
  \item \textbf{多谐振电路（MRC）}：电路中存在多个谐振元件，可同时实现开关与二极管的软开关。
\end{itemize}

谐振频率与特性阻抗决定了谐振过程的快慢与峰值：

\keyformula{\omega_r=\frac{1}{\sqrt{L_rC_r}},\qquad Z_r=\sqrt{\frac{L_r}{C_r}}}

准谐振电路的共同缺点是：\textbf{开关应力大}（谐振电压/电流峰值高）、\textbf{输出与开关频率相关}（只能靠调频调压，不利于滤波设计），且软开关条件与负载有关（轻载时可能失去软开关）。

\subsection{零开关 PWM 电路}

在准谐振电路基础上加辅助开关，使电路\textbf{只在开关动作前后的短时间内谐振}，其余时间按 PWM 方式工作。分为\textbf{零电压开关 PWM 电路（ZVS-PWM）}与\textbf{零电流开关 PWM 电路（ZCS-PWM）}。

优点：既实现了软开关，又保持了 PWM 控制（恒频、调宽）的优点，输出电压与占空比成线性关系。

\subsection{零转换 PWM 电路}

谐振网络\textbf{与主开关并联}（而非串联），且只在开关动作时投入工作，主电路电流不流过谐振网络，因此\textbf{导通损耗小、软开关范围宽（轻载也能实现）}。分为\textbf{零电压转换 PWM 电路（ZVT-PWM）}与\textbf{零电流转换 PWM 电路（ZCT-PWM）}。

\begin{table}[htbp]
\centering
\caption{三类软开关电路的比较}
\label{tab:8-1}
\begin{tabular}{p{3.2cm}p{4.2cm}p{6.6cm}}
\toprule
\textbf{类型} & \textbf{谐振元件位置} & \textbf{特点} \\
\midrule
准谐振（QR） & 与开关串联或并联 & 谐振贯穿整个开关周期，应力大、调频控制、软开关范围受负载影响 \\
零开关 PWM（ZS-PWM） & 与开关串联或并联，加辅助开关 & 只在开关瞬间谐振，恒频 PWM，应力仍较大 \\
零转换 PWM（ZT-PWM） & 与主开关并联，辅助支路 & 主电路电流不经过谐振网络，损耗小、软开关范围宽，是主流方案 \\
\bottomrule
\end{tabular}
\end{table}

% ============================================================
\section{典型的软开关电路}

\subsection{零电压开关准谐振电路}

这是结构较为简单的软开关电路，以 Buck 型为例：在开关 $S$ 两端并联谐振电容 $C_r$，在主回路串入谐振电感 $L_r$。

\begin{itemize}
  \item $S$ 关断时，$L_r$ 中的电流向 $C_r$ 充电，$S$ 两端电压 $u_S$ 从零开始上升（\textbf{零电压关断}，无容性关断损耗）；随后 $L_r$、$C_r$ 谐振，$u_S$ 形成正弦半波；
  \item 当 $u_S$ 谐振回到零时，反并联二极管先导通把 $u_S$ 钳在零，此时再给 $S$ 开通信号，即实现\textbf{零电压开通}；
  \item 谐振电压峰值约为 $2E$，因此\textbf{开关管耐压须按 $2E$ 以上选取}，这是 ZVS-QR 的主要缺点；此外输出电压与开关频率成反比（调频调压），轻载时可能失去 ZVS。
\end{itemize}

\subsection{谐振直流环}

谐振直流环（Resonant DC Link，RDCL）适用于变频器：在直流母线与逆变桥之间串入小电感 $L_r$，在逆变桥的直流侧并联电容 $C_r$，使\textbf{直流母线电压周期性谐振到零}，从而为逆变桥的所有开关提供零电压开通条件。

\begin{itemize}
  \item 母线电压每谐振到零一次，逆变桥就可以“择机”进行一次开关动作（离散脉冲调制），从而实现全桥 ZVS；
  \item 优点：所有开关都能零电压开通，控制灵活；
  \item 缺点：母线电压峰值高（约 $2E$ 以上）、器件电压应力大，且谐振过程增加了额外的导通损耗与 EMI。
\end{itemize}

\subsection{移相全桥型零电压开关 PWM 电路}

这是目前中大功率 DC-DC 变换器最常用的软开关拓扑：全桥的四个开关中，\textbf{超前桥臂}与\textbf{滞后桥臂}的驱动信号之间加入\textbf{移相角}，利用变压器\textbf{漏感}（或外加谐振电感）与开关管的\textbf{结电容}谐振实现 ZVS。

\begin{itemize}
  \item 每个桥臂的上下管互补导通（各约 50\%），两桥臂之间移相，移相角决定输出电压（占空比）；
  \item \textbf{死区时间内}，原导通的管子关断后，负载电流经另一管的结电容与漏感谐振，把即将开通的管子两端电压谐振到零，再由其反并联二极管钳位，从而实现 ZVS；
  \item 特点：\textbf{超前桥臂容易实现 ZVS}（利用滤波电感能量），\textbf{滞后桥臂较难}（只靠漏感能量，轻载时可能失去 ZVS）；
  \item 存在\textbf{占空比丢失}：换相期间二次侧电压被“短路”，有效占空比小于一次侧给定占空比，输出电压比理论值低。
\end{itemize}

\subsection{零电压转换 PWM 电路}

零电压转换（ZVT-PWM）电路在主开关两端并联一个由辅助开关 $S_a$、谐振电感 $L_r$、谐振电容 $C_r$ 组成的辅助支路：

\begin{itemize}
  \item 主开关开通前，先开通辅助开关 $S_a$，$L_r$ 中的电流线性上升，把主开关的电流“抽走”，使主开关反并联二极管导通、电压钳在零；
  \item 然后开通主开关，即为\textbf{零电压开通}；随后关断 $S_a$，$L_r$ 的能量经二极管回馈，辅助开关也可实现软关断；
  \item 优点：\textbf{辅助支路只在开关瞬间工作}，主电路电流不经过 $L_r$，导通损耗小；\textbf{软开关范围宽}，轻载仍能实现 ZVS；辅助器件容量小、成本低。
\end{itemize}

% ============================================================
\section{软开关技术新进展}

软开关技术的发展方向主要有：

\begin{enumerate}
  \item \textbf{无源无损缓冲电路}：只用电感、电容、二极管构成缓冲网络，不加辅助开关，靠谐振把开关损耗转移到缓冲元件中再回馈，结构简单、可靠性高，但软开关范围有限。
  \item \textbf{有源钳位电路}：用一只辅助开关加钳位电容，把开关两端的电压尖峰钳位并回收能量，同时实现主开关的 ZVS，广泛用于正激、反激等单端电路。
  \item \textbf{LLC 谐振变换器}：用谐振网络（$L_r$、$C_r$ 与励磁电感 $L_m$）实现全负载范围的 ZVS 与二次侧 ZCS，效率高、EMI 小，是目前服务器电源、充电器的热门拓扑。
  \item \textbf{复合软开关与多电平软开关}：把软开关技术与多电平、多重化技术结合，用于高压大容量场合；也采用“几个简单高效的开关电路级联、并联、串并联组合”的思路，用成熟单元构成高性能系统。
  \item \textbf{软开关与数字控制结合}：用数字信号处理器实时调整死区与移相角，保证全工况下的软开关条件（自适应死区控制）。
\end{enumerate}

\corepoint{本章的考点是“概念辨析 + 分类归属”：\textbf{（1） 硬开关损耗来自电压电流交叠，软开关靠谐振消除交叠}；\textbf{（2） ZVS 对应开通、ZCS 对应关断}；\textbf{（3） 三类电路的归属}——准谐振（谐振贯穿全周期、应力大、调频）→ 零开关 PWM（只瞬间谐振、恒频、应力仍大）→ 零转换 PWM（谐振网络与主开关并联、损耗小、范围宽）；\textbf{（4） 典型电路的软开关机理}：ZVS-QR 靠 $L_rC_r$ 谐振到零、谐振直流环靠母线电压周期性到零、移相全桥靠漏感与结电容谐振（滞后桥臂轻载难 ZVS、有占空比丢失）、ZVT 靠辅助支路先抽走电流。}

\chapter{电力电子器件应用的共性问题}

第 2 章讲了器件本身，本章讲器件\textbf{用到电路里}时的共性问题：怎么驱动它开通与关断、怎么保护它不被过电压过电流损坏、电压或电流不够时怎么串并联使用。这三个问题是所有电力电子装置都要面对的工程问题，与具体电路拓扑无关。

% ============================================================
\section{电力电子器件的驱动}

\subsection{电力电子器件驱动电路概述}

\textbf{驱动电路}是主电路与控制电路之间的接口，其作用是：把信息电子电路发出的控制信号\textbf{放大}到足以使器件可靠开通或关断的功率水平，并实现主电路与控制电路之间的\textbf{电气隔离}。

\begin{itemize}
  \item \textbf{隔离方式}：主电路为高压、控制电路为低压，必须隔离。常用\textbf{光耦合器}（普通型、高速型、高传输比型）与\textbf{脉冲变压器}（磁隔离）。光耦的响应时间可做到 $1.5\,\mu$s 以下；脉冲变压器在脉冲较宽时须采用\textbf{高频调制与解调}，以避免铁心饱和。
  \item \textbf{驱动电路的性能直接决定器件的开关特性}：驱动电流/电压的幅值、前沿陡度、上升下降时间，影响开通关断损耗与开关时间；驱动电路还应具备\textbf{保护功能}（如过流时软关断、欠压时封锁）。
  \item 工程上应优先选用器件生产厂家配套的\textbf{专用驱动电路/驱动模块}，以获得参数的最佳配合。
\end{itemize}

驱动电路按控制信号性质分为\textbf{电流驱动型}（晶闸管、GTO、GTR）与\textbf{电压驱动型}（电力 MOSFET、IGBT）两大类。

\subsection{晶闸管的触发电路}

晶闸管的驱动电路常称为\textbf{触发电路}，其作用是产生符合要求的门极触发脉冲，保证晶闸管在需要的时刻由阻断转为导通。对触发电路的要求（也是第 3 章 3.8 节的要求）：

\begin{enumerate}
  \item 触发脉冲应具有足够的\textbf{幅值和功率}（一般门极电压 $4\sim10$ V、电流几百毫安到几安）；
  \item 脉冲\textbf{前沿要陡}（上升时间短），以提高触发时刻的准确性、避免局部过热；
  \item 脉冲\textbf{宽度足够}（阻感负载需宽脉冲或脉冲列），保证阳极电流上升到擎住电流以上；
  \item 具有足够的\textbf{移相范围}，并与主电路电源\textbf{同步}；
  \item 各相脉冲应\textbf{对称}（间隔均匀），否则会造成输出电压不对称、产生直流分量。
\end{enumerate}

常用实现方式：阻容移相桥（简单但精度低）、单结晶体管触发电路（简单、用于小功率）、\textbf{锯齿波同步触发电路}（应用最广）、集成触发器（KJ004 系列）、数字触发器（精度可达 $0.7^\circ$ 以内）。

\subsection{典型全控型器件的驱动电路}

\subsubsection*{1. 电流驱动型器件的驱动电路}

\textbf{GTO 的驱动}：GTO 需要\textbf{正、负两个电源}——开通时加正向门极电流（幅值较小、持续整个导通期），关断时加\textbf{幅值很大的负门极电流脉冲}（因电流关断增益 $\beta_{off}$ 只有 5 左右，关断电流 $I_{GM}\approx I_{ATO}/5$）。因此 GTO 驱动电路复杂、功耗大，须用强脉冲加前沿陡化措施。

\textbf{GTR 的驱动}：基极驱动电流应使 GTR 工作在\textbf{准饱和状态}（防止过饱和导致关断时间变长），常用\textbf{抗饱和（贝克箝位）电路}；关断时须加\textbf{反向基极电流}以加速存储电荷抽出。GTR 驱动电路还需具备过流保护（退饱和检测）。

\subsubsection*{2. 电压驱动型器件的驱动电路}

\textbf{电力 MOSFET 的驱动}：本质是给栅极电容充放电。要求：
\begin{itemize}
  \item 开通时加足够幅值的正栅压（一般 $10\sim15$ V），保证充分导通、$R_{on}$ 最小；
  \item 关断时加\textbf{负栅压}（$-5\sim-15$ V），防止 du/dt 引起的误导通；
  \item 串接\textbf{栅极电阻} $R_g$ 控制充放电速度：$R_g$ 小则开关快、损耗小但 EMI 大；$R_g$ 大则相反；
  \item 驱动电路应能提供足够的峰值电流（$\pm$ 数安），并具有隔离。
\end{itemize}

\textbf{IGBT 的驱动}：与 MOSFET 类似但要求更高——栅极电荷大、有\textbf{米勒效应}（关断时集电极电压上升通过米勒电容耦合到栅极，可能引起误导通）。典型要求：正栅压 $+15$ V、负栅压 $-5\sim-15$ V；具备\textbf{退饱和（$U_{CE(sat)}$）检测}以实现短路保护，并采用\textbf{软关断}（缓慢降低栅压）限制关断过电压。常用集成驱动模块（如 EXB841、M57962 等）把这些功能集成在一起。

\corepoint{驱动电路的三句话：\textbf{（1） 晶闸管要“脉冲”（够宽够陡），GTO/GTR 要“正负双电源、强脉冲”（电流驱动），MOSFET/IGBT 要“正负栅压 + 栅极电阻”（电压驱动）}；\textbf{（2） 隔离方式看场合：光耦最常用，脉冲变压器需高频调制}；\textbf{（3） 驱动电路不只是“放大”，它还承担保护（退饱和检测、软关断、欠压封锁）}。}

% ============================================================
\section{电力电子器件的保护}

\subsection{过电压的产生及过电压保护}

电力电子装置中可能发生的过电压分为\textbf{外因过电压}与\textbf{内因过电压}两类。

\begin{table}[htbp]
\centering
\caption{过电压的分类与保护措施}
\label{tab:9-1}
\begin{tabular}{p{2.6cm}p{4.6cm}p{6.4cm}}
\toprule
\textbf{类别} & \textbf{产生原因} & \textbf{保护措施} \\
\midrule
外因过电压（来自外部） & 雷击过电压；供电系统操作过电压（断路器分合、熔断器熔断引起的电网电压波动） & 装置入口加\textbf{避雷器、压敏电阻、硒堆}；变压器屏蔽与静电屏蔽；进线加 RC 吸收 \\
内因过电压（装置内部） & \textbf{换相过电压}（器件关断时回路电感储能释放）；\textbf{关断过电压}（di/dt 大、线路电感产生尖峰）；器件反向恢复引起的尖峰 & 器件两端加 \textbf{RC 吸收电路}、\textbf{RCD 钳位}；合理安排布线以减小线路电感；选用软恢复二极管 \\
\bottomrule
\end{tabular}
\end{table}

\subsection{过电流保护}

电力电子电路运行不正常或发生故障时可能发生过电流，其原因有：\textbf{负载过载}、\textbf{直流侧短路}、\textbf{逆变失败（逆变颠覆）}、\textbf{器件损坏造成直通}等。过电流保护手段及其配合：

\begin{enumerate}
  \item \textbf{快速熔断器}：直接串在器件支路中，动作最快，用于短路保护；但熔断后须更换，且熔断过程中器件仍承受短路电流冲击。
  \item \textbf{过电流继电器 / 直流快速断路器}：用于负载过载与较大故障电流，动作较慢。
  \item \textbf{电子保护（电流检测 + 封锁脉冲）}：用霍尔电流传感器或分流器检测电流，超过阈值时立即\textbf{封锁触发脉冲}（或关断全控器件），动作最快、不损坏器件，是装置的主要保护手段；缺点是须与器件电流容量配合整定。
  \item \textbf{过载保护（热保护）}：按器件结温或散热器温度动作，用于长期过载。
\end{enumerate}

保护配合的原则是\textbf{“分级配合、由快到慢”}：电子保护最先动作（微秒级），快速熔断器作为后备（毫秒级），断路器最后（几十毫秒到秒级）。

\subsection{缓冲电路}

\textbf{缓冲电路}（Snubber，又称吸收电路）的作用是：抑制器件的 du/dt、di/dt 与电压电流尖峰，\textbf{改变器件的开关轨迹，使之落在安全工作区（SOA）内}，并可减少开关损耗（把损耗从器件转移到缓冲元件上）。

\begin{itemize}
  \item \textbf{关断缓冲电路（RC / RCD）}：并联在器件两端，器件关断时电容吸收回路电感释放的能量、限制 du/dt；开通时电容经电阻放电（RCD 中用二极管旁路电阻，使放电加快、减少开通损耗）。
  \item \textbf{开通缓冲电路}：串联电感限制开通时的 di/dt，避免电流集中与反向恢复尖峰。
  \item \textbf{复合缓冲电路}：同时具有开通与关断缓冲功能（如 $L$、$C$、$D$ 组合）。
  \item \textbf{无损缓冲电路}：用电感、电容、二极管构成，缓冲能量回馈到负载或电源而不消耗在电阻上，效率高，但结构复杂。
\end{itemize}

\corepoint{保护这一节的考点是“对号入座”：\textbf{外因过电压→避雷器/压敏电阻；内因过电压→RC 吸收、RCD 钳位；短路→快速熔断器；过载→电子保护/热保护；du/dt 过大→RC 缓冲；di/dt 过大→串联电感；二次击穿→GTR 的 SOA 与缓冲电路}。缓冲电路的核心作用是\textbf{把开关轨迹限制在 SOA 内}。}

% ============================================================
\section{电力电子器件的串联和并联使用}

当单个器件的电压或电流定额不能满足要求时，需将器件\textbf{串联}以提高耐压、\textbf{并联}以提高通流能力。但器件特性不可能完全一致，必须解决\textbf{均压}与\textbf{均流}问题。

\subsection{晶闸管的串联}

晶闸管串联时，由于各器件\textbf{反向漏电流与正反向阻断特性不同}，承受的电压不相等，须采取均压措施：

\begin{itemize}
  \item \textbf{静态均压}：在每只器件两端并联\textbf{阻值相等的均压电阻}，使静态电压按电阻分压而均匀；
  \item \textbf{动态均压}：在每只器件两端并联 \textbf{RC 吸收电路}，利用电容使开通关断过程中的电压变化均匀（动态均压电容同时起吸收作用）；
  \item \textbf{触发同时性}：串联器件必须\textbf{同时触发}，否则先导通的器件承受全部电压而损坏——须用强触发、同一触发信号源。
\end{itemize}

串联器件数按下式估算（考虑均压系数 $K_u$）：

\keyformula{n\ge\frac{U_m}{K_u\,U_{RRM}},\qquad K_u=0.8\sim0.9}

\subsection{晶闸管的并联}

晶闸管并联时，由于各器件\textbf{正向伏安特性不同}，同一电压下流过的电流不相等，须采取均流措施：

\begin{itemize}
  \item \textbf{串联电阻均流}：在每只器件支路串入小电阻，用电阻的压降强迫电流均匀（简单但损耗大）；
  \item \textbf{串联电抗器均流}：串入铁心电抗器（或空心电抗器），利用电感对电流变化的阻碍实现动态均流；
  \item \textbf{选配器件}：挑选\textbf{同一型号、同一批次、正向压降相近}的器件并联；
  \item \textbf{降额使用}：并联后总电流定额按均流系数 $K_i=0.8\sim0.9$ 折算。
\end{itemize}

\keyformula{n\ge\frac{I_m}{K_i\,I_{T(AV)}},\qquad K_i=0.8\sim0.9}

\subsection{电力 MOSFET 的并联和 IGBT 的并联}

电力 MOSFET 的通态电阻 $R_{on}$ 具有\textbf{正温度系数}：某只管子电流偏大时温度升高、$R_{on}$ 增大，电流自动减小——这种\textbf{自动均流}特性使 MOSFET 的并联比晶闸管容易得多。

IGBT 的通态压降 $U_{CE(sat)}$ 也有一定的正温度系数（在较大电流下），故也可并联，但仍需注意：

\begin{itemize}
  \item 各管\textbf{阈值电压 $U_{GE(th)}$ 与跨导 $g_{fs}$ 的差异}会造成开通时刻与电流分配不一致，须\textbf{各自串接栅极电阻}（不要共用一只电阻），必要时单独整定；
  \item 并联时须\textbf{对称布线}，使各支路的引线电感尽量相等，否则动态电流分配不均；
  \item 一般按 \textbf{90\%} 降额使用，并保证散热条件一致（同一散热器、对称安装）。
\end{itemize}

\corepoint{本节抓三条：\textbf{（1） 晶闸管串联靠“均压电阻 + RC 动态均压 + 同时触发”，并联靠“串电阻/电抗 + 选配 + 降额”}；\textbf{（2） 均压系数、均流系数一般取 0.8$\sim$0.9，串联/并联数按此折算}；\textbf{（3） MOSFET 因 $R_{on}$ 正温度系数可自动均流，IGBT 并联须各自串栅极电阻、对称布线}。}

\chapter{电力电子技术的应用}

前九章讲的是器件与电路，本章讲\textbf{装置}：把器件、电路、控制组合起来，解决实际的工程问题。教材按应用领域分了七节，主线是“电动机传动（最大市场）→ 电源（UPS、开关电源、PFC）→ 电力系统 → 其他”。

% ============================================================
\section{晶闸管直流电动机系统}

晶闸管可控整流装置带直流电动机负载组成的系统，习惯称为\textbf{晶闸管直流电动机系统}，是直流调速的主要形式。直流电动机的电枢回路有电阻 $R$、电感 $L$ 和反电动势 $E$，$E=C_e\Phi n$。

\subsection{工作于整流状态时}

电流连续时，整流输出平均电压 $U_d$ 与电动机反电动势平衡：

\keyformula{U_d=E+I_dR_\Sigma=C_e\Phi n+I_dR_\Sigma}

其中 $R_\Sigma$ 为电枢电阻与整流装置等效内阻（含换相压降折算）之和。把 $U_d=2.34U_2\cos\alpha-\Delta U_d$ 代入，得\textbf{机械特性}：

\keyformula{n=\frac{2.34U_2\cos\alpha-\Delta U_d-I_dR_\Sigma}{C_e\Phi}}

\begin{itemize}
  \item 改变 $\alpha$ 即改变理想空载转速 $n_0$，得到一族\textbf{平行下斜的直线}——这就是晶闸管直流调速的调压调速原理；
  \item 与直流发电机供电相比，晶闸管系统的机械特性斜率更大（等效内阻大），且\textbf{电流断续时特性明显变软}（$n$ 随 $I_d$ 减小而迅速升高，理想空载转速虚高），这是它特有的现象，分析时必须先判断电流是否连续。
\end{itemize}

\subsection{工作于有源逆变状态时}

$\alpha>90^\circ$（即逆变角 $\beta<90^\circ$）时 $U_d<0$，电动机作\textbf{发电机运行}（回馈制动），能量经变流器送回电网。此时

\keyformula{n=\frac{-2.34U_2\cos\beta-\Delta U_d-I_dR_\Sigma}{C_e\Phi}}

负号表示\textbf{逆变时电动机的转向与整流时相反}（相对于同一 $I_d$ 方向）。对应不同的 $\beta$ 可得到一族平行的机械特性，位于第二象限（正转制动）或第四象限（反转制动）。

\subsection{直流可逆电力拖动系统}

要实行四象限运行（正转电动、正转制动、反转电动、反转制动），须使电动机的\textbf{电压和电流都能反向}，常用两套变流装置\textbf{反并联}供电：

\begin{itemize}
  \item \textbf{有环流接线}：两组同时加触发脉冲，用\textbf{环流电抗器}限制两组之间的环流；切换平滑、无死区，但需笨重电抗器、损耗大（如三相半波有环流接线）。
  \item \textbf{无环流接线}：任何时刻只让一组工作，另一组封锁脉冲；切换时须先封锁、等电流过零并留死区，再开放另一组，故\textbf{切换有死区}，但效率高、无需环流电抗器（如三相全控桥无环流接线）——实际多用。
\end{itemize}

% ============================================================
\section{变频器和交流调速系统}

过去调速传动的主流是晶闸管直流电动机系统，但直流电动机有电刷与换向器、维护麻烦、容量受限。随着变频技术发展，\textbf{交流电动机的调速性能已可与直流电动机相媲美}，交流调速逐渐占据主导地位。

\subsection{交—直—交变频器}

\textbf{交—直—交变频器}（间接交流变流电路）由“整流器 + 中间直流环节 + 逆变器”构成，是通用变频器的主流结构。按中间直流环节的储能元件分为电压型与电流型。

\begin{itemize}
  \item \textbf{电压型间接交流变流电路}：中间为大电容，逆变器为电压型 PWM 逆变。若整流部分用不可控二极管整流，则\textbf{不能处理再生反馈的电力}（能量只能单方向流动），制动时需加制动电阻或制动单元把能量消耗掉，否则直流母线电压会升高（“泵升电压”）。
  \item \textbf{可逆的电压型电路}：整流和逆变均为 PWM 控制（双 PWM 变频器），可实现能量的双向流动，能再生制动、且输入电流正弦、功率因数可调。
  \item \textbf{电流型间接交流变流电路}：中间为大电感，可方便地实现再生反馈（通过改变整流桥的极性或采用可逆整流），但动态响应较慢。
\end{itemize}

\subsection{交—交变频器}

\textbf{交—交变频器}（直接变频、周波变流器）不经中间直流环节，用两组反并联相控整流直接把电网频率变成较低频率（见第 6 章）。其特点是\textbf{效率高、低频输出波形好、可四象限运行}，适合\textbf{低速大功率}传动（轧机、矿井提升机、水泥磨机）；缺点是输出频率不能超过输入频率的 $1/3\sim1/2$、功率因数低、输入谐波大。

\subsection{变频调速的基本控制方式}

由 $n_0=\dfrac{60f_1}{p}$ 可知，改变频率 $f_1$ 即可改变同步转速。但只变频不变压会引起磁通变化（$U_1\approx E_1=4.44f_1N_1\Phi_m$），因此须\textbf{变压变频（VVVF）}配合：

\keyformula{\frac{U_1}{f_1}=\text{const}\quad\Longrightarrow\quad \Phi_m\approx\text{const}}

\begin{itemize}
  \item \textbf{基频以下}：保持 $U_1/f_1=$ 常数（恒磁通），实现\textbf{恒转矩调速}；低频时定子电阻压降占比大，须做\textbf{低频电压补偿}（$U_1$ 略抬高）；
  \item \textbf{基频以上}：电压不能超过额定值，只能升频不升压，磁通随频率升高而减小，实现\textbf{恒功率调速}（弱磁调速）；
  \item 高性能调速还需\textbf{转差频率控制}、\textbf{矢量控制}（把定子电流分解为励磁分量与转矩分量分别控制）或\textbf{直接转矩控制}（DTC），以获得快速转矩响应。
\end{itemize}

% ============================================================
\section{不间断电源}

\textbf{不间断电源}（UPS）是当交流输入电源（习惯称为市电）发生中断或异常时，能持续向负载供电的装置，广泛用于计算机系统、通信机房、医院手术室、工业控制系统等不能断电的场合。

\begin{itemize}
  \item \textbf{组成}：整流器（把市电变成直流并向蓄电池充电）、蓄电池组（储能）、逆变器（把直流变成高质量交流）、静态开关（市电与逆变输出的切换）。
  \item \textbf{后备式 UPS}：市电正常时由市电直接供电（经简单滤波），断电时切换到逆变供电，切换时间约几毫秒；结构简单、成本低、效率高，但输出波形与电压质量一般。
  \item \textbf{在线式 UPS}：市电经“整流—逆变”始终由逆变器供电，输出波形纯净、无切换时间，是最可靠的方案，但效率较低、成本高。
  \item \textbf{在线互动式 UPS}：介于两者之间，市电正常时经变压器（可调压）供电，异常时逆变供电。
\end{itemize}

小容量 UPS 的主电路中，整流部分常用\textbf{二极管整流 + 直流斩波器}（兼作 PFC），以降低成本并提高输入功率因数。

% ============================================================
\section{开关电源}

各种电子设备需要多路不同电压的直流供电（如数字电路 5 V、3.3 V、1.8 V 等）。早期采用“工频变压器降压 + 整流滤波 + 线性调压”的\textbf{线性电源}：结构简单、纹波小，但\textbf{效率低（$30\%\sim50\%$）、体积大、重量重}。

采用\textbf{高频开关方式}进行电能变换的电源称为\textbf{开关电源}：市电整流滤波后经高频 DC-DC 变换器（几十到几百千赫兹）降压，再用高频变压器隔离、输出滤波。

\subsection{开关电源的结构}

典型结构：输入滤波与整流 → 功率因数校正（可选）→ 高频逆变（半桥/全桥/正激/反激）→ 高频变压器隔离 → 输出整流滤波 → 反馈控制。

与线性电源相比的优点：\textbf{效率高（可达 90\% 以上）、体积小、重量轻、输入电压范围宽}；缺点：\textbf{输出纹波与噪声较大、电磁干扰强、电路复杂}。

\subsection{开关电源的控制方式}

\begin{itemize}
  \item \textbf{PWM 方式}：开关频率固定，通过改变占空比稳压——最常用；
  \item \textbf{PFM 方式}：导通脉宽固定，通过改变开关频率稳压——轻载效率高，但滤波器设计困难；
  \item \textbf{混合方式}：重载用 PWM、轻载用 PFM（或跳周期），兼顾效率与纹波。
\end{itemize}

典型的开关电源控制系统为\textbf{电压闭环}（采样输出电压与基准比较，经 PI 调节器产生占空比指令）；要求高的场合再加\textbf{电流内环}（电流模式控制），以改善动态响应与均流。

\subsection{开关电源的应用}

广泛用于计算机、通信设备、家用电器、仪器仪表、LED 照明、电动汽车充电器等。它是“电力电子技术进入信息电子领域”的桥梁——所有电子装置的电源部分都是电力电子电路。

% ============================================================
\section{功率因数校正技术}

传统的二极管整流 + 电容滤波电路，输入电流是\textbf{窄脉冲}，功率因数只有 0.6 左右、谐波电流大（见 3.4、3.5 节）。为满足电网谐波标准（如 IEC 61000-3-2），须采用\textbf{功率因数校正（PFC）}技术。

\subsection{功率因数校正电路的基本原理}

目标：使输入电流\textbf{跟随输入电压的波形}（同频同相的正弦波），从而使 $\nu=I_1/I\to1$、$\cos\varphi_1\to1$，即 $\lambda\to1$。

\begin{itemize}
  \item \textbf{无源 PFC}：在整流桥后加电感（或“填谷”电路），把窄脉冲电流展宽，结构简单、成本低，可将 PF 提高到 0.8 左右；
  \item \textbf{有源 PFC（APFC）}：在整流桥与负载之间插入一个\textbf{Boost 变换器}，用控制环使输入电流跟踪输入电压的正弦波形：
  \begin{enumerate}
    \item 电压外环：采样输出直流电压，与基准比较后经误差放大器得到\textbf{电流指令的幅值}；
    \item 乘法器：把该幅值与整流后的输入电压波形相乘，得到\textbf{正弦电流指令}；
    \item 电流内环：控制开关使输入电流跟踪该指令，于是输入电流为正弦且与电压同相。
  \end{enumerate}
  \item 效果：PF 可达 0.99 以上、输入电流 THD 小于 5\%，且输出电压稳定（Boost 的升压特性使输出约为 400 V 直流）。
\end{itemize}

\keyformula{\lambda=\frac{I_1}{I}\cos\varphi_1\to1\qquad(\text{APFC 的目标})}

\subsection{单级功率因数校正技术}

上述基于 Boost 的有源 PFC 属“两级”结构（PFC 级 + DC-DC 级），成本较高。\textbf{单级 PFC} 把 PFC 功能与 DC-DC 变换合并到一级完成：在反激等变换器的输入端加少量元件（如填谷电容、耦合电感），使输入电流在半个工频周期内近似跟随输入电压，用一个开关同时完成 PFC 与稳压。

\begin{itemize}
  \item 优点：器件少、成本低，适合小功率（如 LED 驱动、适配器）；
  \item 缺点：储能电容电压随负载与输入电压变化较大、PF 与 THD 指标不如两级方案，输出纹波较大。
\end{itemize}

% ============================================================
\section{电力电子技术在电力系统中的应用}

\subsection{高压直流输电}

\textbf{高压直流输电（HVDC）}在长距离、大容量、跨区域输电以及海底电缆输电中有明显优势：没有交流线路的电容充电电流与稳定极限问题，线路损耗小、可控性强。

\begin{itemize}
  \item 送端\textbf{整流站}把交流变成直流，受端\textbf{逆变站}把直流变回交流，两端均为晶闸管换流阀（相控换流）。
  \item 传统的晶闸管换流阀需要电网提供换相电压（电网换相），故不能向无源网络供电；\textbf{轻型直流输电（柔性直流）}采用全控型器件 IGBT 构成的电压型换流器，可向弱电网甚至无源网络供电，是有源配电网、海上风电并网的重要技术。
  \item 换流站会产生大量谐波与无功消耗，须配\textbf{交流滤波器}与\textbf{无功补偿装置}。
\end{itemize}

\subsection{无功功率控制}

无功功率直接影响电网电压水平与线路损耗，对无功功率的控制是电力系统运行的核心问题之一。电力电子型无功补偿装置：

\begin{itemize}
  \item \textbf{晶闸管控制电抗器（TCR）}：用晶闸管的相控调节电抗器电流，连续吸收感性无功；
  \item \textbf{晶闸管投切电容器（TSC）}：分组投切电容器，分级提供容性无功（在电压过零投入以避免冲击）；
  \item \textbf{静止无功补偿器（SVC）}：TCR 与 TSC（或固定电容器 FC）组合，可连续调节；
  \item \textbf{静止无功发生器（SVG / STATCOM）}：用全控型器件构成的电压型逆变器，输出与电网同频的电压，通过调节其幅值控制无功电流的方向与大小，响应快、谐波小、低压时仍能输出额定无功，性能优于 SVC。
\end{itemize}

\subsection{谐波抑制与电能质量控制}

\begin{itemize}
  \item \textbf{无源滤波器}：LC 谐振支路，对特定次谐波呈低阻抗，结构简单但只能滤固定次谐波、易与电网谐振；
  \item \textbf{有源电力滤波器（APF）}：检测负载谐波电流，用 PWM 逆变器产生\textbf{大小相等、方向相反}的补偿电流将其抵消，可动态补偿各次谐波与无功，是治理谐波最有效的手段；
  \item \textbf{统一电能质量调节器、动态电压恢复器（DVR）}：用于抑制电压跌落、闪变、瞬时中断，改善供电质量；
  \item \textbf{柔性交流输电（FACTS）}：用电力电子装置控制线路的阻抗、电压与相位，提高输电能力与稳定性。
\end{itemize}

% ============================================================
\section{电力电子技术的其他应用}

\begin{itemize}
  \item \textbf{焊接电源}：逆变式焊机把工频整流后经高频逆变（20 kHz 以上）再整流输出，变压器体积重量大幅减小、动态响应快、焊接工艺性好，已取代传统工频焊机。
  \item \textbf{照明电源}：电子镇流器（高频点亮荧光灯，提高光效、消除频闪）、LED 驱动电源（恒流输出，通常含 PFC 级）；电力电子照明电源被称为“节能灯”，正逐步取代白炽灯与荧光灯。
  \item \textbf{电解与电镀}：需要低电压大电流的直流电源，多用双反星形或十二脉波整流（见 3.6 节）。
  \item \textbf{感应加热}：中频/高频电源（并联谐振式逆变电路，见 4.3 节）用于熔炼、淬火、焊接。
  \item \textbf{新能源与储能}：光伏并网逆变器、风力发电变流器、电池储能 PCS、电动汽车充电机——都是电力电子装置，且都要求高效率、高功率密度、单位功率因数与能量双向流动。
\end{itemize}

\corepoint{本章复习抓“装置—电路—指标”的对应关系：\textbf{直流调速}对应相控整流 + 电动机机械特性（连续/断续、整流/逆变）；\textbf{变频器}对应交直交（电压型/电流型、能否再生）与 $U_1/f_1=$ 常数；\textbf{UPS} 对应后备式/在线式与静态开关；\textbf{开关电源}对应高频 DC-DC + PWM 控制；\textbf{PFC} 对应“输入电流跟随输入电压”的 Boost APFC 与乘法器控制；\textbf{电力系统}对应 HVDC、SVC/SVG、APF。答题时先说\textbf{用什么电路}，再说\textbf{靠什么控制量调节}，最后说\textbf{优缺点与适用场合}。}

\end{document}